\documentclass[11pt,a4paper]{article}
\usepackage[T1]{fontenc}
\usepackage[utf8]{inputenc}
\usepackage{lmodern,amsmath,amssymb,amsthm,mathtools,microtype,mathrsfs}
\usepackage[margin=25mm,top=23mm,bottom=25mm]{geometry}
\usepackage{booktabs,array,longtable,enumitem,xcolor,fancyhdr}
\usepackage[colorlinks=true,linkcolor=blue!45!black,urlcolor=blue!45!black,citecolor=blue!45!black]{hyperref}
\hypersetup{pdftitle={Reading the Alpoge--Claude RH and Dirichlet-GRH proof: working version 0.44},pdfauthor={ChatGPT},pdfsubject={Partial human-readable reconstruction with explicit arithmetic inputs},pdfkeywords={Riemann hypothesis, Lean, theta lifts, prime sums, Arakelov slopes}}
\setlength{\parindent}{0pt}
\setlength{\parskip}{5pt plus 1pt}
\setlist{topsep=3pt,itemsep=3pt,leftmargin=*}
\setlength{\headheight}{14pt}
\pagestyle{fancy}\fancyhf{}
\fancyhead[L]{\small RH and Dirichlet GRH: a reconstruction}
\fancyhead[R]{\small Working version 0.44}
\fancyfoot[C]{\thepage}
\renewcommand{\headrulewidth}{0.3pt}
\numberwithin{equation}{section}
\newtheorem{theorem}{Theorem}[section]
\newtheorem{proposition}[theorem]{Proposition}
\newtheorem{lemma}[theorem]{Lemma}
\newtheorem{corollary}[theorem]{Corollary}
\theoremstyle{remark}\newtheorem{remark}[theorem]{Remark}
\newcommand{\Rea}{\operatorname{Re}}
\newcommand{\HS}{\mathrm{HS}}
\newcommand{\tr}{\operatorname{tr}}
\newcommand{\codim}{\operatorname{codim}}
\newcommand{\covol}{\operatorname{covol}}
\newcommand{\OK}{\mathcal O_K}
\newcommand{\src}[2]{\href{https://github.com/CocoSgt/riemann-hypothesis-lean/blob/7cbf7a0c32fca1badf05931879a37c3f8a2b03b4/src/#1}{\texttt{#2}}}
\newcommand{\status}[1]{\textbf{Status:} #1}
\begin{document}
\begin{center}
{\LARGE\bfseries RH and Dirichlet GRH: new reconstruction chapters}\\[7pt]
{\large Version 0.44 --- 23 November 2026}\\[5pt]
{\large Automorphic parts, ordinary points, and the same ideal}
\end{center}
\textbf{Scope.} This edition contains the current audit and the seven new
chapters. It remains a partial human reconstruction. Earlier named
claims are in the cumulative manuscript and preserved component files.
The underlying formal proof is accepted as correct; the exact remaining
human reconstruction boundaries are stated below.
\setcounter{tocdepth}{1}
\tableofcontents
\clearpage
\begingroup\emergencystretch=3em
\section{Current dependency audit: automorphic parts, ordinary points, and the same ideal}
\label{sec:dependency-audit-v44}

\subsection{Current status and reading order}

Version 0.44 reconstructs the original Galois-decorated automorphic
parts and installs them in the ordinary Hecke construction.  It also
proves the complete ordinary Hecke-point record and the original
ring-properties and density statements for the same ideal $J$.
The named base-change, polarized cuspidal attachment, local
representation and period foundations retain their exact original
scope and are identified in the chapters.

The pinned Alpoge--Claude formal proof is accepted as fully checked
and correct.  This remains a \textbf{partial human reconstruction}
of RH and unrestricted Dirichlet GRH.  The four subsequent Hida
endpoints, the all-level same-ideal theorem on both rank-four Links,
the common geometric and arithmetic families, and the remaining
global quantitative producers have not all been reconstructed.
No Lean build or kernel replay is claimed.

There are 119 new and 2560 cumulative claims.  All 2441
predecessor claim records, source arrays and mathematical dependency
edges are preserved.  The selected sufficient-route graph keeps
121 nodes and 161 edges.  Its 31 whole open obligations remain
22 class-C and 9 class-B tasks: 24 are shared by RH and Dirichlet
GRH, and 7 are additional for Dirichlet GRH.  This increment
completes several substantial nested producers inside the compound
support obligation; it does not close that entire obligation.
These inventories are not percentages of proof completion.

Read LTP, BCS and GUA for local transport, the base-change-to-Speh
bridge, and the single assembled Galois representation.  UPA
constructs the complete original automorphic parts and installs
same-ideal density.  OFL proves the ordered ordinary-flag argument;
HEP and JRP then give the complete Hecke-point output and ring
properties.  The separate reader's edition contains this audit and
these seven chapters.  The cumulative edition also preserves every
earlier proof chapter.  CMP1--CMP6 remain the full conditional
terminal deduction, with their original hypotheses.

\subsection{One actual automorphic space and its local representations}

LTP1--LTP21 work on the original coefficient-valued functions and
their exact scalarizations.  A finite separating family transports
isotypy, and the original compact level equation retains its entire
coefficient matrix.  This gives all seven PartUR2 suppliers,
including the literal dual infinity weight, every original
spherical index, and the depth-one $K_c(1)$ condition.
The first six suppliers add no open-kernel hypothesis on the
integral type.  Outside-place unramifiedness is proved even at
nonsplit places.

Actual split local components are constructed from the positive
invariant form and finite multiplicity on the original automorphic
space.  The argument produces an admissible irreducible submodule
and a nonzero evaluation intertwiner.  It also proves the
unramified--spherical equivalence for every local component.
Choice is made once, over all split places, before that family is
used in the different fields of the arithmetic part.

\subsection{The complete base-change and Galois consumers}

BCS1--BCS19 consume the original
\texttt{Base\allowbreak Change\allowbreak Isobaric\allowbreak N} foundation.  Unitarity of the actual
discrete local factors and the retained unitary-product theorem
replace the selected local subquotient by the entire normalized
product.  The original base-change witnesses already contain the
ascending embedding needed for each residual constituent; consuming
it directly avoids an extra global residual-spectrum classification.

Smooth Hermitian duality changes that ascending local injection
into the exact descending quotient.  The explicitly matched
generic-unitary classification, together with an elementary
centered-segment inequality, proves the unique cosocle for every
positive Speh length.  The Capelli computation keeps the infinity
shifts, while compact-fixed exactness and the finite Satake
calculation keep the split, inert and all-unramified clauses.
The result is the entire original \texttt{Bc\allowbreak Speh} record.

GUA1--GUA18 then use the original polarized cuspidal attachment
\texttt{Row\allowbreak T}/\texttt{Props\allowbreak B}.  If a cusp has rank $a_i$ and
Speh length $d_i$, its $k$th parameter occurs with twist
\[
                 \frac{d_i-2k-n}{2},\qquad 0\le k<d_i.
\]
The half shifts are realized by actual automorphic determinant
twists before applying the attachment theorem.  No square root of
the cyclotomic character is presumed to exist as a Galois character.
Finite direct sums are verified on the actual period modules,
filtrations and coefficient embeddings.  Goodness of the total
parameter implies goodness of each constituent for the same
correspondence, preserving the source's conditional local
compatibility and both universal quantifiers.

The spherical reading applies to every original \texttt{Is\allowbreak Rec\allowbreak T}.
It gives trivial original inertia and the geometric Frobenius
polynomial at every required unramified place.  Conjugate
self-duality of the Satake multisets, followed by Chebotarev and
semisimple character comparison, gives the pairing with multiplier
$\epsilon^{1-n}$ for every original complex conjugation.
If the assembled representation is irreducible, its positive
direct-sum blocks force exactly one cusp and one Speh copy.  The
actual cusp then has all the additional local and infinity clauses
of the original endpoint.  Speh length one is a conclusion of this
conditional branch, not a restriction on the construction.

\subsection{The actual part record and bare same-ideal density}

UPA1--UPA18 descend the same representation through the actual
$V.\mathrm{pr}$, using the closed normal subgroup generated by
the original outside inertia groups.  Surjectivity of this map
preserves the invariant-subspace lattice, hence semisimplicity and
irreducibility.  The pairing descends with the same invertible
matrix.  The spherical computation keeps
\[
 \operatorname{heckePol}(n,q,a)
   =\sum_{j=0}^n(-1)^jq^{j(j-1)/2}a_jX^{n-j}
   =\prod_{\beta\in B}(X-q^{(n-1)/2}\beta),
\]
and compares it to $R(g^{-1})$ for an arithmetic Frobenius lift
$g$.  The inverse is retained explicitly.

At the type places, zero monodromy kills an open inertia subgroup.
Every original inertia matrix is then of finite order, and the
commuting semisimple/nilpotent comparison identifies its action
with the specified type.  Pulling one complex intertwiner through
the fixed algebraic isomorphism $\iota$ gives one coefficient-field
matrix for all inertia elements.  This works for every original
Teichmuller family, including repeated characters.

The full local Hodge comparison gives the original
\texttt{At\allowbreak LT}.  These constructions fill all sixteen
\texttt{Part\allowbreak Alt\allowbreak Alt} fields and all four
\texttt{Unitary\allowbreak Point\allowbreak Properties} with one representation, one
automorphic space and one family of local components.
They prove the strengthened \texttt{part\allowbreak X\_\allowbreak alt} and its
\texttt{part\allowbreak U\_\allowbreak alt} projection.  The actual decomposition ADP
and the integral type ITA therefore install
\texttt{Au\allowbreak Teich\allowbreak Ex}'s \texttt{Automorphic\allowbreak Hecke\allowbreak Set}.

The AFS separation proof now applies with that input supplied.
Consequently the bare original \texttt{TJDense} follows from
\texttt{Is\allowbreak J} on the same original $B,\Lambda,K,J$, retaining
$V.S'=s.S$.  In the nonzero branch $J=\ker(B\to T)$;
in the zero-module branch $J=B$.  No replacement radical or
support ideal is introduced.

\subsection{Ordered flags, Hecke points, and all ring properties}

OFL1--OFL16 supply the complete ordinary-flag consumer.
The original rank-one Hecke-character inputs give the exact local
\texttt{DRat} after globalizing a weight on an auxiliary CM field
with the same local completion.  Distinct inertia characters have
one-dimensional finite-group projectors, so any nilpotent operator
commuting with inertia is zero.  The proof keeps the raw Weil
action separate from its Frobenius semisimplification: their
characteristic polynomials agree, and the scalarization map
intertwines the raw action.

Finite averaging and all Frobenius translates descend an invariant
to the original period module.  A Vandermonde sum of a periodic
period produces a nonzero vector fixed by Frobenius and Galois.
When the labeled weights are nonnegative, filtration zero and
the fundamental crystalline intersection recover a nonzero
Galois-invariant vector.  Twisting by the inverse of the first
ordinary character makes precisely those weights nonnegative.
Removing the resulting eigenline and repeating gives the entire
flag in its prescribed order.  Potential crystallinity holds for
every original bridge; no semisimplicity premise is added.

HEP proves the actual Hecke-eigenpoint dictionary.  Prefix products
of normalized $U$-eigenvalues telescope to individual characters,
and the original Artin and cyclotomic formulas identify those
characters on the full local Weil group.  Dense Frobenius values
identify the determinant law of the same part representation.
Together with AFS and OFL this completes the original
\texttt{Out.\allowbreak pt}, including the exact good-point guard on its
flag and potential-crystalline conclusion, and hence the entire
\texttt{Ord\allowbreak Core\allowbreak H.\allowbreak Out} record.

JRP descends all eight original \texttt{Ring\allowbreak Props} fields through
the jointly separating Hecke points.  Determinant identities are
checked as polynomial laws after arbitrary scalar extension, by
injectivity of polynomial coefficients; equality on group elements
alone is not substituted for equality of laws.  The argument keeps
the form identity, both kernel relations, the ordinary
characteristic-polynomial and ordered-product conditions, and every
Teichmuller inertia type.  The original auxiliary-place
$\mathrm{V0ok}$ condition supplies the exact inertia-vanishing
consumer.  Finiteness and torsionfreeness pass through the actual
quotient isomorphism $B/J\simeq T$.  The top-ideal branch is
handled separately.

Thus the \texttt{ring} and \texttt{dense} fields of the later
six-field \texttt{Hida\allowbreak Props\allowbreak Alt} record are now constructed.
Its \texttt{pt}, \texttt{end1}, \texttt{end2}, and
\texttt{twist} fields remain distinct reconstruction targets.
The ordinary Hecke-point record proved here is not being confused
with that later \texttt{Point\allowbreak Theorem\allowbreak Base\allowbreak Object} field.

\subsection{Every remaining selected obligation}

Each contract below is the full selected output.  Completed nested
theorems are available inputs; a compound obligation is closed only
when all its remaining suppliers, shared objects and choice-order
conditions are supplied.  The current notes govern the historical
frontier descriptions preserved in the cumulative manuscript.
\paragraph{1. GLOBAL\_\allowbreak DENOMINATOR — C — Actual LinkN local table and binding producer}

Closed(Des $\mathcal C$ 3) LinkN: the actual etaSum is bounded with exact dZ0 and same-host vvTerm, subject only to the literal ShaFinOdd and summability guards.

\emph{Shared objects and choice order.} one FamP; one class $\mathcal A$ and its target-compatible subclass; one coefficient field per family datum; one x and i at each application; original selected rows, words and levels; thresholds chosen before the common family; actual D.nd; actual etaJ and Aj; sha\_\allowbreak eq

\emph{Source.} \href{https://github.com/CocoSgt/riemann-hypothesis-lean/blob/7cbf7a0c32fca1badf05931879a37c3f8a2b03b4/src/Formalism/SharpRestriction/GlobalTheoremInputs.lean#L45-L48}{LinkN}.

\emph{Source.} \href{https://github.com/CocoSgt/riemann-hypothesis-lean/blob/7cbf7a0c32fca1badf05931879a37c3f8a2b03b4/src/Formalism/SharpRestriction/GlobalTheoremInputs.lean#L40-L44}{LinkN\_\allowbreak at}.

\emph{Source.} \href{https://github.com/CocoSgt/riemann-hypothesis-lean/blob/7cbf7a0c32fca1badf05931879a37c3f8a2b03b4/src/Development/AutomorphicArithmetic/CharacterBoxes/DivinNOfHypV.lean#L117-L134}{divinN\_\allowbreak of\_\allowbreak hypV}.



\paragraph{2. GLOBAL\_\allowbreak RETAINED\_\allowbreak DEFECT — C — Original retained-row defect bound LinkIRR}

Rows.DelpNonneg and delp$\le$exp(-x\textasciicircum{}(ez/7)) uniformly for every late actual row datum.

\emph{Shared objects and choice order.} one FamP; one class $\mathcal A$ and its target-compatible subclass; one coefficient field per family datum; one x and i at each application; original selected rows, words and levels; thresholds chosen before the common family; AllRows A; LabelCountGoodInjective; AllJoinA; CountAllRows

\emph{Retained predecessor progress.} FLC1–FLC19, reusing RRT1–RRT16 and GQT1–GQT5. Actual supplied normalized-type compatibility, original HidaSetting, reduction/inflation EtOne with the same Haar normalization, RowsAdd/RowsGood, one common AllRows object and the three arithmetic cylinder counts. All five selected Wg terms now have local producers, with the explicit eventual rank threshold for the last. The stronger unused fourth admLq upper comparison is not asserted.

\emph{Source.} \href{https://github.com/CocoSgt/riemann-hypothesis-lean/blob/7cbf7a0c32fca1badf05931879a37c3f8a2b03b4/src/Formalism/SharpRestriction/GlobalTheoremInputs.lean#L55-L58}{LinkIRR}.

\emph{Source.} \href{https://github.com/CocoSgt/riemann-hypothesis-lean/blob/7cbf7a0c32fca1badf05931879a37c3f8a2b03b4/src/Development/AbelianVarietyTheory/Proofs/GlobalInterpolation/LinkIRROfJoinEx.lean#L17-L38}{linkIRR\_\allowbreak of\_\allowbreak join\_\allowbreak ex}.

\emph{Source.} \href{https://github.com/CocoSgt/riemann-hypothesis-lean/blob/7cbf7a0c32fca1badf05931879a37c3f8a2b03b4/src/Development/AbelianVarietyTheory/Proofs/GlobalInterpolation/LinkIRROfJoinEx.lean#L143-L166}{link\_\allowbreak irr}.



\paragraph{3. GLOBAL\_\allowbreak PERIOD\_\allowbreak HEIGHT — C — Full actual period and height comparison LinkDIVt}

Closed conditional TuSize -> DGamSize -> (LinkEta -> DIVinN -> LinkDIV), with every original LinkDIV\_\allowbreak at clause.

\emph{Shared objects and choice order.} one FamP; one class $\mathcal A$ and its target-compatible subclass; one coefficient field per family datum; one x and i at each application; original selected rows, words and levels; thresholds chosen before the common family; same RR target; same Bu/besS; same Nu and Bk

\emph{Source.} \href{https://github.com/CocoSgt/riemann-hypothesis-lean/blob/7cbf7a0c32fca1badf05931879a37c3f8a2b03b4/src/Formalism/SharpRestriction/GlobalTheoremInputs.lean#L77-L90}{LinkDIV\_\allowbreak at}.

\emph{Source.} \href{https://github.com/CocoSgt/riemann-hypothesis-lean/blob/7cbf7a0c32fca1badf05931879a37c3f8a2b03b4/src/Formalism/SharpRestriction/GlobalTheoremInputs.lean#L94-L96}{LinkDIVt}.

\emph{Source.} \href{https://github.com/CocoSgt/riemann-hypothesis-lean/blob/7cbf7a0c32fca1badf05931879a37c3f8a2b03b4/src/Development/AbelianVarietyTheory/Proofs/GlobalInterpolation/LinkBrest.lean#L7-L25}{GlobalTheory.Top.link\_\allowbreak brest}.

\emph{Source.} \href{https://github.com/CocoSgt/riemann-hypothesis-lean/blob/7cbf7a0c32fca1badf05931879a37c3f8a2b03b4/src/Development/AbelianVarietyTheory/Proofs/GlobalInterpolation/WireSel.lean#L11-L41}{GlobalTheory.Top.wire\_\allowbreak sel}.



\paragraph{4. GLOBAL\_\allowbreak LOCAL\_\allowbreak GAIN — C — Original global LinkG installation}

One uniform LinkG with actual local lattice lengths, served-orbit coverage and MeanKappaBound; delta(x)->0 before varying actual indices.

\emph{Shared objects and choice order.} one FamP; one class $\mathcal A$ and its target-compatible subclass; one coefficient field per family datum; one x and i at each application; original selected rows, words and levels; thresholds chosen before the common family; actual G.Inst; actual GK; actual lenMN; same orbit representatives

\emph{Source.} \href{https://github.com/CocoSgt/riemann-hypothesis-lean/blob/7cbf7a0c32fca1badf05931879a37c3f8a2b03b4/src/Formalism/SharpRestriction/GlobalTheoremInputs.lean#L122-L137}{LinkG\_\allowbreak at}.

\emph{Source.} \href{https://github.com/CocoSgt/riemann-hypothesis-lean/blob/7cbf7a0c32fca1badf05931879a37c3f8a2b03b4/src/Formalism/SharpRestriction/GlobalTheoremInputs.lean#L138-L141}{LinkG}.

\emph{Source.} \href{https://github.com/CocoSgt/riemann-hypothesis-lean/blob/7cbf7a0c32fca1badf05931879a37c3f8a2b03b4/src/Development/AbelianVarietyTheory/Proofs/GlobalInterpolation/LinkG.lean#L14-L27}{GlobalTheory.Top.link\_\allowbreak g}.



\paragraph{5. GLOBAL\_\allowbreak SPECTRAL — C — Original normalized kernel/Gram estimates LinkS}

Closed TuSize -> LinkS with exact Gamma normalization, off-diagonal Hilbert-Schmidt bound, diagonal error and Dall/Dstar estimate, errors exp(-x\textasciicircum{}(epsilon/8)).

\emph{Shared objects and choice order.} one FamP; one class $\mathcal A$ and its target-compatible subclass; one coefficient field per family datum; one x and i at each application; original selected rows, words and levels; thresholds chosen before the common family; actual KD; same S.Cu spaces; GamOf; Dstar; Dall

\emph{Source.} \href{https://github.com/CocoSgt/riemann-hypothesis-lean/blob/7cbf7a0c32fca1badf05931879a37c3f8a2b03b4/src/Formalism/SharpRestriction/GlobalTheoremInputs.lean#L145-L149}{LinkS\_\allowbreak at}.

\emph{Source.} \href{https://github.com/CocoSgt/riemann-hypothesis-lean/blob/7cbf7a0c32fca1badf05931879a37c3f8a2b03b4/src/Formalism/SharpRestriction/GlobalTheoremInputs.lean#L152-L157}{LinkS}.

\emph{Source.} \href{https://github.com/CocoSgt/riemann-hypothesis-lean/blob/7cbf7a0c32fca1badf05931879a37c3f8a2b03b4/src/Development/AbelianVarietyTheory/Proofs/GlobalInterpolation/LinkAlt.lean#L16-L50}{GlobalTheory.Top.link\_\allowbreak alt}.



\paragraph{6. TASK\_\allowbreak PROP\_\allowbreak GK\_\allowbreak INSTALL — B — Install OneChar and the congruence integer into PropGK}

Complete the finite orbit/Frobenius/valuation argument producing one nonzero integer before every later coefficient prime, with the original e,N,kernel,p and fixed-K constants.

\emph{Source.} \href{https://github.com/CocoSgt/riemann-hypothesis-lean/blob/7cbf7a0c32fca1badf05931879a37c3f8a2b03b4/src/Development/AutomorphicArithmetic/AutomorphicVarieties/UnitaryThreeGroups/KatoUse.lean#L36-L42}{PropGK}.

\emph{Source.} \href{https://github.com/CocoSgt/riemann-hypothesis-lean/blob/7cbf7a0c32fca1badf05931879a37c3f8a2b03b4/src/Development/AutomorphicArithmetic/AutomorphicVarieties/UnitaryThreeGroups/OrbEquiv.lean#L111-L173}{UpperTheory.V.U3.propGK}.



\paragraph{7. TASK\_\allowbreak ONE\_\allowbreak CHAR\_\allowbreak PROOF — C — Reconstruct the original one-character argument}

From exact InpZ and AutomorphyUnramifiedCoefficients, prove the original OneChar trace/length estimate on the supplied coefficient datum and GD input.

\emph{Source.} \href{https://github.com/CocoSgt/riemann-hypothesis-lean/blob/7cbf7a0c32fca1badf05931879a37c3f8a2b03b4/src/Development/AutomorphicArithmetic/AutomorphicVarieties/UnitaryThreeGroups/KatoUse.lean#L54-L64}{OneChar}.

\emph{Source.} \href{https://github.com/CocoSgt/riemann-hypothesis-lean/blob/7cbf7a0c32fca1badf05931879a37c3f8a2b03b4/src/Development/AutomorphicArithmetic/AutomorphicVarieties/ModularityArgument/KatoQuadraticCharacters/OneChar.lean#L158-L178}{UpperTheory.V.Z.Kato.oneChar}.



\paragraph{8. TASK\_\allowbreak INP\_\allowbreak Z\_\allowbreak ASSEMBLY — B — Install the exact Kato input record}

Build the original CohInp/InpZ record using exactly lemM,lemMprime,lemP,RigII,AdVanishK and number on the same arithmetic objects. No rank-four replacement is permitted.

\emph{Source.} \href{https://github.com/CocoSgt/riemann-hypothesis-lean/blob/7cbf7a0c32fca1badf05931879a37c3f8a2b03b4/src/Development/AutomorphicArithmetic/AutomorphicVarieties/ModularityArgument/KatoQuadraticCharacters/KatoTwoOfUseN.lean#L55-L61}{InpZ}.



\paragraph{9. RIG\_\allowbreak II — C — Type-II polarized adjoint rigidity}

Marked original datum, actual type-II point and modularity of all source extensions imply vanishing for its exact adRep.

\emph{Source.} \href{https://github.com/CocoSgt/riemann-hypothesis-lean/blob/7cbf7a0c32fca1badf05931879a37c3f8a2b03b4/src/Development/AutomorphicArithmetic/AutomorphicVarieties/ModularityArgument/KatoQuadraticCharacters/KatoTwoOfUseN.lean#L49-L53}{RigII}.

\emph{Source.} \href{https://github.com/CocoSgt/riemann-hypothesis-lean/blob/7cbf7a0c32fca1badf05931879a37c3f8a2b03b4/src/Development/AutomorphicArithmetic/AutomorphicVarieties/ModularityArgument/KatoQuadraticCharacters/KatoTwoOfUseN.lean#L43-L47}{KatoTwo}.

\emph{Source.} \href{https://github.com/CocoSgt/riemann-hypothesis-lean/blob/7cbf7a0c32fca1badf05931879a37c3f8a2b03b4/src/Development/AutomorphicArithmetic/AutomorphicVarieties/ModularityArgument/KatoQuadraticCharacters/KatoTwoOfUseN.lean#L11-L41}{Used}.



\paragraph{10. AUT\_\allowbreak UNRAM\_\allowbreak COEF — C — Actual unramified-coefficient automorphy}

For each actual Dat with coefficient ramification one and fixed iota, choose one marked g1 supplying FinA, AutB and LemW together.

\emph{Source.} \href{https://github.com/CocoSgt/riemann-hypothesis-lean/blob/7cbf7a0c32fca1badf05931879a37c3f8a2b03b4/src/Development/AutomorphicArithmetic/AutomorphicVarieties/AutomorphicCases/AutomorphyUnramifiedCoefficients.lean#L24-L27}{AutomorphyUnramifiedCoefficients}.

\emph{Source.} \href{https://github.com/CocoSgt/riemann-hypothesis-lean/blob/7cbf7a0c32fca1badf05931879a37c3f8a2b03b4/src/Development/AutomorphicArithmetic/AutomorphicVarieties/AutomorphicCases/AutomorphyUnramifiedCoefficients.lean#L12-L13}{FinA}.

\emph{Source.} \href{https://github.com/CocoSgt/riemann-hypothesis-lean/blob/7cbf7a0c32fca1badf05931879a37c3f8a2b03b4/src/Development/AutomorphicArithmetic/AutomorphicVarieties/AutomorphicCases/AutomorphyUnramifiedCoefficients.lean#L15-L17}{AutB}.

\emph{Source.} \href{https://github.com/CocoSgt/riemann-hypothesis-lean/blob/7cbf7a0c32fca1badf05931879a37c3f8a2b03b4/src/Development/AutomorphicArithmetic/AutomorphicVarieties/AutomorphicCases/AutomorphyUnramifiedCoefficients.lean#L19-L22}{LemW}.



\paragraph{11. CNR3 — C — Uniform rank-three and rank-two congruence integer}

Fixed quadratic K -> c,A -> e,N,p,w,finite Galois-stable roots-of-unity orbit -> nonzero integer NN -> every later coefficient prime, closure, embedding, Frobenius and finite admissible collection.

\emph{Source.} \href{https://github.com/CocoSgt/riemann-hypothesis-lean/blob/7cbf7a0c32fca1badf05931879a37c3f8a2b03b4/src/Development/AutomorphicArithmetic/AutomorphicVarieties/AutomorphicForms/CongruenceNumberRankThree.lean#L125-L140}{CongruenceNumberRankThree}.



\paragraph{12. COUNT\_\allowbreak II\_\allowbreak SUPPORT — C — Two actual rank-four Links with common threshold and same-ideal Hida/P1}

For each of the TWO original QX rank-four Links, at fixed iota choose ONE common coefficient threshold before every later coefficient DVR and actual universal B,Lambda; produce one ideal J carrying ring,dense,pt,end1,end2,twist AND P1core. Inherited JSP/CIP/CAP/CCH/CFO supply the same-J consumer, every-map reciprocity/universal NiceI Norm, ArPl under unchanged CohP, supplied-type IsRes, and the full selected h=0,a=n=4,g=0 Choice with coupled B1/B2, D=4[L+:Q]+4T.q and B2 constant zero. The dualN bound retains one C before every m and an eventual N threshold uniform over each original admissible rep-prime under its finiteness premise, on actual TWPl. LIM constructs the full original S4. PRG now constructs the full original RRObj/RMaps, every finite-level map and actual ordered splitting target; PMS constructs the actual all-level PMod and full S5 with exact support, dimension, fiber and move fields, plus piKer/piSurj and the same-J annihilator equivalence. SCI constructs the actual quotient determinant, OoPair/ooMaps and point/OoPt on one splitting prime. The selected CRW full-block cotangent bound forces the support ideal to vanish, proving selected normalized Core. BCN supplies full Core=>CoreBC and its selected normalization/Unit consumer; PQL supplies full original TWQuot and matched pro-ell descent. Together SCI13 proves selected original Raw at h=0,a=n=4,k.g=0, with final B,J,K,f and every Selmer quantifier unchanged. NTL17 now proves full original neatU for arbitrary integral tau0 and every level. OTW12, ITA1–ITA13 and CAP17 supply the original theta, integral TypeForAlt and UnitaryHeckeData over the unchanged O. FNL, OST, WRP, OCT, OHA, SMP and ASP supply the original finite-level/HSet, ordinary/freeness/base-change, weight, control, Hecke-algebra and specialization fields. ADP and CFD give the actual automorphic and classical decompositions; AFS gives the original domain, DenseWt, TWtBridge, faithful fibers, scalar action and separation. LTP now gives all original local transports, actual admissible split components and one simultaneous family. BCS proves BaseChangeIsobaricN to DefiniteN to full BcSpeh under the original bc foundation, with all Speh lengths. GUA constructs the literal GalUnitaryT from the original RowT/PropsB foundation, with full IsRecT and bridge quantifiers, exact half twists, all required unramified places, pairing and conditional cusp output. UPA fills all sixteen PartAltAlt fields and all four point properties on one representation descending through the actual V.pr; it proves full partX\_\allowbreak alt/partU\_\allowbreak alt, installs AuTeichEx AutomorphicHeckeSet and the bare same-J TJDense endpoint. OFL and HEP give the original ordered ordinary flags and complete Out point record, and JRP gives every original RingProps field for that same IsJ ideal. The remaining arithmetic producers are the four exact Hida endpoints PointTheoremBaseObject, ArithmeticPointDecomposition, IrrArithPointsAutomorphicAway and TwistOk, followed by all-level same-ideal LevelAlt/TWAut/HasAut on both original Links. The explicit published base-change, polarized attachment, local representation and period foundations retain their original scope. The selected Raw patching, normalization and descent are no longer open producers. Preserve full CohP, LocCFT, N4, exact determinant/pro-ell conditions, nontrivial residual cyclotomic character, ell=3, fixed iota and one common coefficient threshold. No arbitrary-Kap B1/B2 theorem, Sel1=0 assertion, arbitrary-chain Core or complete same-J arithmetic support producer is claimed.

\emph{New completed inputs.} LTP1–LTP21, BCS1–BCS19, GUA1–GUA18, UPA1–UPA18, OFL1–OFL16, HEP full Out point endpoint, JRP full same-J RingProps

\emph{Current remaining subfrontier.} The four exact Hida endpoints PointTheoremBaseObject, ArithmeticPointDecomposition, IrrArithPointsAutomorphicAway and TwistOk, followed by all-level same-J automorphy on both original Links. The original BaseChangeIsobaricN, RowT/PropsB, local representation and period foundations retain their explicit foundation scope.

\emph{Source.} \href{https://github.com/CocoSgt/riemann-hypothesis-lean/blob/7cbf7a0c32fca1badf05931879a37c3f8a2b03b4/src/Development/AutomorphicArithmetic/AutomorphicVarieties/WeakShimuraTheory/IsAlgClosureField.lean#L19-L52}{Link}.

\emph{Source.} \href{https://github.com/CocoSgt/riemann-hypothesis-lean/blob/7cbf7a0c32fca1badf05931879a37c3f8a2b03b4/src/Development/AutomorphicArithmetic/AutomorphicVarieties/ModularityArgument/ThresholdArguments/DenseThresholdCases/PointTheoremBaseObject.lean#L32-L38}{HidaPropsAlt}.

\emph{Source.} \href{https://github.com/CocoSgt/riemann-hypothesis-lean/blob/7cbf7a0c32fca1badf05931879a37c3f8a2b03b4/src/Development/AutomorphicArithmetic/AutomorphicVarieties/WeilGroupCharacters/UniversalOrdinaryDeterminantRing.lean#L284-L286}{P1core}.



\paragraph{13. TASK\_\allowbreak CASE3\_\allowbreak AUT\_\allowbreak JOIN — B — Complete the original Case-III automorphy/image join}

Install the actual Case-III ProModAlt, SeedIVAlt and shared PMExAlt with the inherited CohomologyHyp/Ferrero-Washington/HFit/SeedII facts into the original whole Case-III automorphy/image conclusion.

\emph{Source.} \href{https://github.com/CocoSgt/riemann-hypothesis-lean/blob/7cbf7a0c32fca1badf05931879a37c3f8a2b03b4/src/Development/AutomorphicArithmetic/AutomorphicVarieties/ModularityArgument/ThresholdArguments/DenseThresholdCases/PromodularityCaseThree.lean#L43-L44}{AutomorphyRankFourDense}.

\emph{Source.} \href{https://github.com/CocoSgt/riemann-hypothesis-lean/blob/7cbf7a0c32fca1badf05931879a37c3f8a2b03b4/src/Development/AutomorphicArithmetic/AutomorphicVarieties/ModularityArgument/ThresholdArguments/DenseThresholdCases/CuFinal.lean#L44-L44}{main4\_\allowbreak alt}.



\paragraph{14. TASK\_\allowbreak CASE3\_\allowbreak PROMOD\_\allowbreak JOIN — C — Reconstruct the actual Case-III chain-to-promodularity proof}

Use the original tensor-factor tower, ChainRankFour, ChainDesc4, seeded rank-four chain reduction, actual PMExAlt and SeedIVAlt to construct ProModAlt on every original Ctx3 star point at the same fixed iota, preserving all finite-coefficient, inertia and support maps.

\emph{Source.} \href{https://github.com/CocoSgt/riemann-hypothesis-lean/blob/7cbf7a0c32fca1badf05931879a37c3f8a2b03b4/src/Development/AutomorphicArithmetic/AutomorphicVarieties/ModularityArgument/ThresholdArguments/DenseThresholdCases/PromodularityCaseThreeOfLemmas.lean#L131-L153}{UpperTheory.V.Cu.promodularity\_\allowbreak case\_\allowbreak three\_\allowbreak of\_\allowbreak chain\_\allowbreak reduction}.

\emph{Source.} \href{https://github.com/CocoSgt/riemann-hypothesis-lean/blob/7cbf7a0c32fca1badf05931879a37c3f8a2b03b4/src/Development/AutomorphicArithmetic/AutomorphicVarieties/ModularityArgument/ThresholdArguments/DenseThresholdCases/PromodularityCaseThree.lean#L40-L42}{PromodularityCaseThreeDense}.



\paragraph{15. SEED\_\allowbreak IV\_\allowbreak ALT — C — Actual Case-III seed with prechosen residual datum}

Base Ctx3/Four tower, iota and Amb -> w,s before every later Dies tower -> compatible finite-coefficient seed with exact residual support/root entries.

\emph{Source.} \href{https://github.com/CocoSgt/riemann-hypothesis-lean/blob/7cbf7a0c32fca1badf05931879a37c3f8a2b03b4/src/Development/AutomorphicArithmetic/AutomorphicVarieties/ModularityArgument/ThresholdArguments/DenseThresholdCases/PromodularityCaseThree.lean#L19-L25}{SeedIVAlt}.



\paragraph{16. TENSOR\_\allowbreak FACTOR\_\allowbreak TOWER — B — Actual tensor-factor tower and characters}

Original finite-field, tame, factor and character constraints; source-level construction is sizeable and has not been reconstructed here.

\emph{Source.} \href{https://github.com/CocoSgt/riemann-hypothesis-lean/blob/7cbf7a0c32fca1badf05931879a37c3f8a2b03b4/src/Development/AutomorphicArithmetic/AutomorphicVarieties/ModularityArgument/ThresholdArguments/DenseThresholdCases/PromodularityCaseThreeOfLemmas.lean#L131-L153}{UpperTheory.V.Cu.promodularity\_\allowbreak case\_\allowbreak three\_\allowbreak of\_\allowbreak chain\_\allowbreak reduction}.



\paragraph{17. CHAIN\_\allowbreak RANK\_\allowbreak FOUR — C — Actual SomeIn-to-AllIn propagation}

Same source Space and PM ideal, residual point Good4, original bound 30(C+1)\textasciicircum{}2 <= d; geometry must preserve all identities.

\emph{Source.} \href{https://github.com/CocoSgt/riemann-hypothesis-lean/blob/7cbf7a0c32fca1badf05931879a37c3f8a2b03b4/src/Development/AutomorphicArithmetic/AutomorphicVarieties/ModularityArgument/ThresholdArguments/DenseThresholdCases/PromodularityCaseThree.lean#L26-L32}{ChainRankFour}.



\paragraph{18. CHAIN\_\allowbreak DESC4 — B — Original residual chain descent}

Finite residual descent retaining source Chain4 compatibility.

\emph{Source.} \href{https://github.com/CocoSgt/riemann-hypothesis-lean/blob/7cbf7a0c32fca1badf05931879a37c3f8a2b03b4/src/Development/AutomorphicArithmetic/AutomorphicVarieties/ModularityArgument/ThresholdArguments/DenseThresholdCases/PromodularityCaseThreeOfLemmas.lean#L131-L153}{UpperTheory.V.Cu.promodularity\_\allowbreak case\_\allowbreak three\_\allowbreak of\_\allowbreak chain\_\allowbreak reduction}.



\paragraph{19. RANK4\_\allowbreak CHAIN\_\allowbreak REDUCTION — C — Seeded reduction to the original residual chain}

Select T1 before every later Deep tower; retain one compatible seed and Chain4 for every original residual lift over Q.

\emph{Source.} \href{https://github.com/CocoSgt/riemann-hypothesis-lean/blob/7cbf7a0c32fca1badf05931879a37c3f8a2b03b4/src/Development/AutomorphicArithmetic/AutomorphicVarieties/ModularityArgument/ThresholdArguments/DenseThresholdCases/PromodularityCaseThree.lean#L33-L39}{RankFourChainReductionDense}.

\emph{Source.} \href{https://github.com/CocoSgt/riemann-hypothesis-lean/blob/7cbf7a0c32fca1badf05931879a37c3f8a2b03b4/src/Development/AutomorphicArithmetic/AutomorphicVarieties/ModularityArgument/ThresholdArguments/DenseThresholdCases/CuPChainIV.lean#L16-L82}{UpperTheory.V.Cu.chainIV\_\allowbreak alt}.



\paragraph{20. TASK\_\allowbreak COUNT\_\allowbreak V\_\allowbreak INSTALL — B — Complete the original CountV finite assembly}

From exactly the listed source host/cohomology inputs, complete countV\_\allowbreak of\_\allowbreak finiteness\_\allowbreak duality\_\allowbreak euler\_\allowbreak host, including original host/frame formation, normalized length, the CaseV guard and R1<=2\textasciicircum{}c2.

\emph{Source.} \href{https://github.com/CocoSgt/riemann-hypothesis-lean/blob/7cbf7a0c32fca1badf05931879a37c3f8a2b03b4/src/Development/AutomorphicArithmetic/AutomorphicVarieties/HeightAndVanishing/Counts.lean#L69-L69}{CountV}.

\emph{Source.} \href{https://github.com/CocoSgt/riemann-hypothesis-lean/blob/7cbf7a0c32fca1badf05931879a37c3f8a2b03b4/src/Development/AutomorphicArithmetic/AutomorphicVarieties/ModularityArgument/ThresholdArguments/DenseThresholdCases/CountVOf.lean#L33-L60}{UpperTheory.V.UnitaryThreeInertia.countV\_\allowbreak of\_\allowbreak finiteness\_\allowbreak duality\_\allowbreak euler\_\allowbreak host}.



\paragraph{21. THM\_\allowbreak S6\_\allowbreak S — C — Original unitary-three ordinary support and automorphy}

Original host/frame and fixed iota -> one marked g1 -> every universal R,u -> weight-zero finiteness and every Type-I BPrime\_\allowbreak s point.

\emph{Source.} \href{https://github.com/CocoSgt/riemann-hypothesis-lean/blob/7cbf7a0c32fca1badf05931879a37c3f8a2b03b4/src/Development/AutomorphicArithmetic/AutomorphicVarieties/ModularityArgument/ThresholdArguments/DenseThresholdCases/ThmS6CoreS.lean#L17-L26}{ThmS6\_\allowbreak s}.

\emph{Source.} \href{https://github.com/CocoSgt/riemann-hypothesis-lean/blob/7cbf7a0c32fca1badf05931879a37c3f8a2b03b4/src/Development/AutomorphicArithmetic/AutomorphicVarieties/ModularityArgument/ThresholdArguments/DenseThresholdCases/ThmS6CoreS.lean#L12-L16}{BPrime\_\allowbreak s}.



\paragraph{22. HOST\_\allowbreak FRAME\_\allowbreak LENGTH — B — Original Case-V host/frame and normalized-length installation}

Actual HostData, frame, local-condition equality, abstract count and normalized original Selmer length must all refer to the same objects.

\emph{Source.} \href{https://github.com/CocoSgt/riemann-hypothesis-lean/blob/7cbf7a0c32fca1badf05931879a37c3f8a2b03b4/src/Development/AutomorphicArithmetic/AutomorphicVarieties/ModularityArgument/ThresholdArguments/DenseThresholdCases/CountVOf.lean#L33-L60}{UpperTheory.V.UnitaryThreeInertia.countV\_\allowbreak of\_\allowbreak finiteness\_\allowbreak duality\_\allowbreak euler\_\allowbreak host}.



\paragraph{23. HOST\_\allowbreak NUMBER\_\allowbreak S — C — Uniform original-host rank-three catalogue}

One exponent before all closures, coefficient rings, levels, characters, frames, embeddings and marked points; one separating trace function controls all later finite pairwise nonisomorphic collections.

\emph{Source.} \href{https://github.com/CocoSgt/riemann-hypothesis-lean/blob/7cbf7a0c32fca1badf05931879a37c3f8a2b03b4/src/Development/AutomorphicArithmetic/AutomorphicVarieties/ModularityArgument/ThresholdArguments/DenseThresholdCases/HostNumberS.lean#L18-L28}{HostNumber\_\allowbreak s}.



\paragraph{24. TASK\_\allowbreak HOST\_\allowbreak VAN\_\allowbreak E\_\allowbreak INSTALL — B — Transport K/duality to literal HostVanE}

Use the already supplied LemmaK/local-global duality and finiteness on the actual CountV host/frame and exact original coefficient modules.

\emph{Source.} \href{https://github.com/CocoSgt/riemann-hypothesis-lean/blob/7cbf7a0c32fca1badf05931879a37c3f8a2b03b4/src/Development/AutomorphicArithmetic/AutomorphicVarieties/UnitaryThreeDeformations/HostNumber.lean#L9-L12}{HostVanE}.

\emph{Source.} \href{https://github.com/CocoSgt/riemann-hypothesis-lean/blob/7cbf7a0c32fca1badf05931879a37c3f8a2b03b4/src/Development/AutomorphicArithmetic/AutomorphicVarieties/LevelCohomology/SoHostVanE.lean#L11-L18}{UpperTheory.V.X.So.hostVanE}.



\paragraph{25. TASK\_\allowbreak HOST\_\allowbreak XPP\_\allowbreak INSTALL — B — Transport CorX to literal HostXpp}

Preserve Psi\textasciicircum{}-2 eta, the prime-to-ell quotient, actual local conditions and the original CountV host while installing BR7.

\emph{Source.} \href{https://github.com/CocoSgt/riemann-hypothesis-lean/blob/7cbf7a0c32fca1badf05931879a37c3f8a2b03b4/src/Development/AutomorphicArithmetic/AutomorphicVarieties/UnitaryThreeDeformations/HostNumber.lean#L14-L17}{HostXpp}.

\emph{Source.} \href{https://github.com/CocoSgt/riemann-hypothesis-lean/blob/7cbf7a0c32fca1badf05931879a37c3f8a2b03b4/src/Development/AutomorphicArithmetic/AutomorphicVarieties/LevelCohomology/Xpp.lean#L11-L84}{UpperTheory.V.X.xpp}.



\paragraph{26. DESIGN\_\allowbreak CHOICES\_\allowbreak COMMON — C — Construct the original common design family and bundle}

For the one target HostExW, construct the original compatible DesignChoices on Cdes=WithTw(WithSm Cprime), where Cprime=WithChN(A intersect target class), A is the one Built construction class, and cT(eps,y)=cTmin(y). Retain every original field, tuple, word, metric and common bundle duty. DSZ10 and RMP11 give size/rank properties of class members, not member existence. The original RecExP/RecAtG and common object/row route remains coupled. Construct ONE original common design family R and its recursive/geometric half with original Inputs/Pins/Holds/ExAt/InGeo and required CoreP. PKT now reconstructs full exists\_\allowbreak pkT, rowPk\_\allowbreak nonempty and exists\_\allowbreak moved with Copy/Coh on those original given inputs and a supplied fitting later stage: one catalogue, both ChT, complete LfT, actual integral lattices and all metric/measure/period constants are retained; one numerical record precedes every eta and V/S/N/G. RDR supplies one Side0 before all readers, ReadPk/pins and full ReadL with exact o.eta=q.eta. The initial actual geometry, Fits/division/central conditions, primitive rational tests and eta agreement must still be produced. Inherited CSR supplies finite columns, one ReadC before later twists, period/evaluation and scalar/error transports, RowSideDetermined and same-R Bundle/BTw/BWit/BunU consumers. RPI supplies exact RowPairInt for actual pairs and finite translates with original mu\_\allowbreak b/fd0; first literal DivDef3 at the same updated R,F,Y,Z,C,sEE follows from its retained original readings, Pins/Holds, stored nonzero IsItW chain, five real-group predicates and group tie. Remaining common fields include host/frame, MultEq nonzero-part support, that original nonzero chain/group tie, eigensystem/carrier/parts and all-embedding logLam, HD/d3b, NEb, actual rational tests and base rows. Retain every simultaneous family hypothesis and eventual threshold. Stronger second DivRows3 remains separate; universal quotient-height CuspRapidDecay and actual common member existence are not inferred from the comparison constructors. Version 0.41 supplies the finite phase-correct rational theta route: exact coefficient extension, split representatives, square-law/PhiIs, Pins from CoversAt, class-sum/point-identity transfer to OneMapOverL, actual-good-projector RatG and both-level LiftTestB. Its original higher covers, all-tuple class-sum uniqueness including phase rescaling, one PointIdentCore before O, base ClassSumsSpan and geometric LatFull remain explicit producers. MLat is available from the original LatFull input. This sharpens the rational-test work without claiming actual common geometry, Fits or eta agreement.



\paragraph{27. LOCAL\_\allowbreak NBUILT\_\allowbreak COMMON — C — Actual local specialization on the common selected family}

Construct, with the required closed-class property on the same Des(C,3), one eventual family of actual NBuilt certificates for the original Fm.D(x,i), retaining the certificate before forgetting to TSel. HLP now reconstructs full HeckeLift0, all-index commutation, HeckeStd0 and HeckeIntTw from the original MtInput and complex Hecke facts, in the full ambient GenAt range k<=nV. Every maximal ideal is an actual supported residual character class; literal MtSub is the O-span of projections of the original MtLat generators, is finite free, and localizes onto the exact NCov.part. Literal Psi has its exact range and off-space projection error, both original span assertions hold, and ActEx/ActHK are constructed on the same lattice under original semisimplicity. Original MjLat relation quotients, universal coefficient-map relation fields, selected tame generators, coefficient embeddings and all identity transports give the complete nonempty\_\allowbreak spec constructor on its full printed input list. A single compositum combines the four original coefficient requirements with big fixed beforehand. Whole ndSpec\_\allowbreak exists still needs actual common geometric/choice MtInput, complex HeckeStdPlain and the no-Jordan producer, original F3 bridges with valuation/GoodCoef/iq, original F4 common-overfield reading of the entire coefficient field, and exact arithmetic character/reciprocity/admissibility suppliers. Original AutSpec, Galois specialization and the remaining NBuilt fields on that same installation also remain. Do not replace the actual projected lattice or maps by abstract objects with the same dimensions.



\paragraph{28. SECOND\_\allowbreak LAYER\_\allowbreak COMPATIBLE\_\allowbreak SELECTION — C — Original second-layer selector and class duties}

For the same A, supply Duty on A and a Closed common original TSel family whose selected records are the original true-column rec7 records (or have the entire dependent tuple transported), with their real witness and original coefficient-field root data. Existence/compatibility remains arithmetic; merely abstractly isomorphic fields or independently selected columns do not suffice. CSR21–22 now construct the same-R Bundle/BTw/BWit/BunU outputs from the exact selected realization/row-side packet and NEb, with the original EvP/ParEv quantifier order. Actual common selection, the five stated realization packets and NEb remain upstream; none is replaced by an independently chosen compatible record. Version 0.38 closes the selected analytic extension and reconstructs the first updated DivDef3 consumer; it leaves the original common selection, geometric readings, rational tests, base rows, host/parts/NEb and exact retained DI premises upstream, with unchanged EvP/ParEv thresholds. Version 0.39 supplies the original given-data reader constructors and exact RowPairInt/first DivDef3 on original premises. It leaves one actual common geometric/comparison/test member, eta agreement, compatible selection, rational tests, base rows, host/parts/NEb and original EvP/ParEv thresholds upstream. Version 0.40 also supplies original comparison/RowPk/fitting-stage Moved algorithms and exact common data transport. The actual initial geometric member, fitting tests, eta agreement and compatible second-layer selector with original eventual thresholds remain upstream. Version 0.41 supplies the finite phase-correct rational theta route: exact coefficient extension, split representatives, square-law/PhiIs, Pins from CoversAt, class-sum/point-identity transfer to OneMapOverL, actual-good-projector RatG and both-level LiftTestB. Its original higher covers, all-tuple class-sum uniqueness including phase rescaling, one PointIdentCore before O, base ClassSumsSpan and geometric LatFull remain explicit producers. MLat is available from the original LatFull input. This sharpens the rational-test work without claiming actual common geometry, Fits or eta agreement.



\paragraph{29. BASE\_\allowbreak THETA\_\allowbreak READINGS — C — Actual base reading on the common bundle}

For every family in the original class, eventually on the SAME selected Bundle and for each i<m, supply the actual R,F,Read,ReadL, weak models/Weil data and phase-one pins. Identify their LiftTie.theta, mtRead, mjRead and ThDef3 basis with that Bundle.P.A3, P.Zs, P.S0 and its stored maps. Exact field, embeddings, metrics, measures, row family and weak choices are retained. This remains one part of the coupled BTH12 producer. PKT now supplies the complete original comparison packet, RowPk and fitting-stage Moved with Copy/Coh on original given geometry, Inputs/Pins/Holds/ExAt/InGeo and Fits. Its ChT/LfT/Reps, actual lattice carriers, original Haar and norm constants, period maps, representatives and all-place RrT are constructed together; its one numerical record precedes every eta/V/S/N/G. RDR gives the original level-zero side, moved-Read and ReadPk/pin constructors on those objects, with one Side0 before every reader and all twelve Read ties. Full ReadL retains the exact o.eta(j,v)=q.eta(j,v) guard for every 1<=j<=m and place v; arbitrary Over does not imply it. Inherited CSR fixes one ReadC before all TwArgs, transports both metrics and all thirteen reading maps, and gives actual period/evaluation with kappa=(c*nP)\textasciicircum{}(-1/2). RPI supplies exact row-pair integrability, so first updated DivDef3 follows from the retained original Pins/Holds, nonzero chain, five real-group predicates and group tie. Initial theta extension remains complete. Remaining work is actual initial common geometry and its original predicates, Fits/primitive rational tests, eta agreement and outer family/eventual/same-Bundle duties, reFace compatibility, MultEq support, original chain/group data, actual base rows and any stronger DivRows3 theta square. No new compatible reading, family threshold or Haar normalization is substituted. Version 0.41 supplies the finite phase-correct rational theta route: exact coefficient extension, split representatives, square-law/PhiIs, Pins from CoversAt, class-sum/point-identity transfer to OneMapOverL, actual-good-projector RatG and both-level LiftTestB. Its original higher covers, all-tuple class-sum uniqueness including phase rescaling, one PointIdentCore before O, base ClassSumsSpan and geometric LatFull remain explicit producers. MLat is available from the original LatFull input. This sharpens the rational-test work without claiming actual common geometry, Fits or eta agreement.



\paragraph{30. BASE\_\allowbreak THETA\_\allowbreak RATIONAL\_\allowbreak TESTS — C — Simultaneous original base rational tests}

On the SAME original Bundle/reading/theta map, for every i<m supply one finite separating complex-linear test family whose values lie in the fixed embedded coefficient field on theta of EVERY original rational source vector and on EVERY stored target-coordinate basis vector. BTH8/BTH11 use this same LiftTestB for rational adjoint image and matrix coordinates. RTG now reconstructs prescribed-field coefficient extension, original split class representatives/RepsOK, the self-dual split-lattice square law and finite residue PhiIs; every Pins field follows from the actual CoversAt. With the base coefficient record literally fixed, exact class-sum existence/uniqueness including phase-rescaled O, PointIdentCore with common lambda,z before every O, and base ClassSumsSpan imply all-tuple finite lift span and OneMapOverL. The phase is exactly cj(twAt embLL A3.A.v0)(R.phs i). RatG uses the actual good projector eOp(y), and full LiftTestB follows at BOTH original levels. These finite implication proofs are complete. Remaining original producers are the prescribed higher covers, those actual class sums/uniqueness, PointIdentCore, base ClassSumsSpan and the base/common geometric datum; exact MLat is supplied only from original LatFull, whose geometric production remains separate. Finite dimension or an unrelated rational theta model cannot replace these same-object identities. Original outer eventual-family/Bundle quantifiers remain.



\paragraph{31. SLR\_\allowbreak BASE\_\allowbreak ROWS\_\allowbreak PRODUCER — C — Actual base row identities on that same bundle}

For the SAME selected Bundle, common original fields and readings used by the other three BTH12 inputs, supply actual RowsHold0\_\allowbreak 3 uniformly eventually for every family. Its original primitive row realization, spectrum/type and basis identifications must be retained. BTH supplies a theta consumer only; it does not construct the base row identity. No separate witness may be chosen merely to satisfy this row requirement.



\clearpage
\section{Local components and exact transport on the original automorphic part}
\label{sec:local-transport-v44}

\status{This chapter proves the seven original local-transport
suppliers \texttt{part\allowbreak UR2Iso\allowbreak L}, \texttt{part\allowbreak UR2Iso\allowbreak Q},
\texttt{part\allowbreak UR2Sph}, \texttt{part\allowbreak UR2Unr}, \texttt{part\allowbreak UR2Wt},
\texttt{part\allowbreak UR2V0}, and \texttt{part\allowbreak UR2Type\_\allowbreak alt}.  It also constructs
an admissible local component at every split place, with the exact
unramified--spherical equivalence, on the original automorphic
function space.  The scalarization and the integral type are the
ones already reconstructed in ADP and ITA.  No Galois attachment,
base-change theorem, or completed \texttt{Part\allowbreak Alt\allowbreak Alt} record is
asserted in this chapter.}

\subsection{Objects, conventions, and reusable elementary arguments}

Fix all the original data
\[
 \ell,\ L,\ \Omega,\ s,\ n\ge1,\ V,\ Q,\ c,\ \iota,\ m,\ \tau_0,\
       \lambda,\ d,\ E,\ \xi.
\]
Thus \(L\) is the original CM field, \(\Omega/L\) its specified
algebraic closure, \(s\) the original Hida setting, \(V\) its original
level datum, and \(Q\) its auxiliary-prime datum.  The coefficient DVR
is \(O=c.O\), and the fixed field isomorphism is
\(\iota:\overline{\mathbf Q}_\ell\longrightarrow\mathbf C\).
The integral representation
\(\tau_0:K_T\longrightarrow\operatorname{GL}_m(O)\) is arbitrary
unless a type hypothesis is expressly imposed.  No open-kernel
hypothesis is added to the first six endpoints.

Put \(H=\mathrm{hset}(Q,\tau_0)\),
\(G=U_{\mathbb A}(L,n)\), \(G_f=U_f(L,n)\), and
\(G_\infty=U_\infty(L,n)\).  The group is the identity-form unitary
group of FNL2.  Right translation has the convention
\[
              (\mathscr R_x\phi)(g)=\phi(gx),
              \qquad \mathscr R_x\mathscr R_y=\mathscr R_{xy}.
\]
Let \(I=\operatorname{Emb}_\ell((s.w(v))_{v\in V.\mathrm{frame}.V_\ell})\),
and use exactly the prescribed basis equivalence
\[
 D=\bigotimes_{a\in I}\nabla_{\lambda(a.\mathrm{res})}(\mathbf C),
       \qquad E:\mathbf C^d\xrightarrow{\sim}D .
\]
The source's \(\mathrm{IsBasis}(\iota,E,\xi)\) is assumed.  Write
\(R=\mathrm{infRep}(\iota,\lambda)\) and retain the original projection
\[
 p(g)=\bigl(\mathrm{away}(g),
                 (\mathrm{cLoc}_{s.w(v)}(g))_{v\in V.\mathrm{frame}.V_\ell}\bigr)
             \in G^a\times G^\ell .
\]
For a coefficient function \(f:G^a\times G^\ell\to\mathbf C^d\),
its scalarization is
\begin{equation}\label{eq:ltp-scalarization}
       \mathcal T_f(\alpha)(g)
       =\alpha\bigl(R(g^{-1})E(f(p(g)))\bigr),
                   \qquad\alpha\in D^* .
\end{equation}
For \(F=(F_i)_{i<m}\), use this formula on each \(F_i\).
Let \(\sigma\) be an original \(\mathrm{IsAutRepU}\) subspace of
the actual functions \(G\to\mathbf C\), and set
\(Y=\mathrm{YOf}(Q,\iota,\tau_0,E,\xi,\sigma)\).

The inherited statements used here have written proofs: FNL2--FNL6
and FNL16 for the actual groups and sections; ADP1--ADP3 and
ADP6--ADP11 for the positive form, finite multiplicity and
scalarization; HP1--HP4 for the characteristic-zero polynomial
representations; and ITA9--ITA12 for the supplied integral type.
The local-transport arguments introduce no new local Langlands,
base-change or Galois-attachment premise.  When ITA12 is used, its
expressly recorded classical type and correspondence foundations
remain at exactly their previous scope.  In particular we do not
turn a supplied \(\mathrm{IsTypeAlt}\) into an unproved assertion
about all local representations.

% BEGIN LTP1
\paragraph{LTP1. The literal part and its finite separating scalarizations.}
\label{claim:LTP1}
For every coefficient vector \(F\),
\[
 F\in Y
 \quad\Longleftrightarrow\quad
 F\in H.\mathrm{Big}(\iota,\xi)
       \ \hbox{and}\ 
 \mathcal T_{F_i}(\alpha)\in\sigma
                  \quad\hbox{for every }i,\alpha .
\]
If \(\alpha_1,\ldots,\alpha_d\) is a basis of \(D^*\), the finitely
many linear maps
\(F\mapsto\mathcal T_{F_i}(\alpha_b)\), indexed by \(i<m,b\le d\),
separate coefficient vectors.  These assertions include \(m=0\)
and \(d=0\).
\begin{proof}
Scalarization is complex linear in \(F_i\) and in \(\alpha\).
Consequently the set of vectors for which all scalarizations belong
to the submodule \(\sigma\) is itself a submodule.  Its span, which
occurs in the literal definition of \(\mathrm{YOf}\), is therefore
that same set.  Intersecting with \(H.\mathrm{Big}\) proves the
equivalence, rather than just one implication for generators.

If all displayed coordinates vanish, linearity in \(\alpha\) makes
every scalarization vanish.  ADP7 supplies injectivity on the original
functions: \(p|_{G_f}\) is surjective, \(R(g)\) is invertible, \(E\)
is an isomorphism, and the full dual separates \(D\).  Thus every
\(F_i=0\).  When an index set is empty the corresponding coefficient
space is zero, so the same assertion holds.  This proves the full
source \texttt{part\allowbreak UR2Mem} and the finite separating family needed
below.
\end{proof}
% END LTP1

% BEGIN LTP2
\paragraph{LTP2. Actual local factors, centralizers, and the opposite place.}
\label{claim:LTP2}
For a split finite place \(w\) of \(L\), there is a continuous
homomorphism
\[
       j_w:\operatorname{GL}_n(L_w)\longrightarrow G_f
\]
whose image is exactly \(\mathrm{locU}(n,w)\), with
\(\mathrm{cLoc}_w j_w(g)=g\).  It has the actual finite coordinates
specified in FNL4--FNL6.  Every \(x\in G\) factors as
\[
       x=j_w(\mathrm{cLoc}_w x)\,x^{(w)},
             \qquad x^{(w)}j_w(g)=j_w(g)x^{(w)}
                                     \quad\hbox{for every }g.
\]
The factor \(x^{(w)}\) has identity coordinates at both places above
\(w|_{L^+}\).  Distinct base places give commuting local factors.
If \(w'\) has the same base place as \(w\), then
\[
 \mathrm{cLoc}_w x\in\operatorname{GL}_n(\mathcal O_{L_w})
 \quad\Longrightarrow\quad
 \mathrm{cLoc}_{w'}x\in\operatorname{GL}_n(\mathcal O_{L_{w'}}).
\]
This last implication does not require that the base place split.
\begin{proof}
At a split pair \(w,w^c\), set the \(w\)-coordinate equal to \(g\),
the other equal to its inverse conjugate transpose, and every
remaining finite and infinite coordinate equal to one.  The unitary
identity verifies membership.  The inverse-transpose operation is
a group homomorphism in this order:
\((gh)^{-t}=g^{-t}h^{-t}\).  Thus the coordinate construction is a
homomorphism.  It is continuous because the two nonconstant
coordinate maps are continuous.  Projection at \(w\) is its inverse
on \(\mathrm{locU}(n,w)\): the local-support condition removes all
other base places and the unitary equation determines the opposite
coordinate.  This is precisely the section already checked in FNL.

Removing these two coordinates from \(x\) gives \(x^{(w)}\).
Coordinate multiplication proves the displayed factorization and
commutation.  It also proves commutation for sections supported at
distinct base places.  Thus the centralizer assertion is about the
actual full adelic group, including its infinite coordinates.

In a quadratic CM extension, two places over one finite base place
are equal or are interchanged by conjugation.  In the second case
the unitary equation gives
\[
 \mathrm{cLoc}_{w'}x
       ={}^t\!\bigl(c_{w,w'}(\mathrm{cLoc}_w x)^{-1}\bigr).
\]
Local conjugation carries the valuation ring onto the valuation
ring.  Membership in \(\operatorname{GL}_n(\mathcal O)\) says that
both a matrix and its inverse have integral entries.  The displayed
formula exchanges those two conditions and preserves them.
This proves the last assertion, including the inert case when the
places coincide.  The source's \(KGL\) and \(KO\) both express this
same lattice-preservation condition.
\end{proof}
% END LTP2

% BEGIN LTP3
\paragraph{LTP3. A local component makes the whole restriction isotypic.}
\label{claim:LTP3}
Suppose
\[
 \mathrm{IsCompAlong}\!
 \left(\sigma,\mathrm{locU}(n,w),\mathrm{cLoc}_w,\tau\right)
\]
for a smooth irreducible \(\tau\) at a split place \(w\).
Then the restriction of the full \(\sigma\) along \(j_w\) is
\(\tau\)-isotypic in the literal \(\mathrm{IsoSum}\) sense.
\begin{proof}
The component hypothesis supplies a nonzero equivariant linear map
from \(\tau\) into the original function subspace \(\sigma\).  It
intertwines the action along \(j_w\), by LTP2.  Let \(Z\) be the sum
of the images of all such local intertwiners.  This sum is nonzero.
Each image is stable under \(j_w(\operatorname{GL}_n(L_w))\).
An element \(h\) centralizing that local subgroup sends an intertwiner
\(f\) to the intertwiner \(\mathscr R_hf\).  Therefore it preserves
\(Z\).  The factorization in LTP2 now shows that every \(x\in G\)
preserves \(Z\).  The original global irreducibility of \(\sigma\)
gives \(Z=\sigma\).  This equality is precisely \(\mathrm{IsoSum}\);
no irreducibility of the whole restriction is asserted.
\end{proof}
% END LTP3

% BEGIN LTP4
\paragraph{LTP4. Finite equivariant separating families preserve isotypy.}
\label{claim:LTP4}
Let \(Y\) be a stable subspace of an arbitrary representation
\(\rho\) of a group \(J\).  Suppose \(Z\) is a \(\tau\)-isotypic
subspace of another representation, and finitely many linear maps
\(\phi_a\) take \(Y\) into \(Z\), intertwine the \(J\)-actions on
\(Y\), and jointly have zero kernel on \(Y\).  Then \(Y\) is
\(\tau\)-isotypic.
\begin{proof}
The combined map embeds \(Y\) equivariantly in \(Z^A\).
An isotypic module, defined as the sum of images of maps from a
simple module \(\tau\), is a sum of simple submodules of that type.
The elementary module argument of ITA10 makes it a direct sum of
copies of \(\tau\), and every submodule of that sum is again a sum
of copies of \(\tau\).  Since \(A\) is finite, \(Z^A\) is a finite
direct sum of such modules, hence has that same property.  Apply
it to the image of \(Y\), and transport the embeddings of the
simple summands back through the injective map.
All statements are algebraic; neither admissibility nor finite
dimension of \(Y,Z,\tau\) is used.
\end{proof}
% END LTP4

% BEGIN LTP5
\paragraph{LTP5. Compact fixed vectors and the exact spherical sums.}
\label{claim:LTP5}
Let \(Z\) be a \(\tau\)-isotypic subspace of a representation
\(\rho\) of \(\operatorname{GL}_n(F)\), where \(F\) is a
nonarchimedean local field and \(\tau\) is smooth irreducible.
For every compact subgroup \(K\),
\[
               Z^K\ne0\quad\Longrightarrow\quad\tau^K\ne0 .
\]
Put \(K_0=\operatorname{GL}_n(\mathcal O_F)\).  If the source
operator \(\mathrm{sphOp}_{\tau}(j)\) equals \(a\) on \(\tau^{K_0}\),
then \(\mathrm{sphOp}_{\rho}(j)\) equals \(a\) on \(Z^{K_0}\),
for every \(j\in\mathbf N\).  Every additive equivariant map
commutes with the source spherical operator.  Finally each original
\(K_c(F,n,c)\), for \(c\ge0\), is compact.
\begin{proof}
For a smooth vector \(v\), its \(K\)-orbit is finite: the stabilizer
is open and its cosets cover the compact group.  Averaging over the
finite left cosets of this stabilizer gives a \(K\)-fixed vector
\(e_Kv\).  Refining stabilizers shows that this is a linear operation,
fixes every \(K\)-fixed vector, and commutes with equivariant linear
maps.  Normality of a stabilizer is unnecessary: left multiplication
permutes its left cosets.

Write \(z\in Z^K\) as a finite sum \(\sum_b f_b(v_b)\) of images of
intertwiners from \(\tau\).  Applying the average gives
\(z=\sum_bf_b(e_Kv_b)\).  If \(\tau^K=0\), this makes \(z=0\).
This proves the first assertion and gives the decomposition into
images of fixed vectors needed for the second.

The original spherical operator is the unnormalized double-coset
sum for
\[
 a_j=\operatorname{diag}
       (\underbrace{\varpi,\ldots,\varpi}_{\min(j,n)},
        1,\ldots,1).
\]
For the source's finite right-coset representatives
\(K_0a_jK_0=\coprod_b g_bK_0\), it is
\(\sum_b\rho(g_b)\).  Finiteness follows because
\(K_0\cap a_jK_0a_j^{-1}\) is open in compact \(K_0\);
this is the same finite-sum dictionary proved in OHA1.
An additive intertwiner commutes with this finite sum term by
term, even before one restricts to fixed vectors.  On the
decomposition \(z=\sum_bf_b(v_b)\) with \(v_b\in\tau^{K_0}\),
the assumed scalar equation therefore gives
\(\mathrm{sphOp}_\rho(j)z=az\).
The argument covers \(j=0\) and \(j\ge n\), with no change of
normalization or factor order.

The subgroup \(K_0\) is compact: identify it with the closed set of
pairs \((A,B)\) in two finite matrix powers of the compact ring
\(\mathcal O_F\) satisfying \(AB=BA=1\).
The original \(K_c(F,n,c)\) is the kernel of reduction modulo
\(\mathfrak m_F^c\), inside this group.  Reduction is continuous
with finite target, so the kernel is closed and compact.
For \(c=0\) the reduction condition is vacuous and the group is
\(K_0\).  These are all four clauses of the source
\texttt{part\allowbreak UR2Abs}.
\end{proof}
% END LTP5

\subsection{Transport of the two local actions}

% BEGIN LTP6
\paragraph{LTP6. The exact action at a place above \(\ell\).}
\label{claim:LTP6}
For \(v\in V.\mathrm{frame}.V_\ell\), use the actual split section
\(j_v=j_{s.w(v)}\).  For every coefficient vector \(F\), every
\(g\in\operatorname{GL}_n(L_{s.w(v)})\), every \(i,\alpha\),
\[
 \mathcal T_{(H.\mathrm{atV}(v,g)F)_i}(\alpha)
                 =\mathscr R_{j_v(g)}\mathcal T_{F_i}(\alpha).
\]
Both \(H.\mathrm{Big}(\iota,\xi)\) and \(Y\) are stable under this
action.
\begin{proof}
The chosen place lies above \(v\), and the splitting condition in
the Hida setting gives LTP2's section.  Its away coordinates are
one and its selected \(\ell\)-coordinates are one except at \(v\),
where the coordinate is \(g\).  Thus
\[
             p(xj_v(g))=p(x)(1,\delta_v(g)).
\]
Its infinite coordinates are one, hence \(R(j_v(g))=1\) by ADP6.
In ADP7's right-translation rule the functional
\(\alpha\circ R(j_v(g)^{-1})\) is therefore exactly \(\alpha\).
The result is the displayed identity, with \(xj_v(g)\) in that
order.  This proves the full source \texttt{part\allowbreak UR2At\allowbreak V}.

The local right translation commutes with rational left
translation.  It also commutes with the entire away action,
including its coefficient matrix, since the two group coordinates
are distinct and the matrix is constant.  Therefore it preserves
both defining clauses of \(H.\mathrm{Big}\).  Finally the
scalarization identity and global stability of \(\sigma\) preserve
all the scalarizations required by LTP1, proving stability of \(Y\).
\end{proof}
% END LTP6

% BEGIN LTP7
\paragraph{LTP7. Full original \texttt{part\allowbreak UR2Iso\allowbreak L}.}
\label{claim:LTP7}
For every \(v\in V.\mathrm{frame}.V_\ell\) and every supplied smooth
irreducible local component \(\tau\) of \(\sigma\) along
\(\mathrm{locU}(n,s.w(v))\),
\[
                  \mathrm{IsoSum}(H.\mathrm{atV}(v),Y,\tau).
\]
No nonzero-\(Y\) or open-\(\ker\tau_0\) assumption is required.
\begin{proof}
LTP3 makes the restriction of \(\sigma\) along \(j_v\)
\(\tau\)-isotypic.  LTP6 gives the stability and intertwining
conditions for the finite scalarization family of LTP1.
That family is separating and lands in \(\sigma\).
LTP4 proves the stated literal \(\mathrm{IsoSum}\).
\end{proof}
% END LTP7

% BEGIN LTP8
\paragraph{LTP8. The exact away action, including its coefficient matrix.}
\label{claim:LTP8}
Let \(w\) be split with \(w|_{L^+}\notin s.S\).  For every
coefficient vector \(F\), \(g\in\operatorname{GL}_n(L_w)\),
and \(i,\alpha\),
\[
 \mathcal T_{(H.\mathrm{atW}(w,g)F)_i}(\alpha)
                    =\mathscr R_{j_w(g)}\mathcal T_{F_i}(\alpha).
\]
At an original auxiliary place \(w=Q.wq(q)\), the actual
\(H.\mathrm{Ax}\) makes \(H.\mathrm{Big}\) and \(Y\) stable
under this action.
\begin{proof}
The place lies away from \(\ell\), the type places, and \(v_0\).
The away projection of \(j_w(g)\) belongs to the original
\(\mathrm{locA}(n,s,w)\).  Its \(w\)-coordinate is \(g\);
all other base coordinates are one.  The bijective coordinate
projection checked in FNL4--FNL6 therefore identifies that away
projection with the source's \(\mathrm{sec}(w,g)\), including
the branch of its conditional definition.  Furthermore
\[
 \mathrm{toKT}(\mathrm{sec}(w,g))=1,\qquad
 H.\tau_C(\mathrm{sec}(w,g))=1.
\]
Its selected \(\ell\)-coordinates and all infinite coordinates
are one.  Thus the original matrix-valued away action
\[
  (\mathrm{raF}(u)F)_i(y)
       =\sum_b\tau_C(u)_{ib}F_b(y(u,1))
\]
reduces at \(u=\mathrm{sec}(w,g)\) to the pure right translation
\(F_i(y(u,1))\).  ADP7 gives exactly the asserted scalarization
identity.  This is the full source \texttt{part\allowbreak UR2At\allowbreak W}.

For an auxiliary \(q\), the original axiom \(\mathrm{Q\_\allowbreak UQ}\)
says that \(\mathrm{sec}(Q.wq(q),g)\) commutes with every element
of \(H.UQ\).  Consequently its \(\mathrm{raF}\) action preserves
the \(UQ\)-fixed condition.  It also preserves rational coefficient
equivariance: every coordinate is a finite complex linear
combination of right translates, which commute with rational
left translation.  It preserves \(H.\mathrm{Big}\).
The scalarization identity and LTP1 then prove stability of \(Y\).
The axiom has been verified for this very \(H\) in FNL16.
\end{proof}
% END LTP8

% BEGIN LTP9
\paragraph{LTP9. Full original \texttt{part\allowbreak UR2Iso\allowbreak Q}.}
\label{claim:LTP9}
Under the original \(H.\mathrm{Ax}\), for every auxiliary
\(q\in Q.Q\) and every supplied smooth irreducible local component
\(\tau\) of \(\sigma\) at \(Q.wq(q)\),
\[
             \mathrm{IsoSum}(H.\mathrm{atW}(Q.wq(q)),Y,\tau).
\]
This includes \(Y=0\).
\begin{proof}
The auxiliary-prime conditions give
\(Q.wq(q)|_{L^+}\notin s.S\) and membership of this base place
in \(V.S'\); the original level datum says it splits.
Thus LTP2, LTP3 and LTP8 apply.  The same finite family from
LTP1 satisfies all the hypotheses of LTP4, giving the conclusion.
The original \(H.\mathrm{Ax}\) premise is retained in this generic
endpoint and is available from FNL16 in the actual application.
\end{proof}
% END LTP9

\subsection{The compact level equation and its three consequences}

% BEGIN LTP10
\paragraph{LTP10. Full matrix-valued transport of the base level equation.}
\label{claim:LTP10}
Let \(F\in H.\mathrm{Big}(\iota,\xi)\), and let \(x\in G_f\).
Assume its selected \(\ell\)-coordinates are one.  Let \(k\in K_T\)
have components \(k_u=\mathrm{cLoc}_{s.w(u)}x\) at every type
place \(u\in s.St\).  Suppose also that
\[
\begin{gathered}
 \mathrm{cLoc}_{s.w(v_0)}x\in K_c(L_{s.w(v_0)},n,1),\\
 \mathrm{cLoc}_{w}x\in KO_w
                  \quad(w|_{L^+}\notin V.S'),\qquad
 \mathrm{cLoc}_{Q.wq(q)}x=1\quad(q\in Q.Q).
\end{gathered}
\]
Writing \(t=\iota(\tau_0)\) for the original coefficient extension,
one has, for every \(i,\alpha\),
\begin{equation}\label{eq:ltp-level-matrix}
       \sum_{b<m} t(k)_{ib}\,
            \mathscr R_x\mathcal T_{F_b}(\alpha)
                      =\mathcal T_{F_i}(\alpha).
\end{equation}
This is the full source \texttt{part\allowbreak UR2Fix}.
\begin{proof}
Set \(u=\mathrm{away}(x)\).  The definition of \(H.U_b\) asks that
its type components be integral and its distinguished component
lie in \(K_c(1)\); these follow from the stated conditions.
The definition of \(H.UQ\) adds precisely integrality outside
\(V.S'\) and the identity at each auxiliary place; these are the
remaining displayed conditions.  The source's away-coordinate
map \(\mathrm{cmp}\) equals \(\mathrm{cLoc}\) away from \(\ell\)
and equals one at a place above \(\ell\), so all these subgroup
membership assertions are valid on the literal away projection.
The type places lie away from \(\ell\), by \(s.St_\ell\).
Consequently \(\mathrm{toKT}(u)=k\) exactly.

The \(UQ\)-fixed clause of \(H.\mathrm{Big}\) is therefore the
coordinate equation
\[
       \sum_b t(k)_{ib}F_b(y(u,1))=F_i(y)
                     \quad\hbox{for every }y.
\]
Since \(x\) is finite, \(R(x)=1\); since its selected
\(\ell\)-coordinates are one, \(p(gx)=p(g)(u,1)\).
Apply \(E\), then \(R(g^{-1})\), and finally \(\alpha\) to the
last equation at \(y=p(g)\).  Linearity gives
\eqref{eq:ltp-level-matrix}.  In particular the coefficient
matrix has not been discarded or replaced by a scalar character.
\end{proof}
% END LTP10

% BEGIN LTP11
\paragraph{LTP11. Full original \texttt{part\allowbreak UR2Sph}.}
\label{claim:LTP11}
Assume the original \(H.\mathrm{Ax}\).
Let \(w|_{L^+}\notin V.S'\) be split, let \(\tau\) be any
supplied local component of \(\sigma\) at \(w\), and let
\(j\in\mathbf N\), \(a\in\mathbf C\).  If
\[
       \mathrm{sphOp}_\tau(j)v=av
                      \quad\hbox{for every }v\in\tau^{KO_w},
\]
then, for every \(F\in Y\),
\[
                \mathrm{sphOp}_{H.\mathrm{atW}(w)}(j)F=aF .
\]
\begin{proof}
Here \(w|_{L^+}\notin s.S\), so LTP8 supplies the exact
intertwining formula on all coefficient vectors.  The axiom
\(\mathrm{W\_\allowbreak le}\) says that the image of \(KO_w\) under the
original section lies in \(H.UQ\).  Since \(F\in H.\mathrm{Big}\),
it follows that \(H.\mathrm{atW}(w,k)F=F\) for every \(k\in KO_w\).

LTP3 makes \(\sigma\), restricted along \(j_w\), \(\tau\)-isotypic.
Each scalarization of \(F\) lies in that subspace and is \(KO_w\)-fixed.
LTP5 transfers the scalar equation to it.  The same claim's
finite-sum naturality, applied to the linear scalarization maps
and LTP8, gives
\[
 \mathcal T_{(\mathrm{sphOp}_{H.\mathrm{atW}(w)}(j)F)_i}(\alpha)
                          =a\,\mathcal T_{F_i}(\alpha).
\]
LTP1 separates the two coefficient vectors, proving the result.
We used the chosen source double-coset sum, not a rescaled
Hecke operator.  No nonzero-\(Y\) assumption, upper bound on \(j\),
or extra condition on \(\tau_0\) has entered.
\end{proof}
% END LTP11

% BEGIN LTP12
\paragraph{LTP12. Full original \texttt{part\allowbreak UR2Unr}, also at nonsplit places.}
\label{claim:LTP12}
If \(Y\ne0\), then for every finite place \(w\) of \(L\) with
\(w|_{L^+}\notin V.S'\), the original
\(\mathrm{IsUnramU}(\sigma,w)\) holds.
\begin{proof}
Choose \(0\ne F\in Y\).  By LTP1 some
\(\phi=\mathcal T_{F_i}(\alpha)\) is nonzero, and \(\phi\in\sigma\).
We show that this one function is fixed by every
\(x\in\mathrm{locU}(n,w)\) whose \(w\)-coordinate belongs to
the original \(KGL_w=KO_w\).

Such \(x\) is finite and is supported above the base place
\(w|_{L^+}\).  Since that place lies outside \(V.S'\), it lies
outside every selected \(\ell\)-place, every type place, \(v_0\),
and every auxiliary place.  All of \(x\)'s coordinates there
are one.  At a different place \(w'\) above \(w|_{L^+}\), the
last assertion of LTP2 makes its coordinate integral; at all
other base places its coordinates are one.  Thus every hypothesis
of LTP10 holds with \(k=1\).
As \(t(1)\) is the identity matrix,
\eqref{eq:ltp-level-matrix} gives \(\mathscr R_x\phi=\phi\).

The chosen \(\phi\) is nonzero, belongs to \(\sigma\), and has
exactly the invariance required in \(\mathrm{IsUnramU}\).
The argument used the local-support subgroup directly and did
not use a split section at \(w\); it therefore proves the
literal endpoint at nonsplit places as well.
\end{proof}
% END LTP12

\subsection{The prescribed dual infinity weight}

% BEGIN LTP13
\paragraph{LTP13. The actual CM embeddings and the single-factor dual map.}
\label{claim:LTP13}
The map
\[
       a\longmapsto\operatorname{InfinitePlace}
                       (\iota\circ a.\mathrm{res}),\qquad a\in I,
\]
is injective.  For a finite family \(D_a\) of finite-dimensional
complex spaces, any nonzero linear map
\(\Phi:(\bigotimes_aD_a)^*\to N\), and every \(a\), there is a
linear map
\[
                 e_a:D_a^*\longrightarrow(\bigotimes_bD_b)^*
\]
such that \(\Phi e_a\ne0\), which intertwines precomposition by
every tensor operator acting as the identity outside factor \(a\).
For the actual polynomial module,
\[
 \nabla_{\mu^\vee}(\mathbf C)
       \simeq\nabla_\mu(\mathbf C)^*,\qquad
              \mu^\vee_i=-\mu_{n-1-i},
\]
as representations with contragredient action on the right.
The displayed \(\mu^\vee\) is exactly the source's
\(\mathrm{DomWeight.dual}\).
\begin{proof}
Write \(a=(v,j)\), where \(j:L_{s.w(v)}\to
\overline{\mathbf Q}_\ell\) is one of the original continuous
local embeddings.  Such an embedding recovers the chosen prime:
for \(z\in\mathcal O_L\),
\[
                z\in s.w(v)
                   \quad\Longleftrightarrow\quad |j(z)|_\ell<1.
\]
This is the finite local norm calculation of WRP1.  It also follows
by applying continuity to powers of \(z\), and using maximality
of the two resulting prime ideals.

Equal infinite places give equality or complex conjugacy of the
two complex embeddings.  In the equality case, injectivity of
\(\iota\) identifies their restrictions to \(L\), hence their
primes and their base places by the displayed test.  The local
embeddings then agree by continuity and density of \(L\) in its
completion.

In the conjugacy case the CM identity
\(\overline{\iota(j(z))}=\iota(j(cz))\), for \(z\in L\),
identifies one restriction with the other composed with CM
conjugation.  It first gives equality on \(L^+\), hence the same
base place \(v\).  Both chosen places must consequently be the
same \(s.w(v)\), while the prime test makes one the conjugate of
the other.  This contradicts splitting at \(v\).  Thus conjugacy
cannot occur between two indices.  No continuity of \(\iota\)
has been used; the CM identity is an identity of embeddings of
the number field.

For the tensor assertion, pure tensors of dual vectors span the
dual tensor product, by finite-dimensional duality.  Choose
\((\beta_b)_b\) with
\(\Phi(\bigotimes_b\beta_b)\ne0\), and define \(e_a(\beta)\)
by replacing only \(\beta_a\) with \(\beta\).
Then \(\Phi e_a(\beta_a)\ne0\).  Evaluation on a pure tensor
shows directly that
\[
 e_a(\beta)\circ(\bigotimes_b T_b)
          =e_a(\beta\circ T_a)
                       \quad\hbox{if }T_b=1\ (b\ne a).
\]

Finally HP1--HP4 identify the exact \(\nabla_\mu\) as the
irreducible rational representation with highest weight \(\mu\).
Its dual is irreducible: the annihilator of a proper nonzero
dual-stable subspace would be a proper nonzero original-stable
subspace.  Its highest weight is the negative of the lowest
weight of the original representation, namely
\((-\mu_{n-1},\ldots,-\mu_0)\).
Indeed the longest permutation carries the highest weight line
to the lowest weight line, and the functional on that line
annihilating the other weight spaces is fixed by the upper
unipotents in the contragredient action.  The weight calculation
and classification of HP4 therefore give the displayed
isomorphism on the original polynomial modules.  This calculation
specifies a single dual at the same complex embedding; it does
not silently change to a conjugate CM type.
\end{proof}
% END LTP13

% BEGIN LTP14
\paragraph{LTP14. Full original \texttt{part\allowbreak UR2Wt}.}
\label{claim:LTP14}
If \(Y\ne0\), then for every \(a\in I\),
\[
   \mathrm{HasWeightAt}
       \bigl(\sigma,\iota\circ a.\mathrm{res},
                         \lambda(a.\mathrm{res})^\vee\bigr).
\]
\begin{proof}
Choose \(F\in Y\) with one coordinate \(F_i\ne0\).  By ADP7
the linear map
\[
                 \Phi:D^*\longrightarrow\sigma,\qquad
                        \alpha\longmapsto\mathcal T_{F_i}(\alpha)
\]
is nonzero.  If \(x\in G_\infty\), then \(p(x)=1\), so the
exact translation rule reads
\[
       \mathscr R_x\Phi(\alpha)
                         =\Phi(\alpha\circ R(x^{-1})).
\]
Fix \(a\).  On the source subgroup
\(\mathrm{embU}(n,\iota\circ a.\mathrm{res})\), every other
tensor factor of \(R\) is the identity, by LTP13's injectivity
of infinite places.  Use its map \(e_a\) and the dual
isomorphism there.  The composite
\[
 \nabla_{\lambda(a.\mathrm{res})^\vee}(\mathbf C)
       \longrightarrow D_a^*
       \xrightarrow{e_a}D^*
       \xrightarrow{\Phi}\sigma
\]
is nonzero and intertwines this exact embedded unitary action,
with the source's \(\mathrm{cEmb}\) projection.
Its domain is the irreducible polynomial representation
required in \(\mathrm{IsCompAlong}\), by HP2.
These are precisely the two clauses of
\(\mathrm{HasWeightAt}\).  The inverse in the translation rule
is responsible for the dual weight; no sign or determinant shift
is omitted.  The assertion is vacuous when \(I\) is empty.
\end{proof}
% END LTP14

% BEGIN LTP15
\paragraph{LTP15. Full original \texttt{part\allowbreak UR2V0} at depth one.}
\label{claim:LTP15}
If \(Y\ne0\), then every supplied smooth irreducible local
component \(\tau\) of \(\sigma\) at the distinguished place
\(s.w(v_0)\) satisfies
\[
                     \tau^{K_c(L_{s.w(v_0)},n,1)}\ne0 .
\]
\begin{proof}
The original \(v_0\) has norm equal to a rational prime \(p\)
different from \(\ell\).  In particular \(p\in v_0\).
If also \(\ell\in v_0\), a Bezout relation for the distinct
rational primes \(p,\ell\), viewed in \(\mathcal O_{L^+}\),
would put \(1\) in \(v_0\).  Thus \(v_0\) is away from \(\ell\).
The original setting also says \(v_0\notin s.St\), and it is
distinct from every auxiliary base place because those places
lie outside \(s.S\).  The setting's splitting condition supplies
the actual section \(j_0=j_{s.w(v_0)}\).

Choose a nonzero scalarization \(\phi\in\sigma\) of a nonzero
\(F\in Y\), as in LTP12.  For
\(g\in K_c(L_{s.w(v_0)},n,1)\), the finite element \(j_0(g)\)
has identity selected \(\ell\)-coordinates, identity type
coordinates, identity auxiliary coordinates, and identity
coordinates outside \(V.S'\).  Its distinguished coordinate is
exactly \(g\), so LTP10 applies with \(k=1\) and gives
\(\mathscr R_{j_0(g)}\phi=\phi\).  Consequently the restriction
of \(\sigma\) has a nonzero vector fixed by this exact
depth-one subgroup.

LTP3 makes that restriction \(\tau\)-isotypic.  The subgroup is
compact by LTP5, so LTP5 transfers nonzero fixedness to \(\tau\).
No shrinking to an unspecified deeper subgroup is made.
\end{proof}
% END LTP15

\subsection{The type equation with every local correspondence retained}

% BEGIN LTP16
\paragraph{LTP16. One actual product section at the type places.}
\label{claim:LTP16}
There is a homomorphism
\[
       J:\prod_{u\in s.St}\operatorname{GL}_n(L_{s.w(u)})
                                      \longrightarrow G_f
\]
such that \(J(\delta_u(g))=j_{s.w(u)}(g)\), its \(s.w(u)\)-coordinate
is \(g_u\), and all coordinates at other base places are one.
In particular it has identity selected \(\ell\)-coordinates,
identity distinguished coordinate, and identity auxiliary coordinates.
For every \(F\in H.\mathrm{Big}\), \(k\in K_T\), and \(i,\alpha\),
\[
       \sum_b t(k)_{ib}\,
               \mathscr R_{J(k)}\mathcal T_{F_b}(\alpha)
                          =\mathcal T_{F_i}(\alpha).
\]
\begin{proof}
Every original type place splits, and the chosen place lies above
its indicated base place, by FNL4.  The base places indexed by
\(s.St\) are distinct.  Their sections therefore commute by LTP2.
Define \(J(g)\) as the finite product of their images.
Commutation makes the product independent of order and makes
\(J\) a homomorphism.  Its coordinates have the stated values
by construction.  This remains valid for an empty type set.

The type places lie away from \(\ell\), exclude \(v_0\), and are
disjoint from the auxiliary base places.  They are all contained
in \(V.S'\).  Consequently \(J(k)\) has precisely the coordinates
required by LTP10: its type coordinates are the prescribed \(k\),
its distinguished coordinate is the identity in \(K_c(1)\), and
outside \(V.S'\) it is the identity.  Apply LTP10 to obtain the
displayed matrix equation.
\end{proof}
% END LTP16

% BEGIN LTP17
\paragraph{LTP17. Full original \texttt{part\allowbreak UR2Type\_\allowbreak alt}.}
\label{claim:LTP17}
Assume \(Y\ne0\).  Let
\(\chi_{u,i}:I_{s.A(u)}\to\mathbf C^\times\) be any family
for which the original
\[
      \mathrm{IsTypeAlt}\bigl((s.hw(u))_u,t,\chi\bigr)
\]
holds.  Let \((\tau_u)_{u\in s.St}\) be any family of smooth
irreducible local components of \(\sigma\) at the chosen type
places.  Then, for every \(u\in s.St\) and every map
\(\mathrm{rc}\) satisfying the original
\(\mathrm{IsRecT}(s.hw(u),\mathrm{rc})\),
\[
        \mathrm{IsTypeWD}\!
          \left(s.hw(u),\chi_u,
                         \mathrm{rc}(\mathrm{algNorm}(\tau_u))\right).
\]
\begin{proof}
Choose \(0\ne F\in Y\).  LTP1 gives \(i_0,\alpha\) such that
\(\mathcal T_{F_{i_0}}(\alpha)\ne0\).  Keep this \(\alpha\)
fixed, and put
\(\Phi_i=\mathcal T_{F_i}(\alpha)\in\sigma\).
Thus \(\Phi=(\Phi_i)_i\) is a single nonzero coefficient vector
in the original nonzero vector space \(\sigma\).

Use LTP16's one homomorphism \(J\) to define a representation
of \(\prod_u\operatorname{GL}_n(L_{s.w(u)})\) on \(\sigma\):
\(\rho(g)\phi=\mathscr R_{J(g)}\phi\).
Global stability of \(\sigma\) makes this well defined.
Restriction to its \(u\)-th factor is exactly the action along
\(j_{s.w(u)}\), so LTP3 says that this restricted representation
is \(\tau_u\)-isotypic.  Inclusion of \(\sigma\) into the ambient
functions identifies this assertion with the literal
\(\mathrm{IsoSum}(\rho|_u,\top,\tau_u)\).

LTP16 proves, for every \(k\in K_T\),
\[
                   \sum_b t(k)_{ib}\rho(k)\Phi_b=\Phi_i .
\]
We have therefore verified the entire left side of the equivalence
in the supplied \(\mathrm{IsTypeAlt}\), on this very nonzero
module \(\sigma\) and this product action.  Its right side is
exactly the displayed conclusion, simultaneously for all \(u\)
and all maps \(\mathrm{rc}\) satisfying \(\mathrm{IsRecT}\).
We do not replace that right side by a statement for one chosen
correspondence or by a statement only about unramified parameters.
No finite-length or admissibility hypothesis on the product
module is needed by this application.
\end{proof}
% END LTP17

% BEGIN LTP18
\paragraph{LTP18. Every original Teichmuller family uses the same integral type.}
\label{claim:LTP18}
Let \(r\) be the original residual representation and suppose
\(\tau_0\) satisfies its original \(\mathrm{TypeForAlt}(r,\iota)\),
as supplied by ITA12.  For every original family \(t_0\) with
\(r.\mathrm{IsTeich}(t_0)\), define
\[
 \chi_{u,i}(\gamma)
       =\iota\!\left(\operatorname{algebraMap}_{O,\overline{\mathbf Q}_\ell}
                      \bigl(t_{0,i}(V.\mathrm{pr}(\gamma))\bigr)\right),
                         \qquad\gamma\in I_{s.A(u)} .
\]
For the same \(Y\ne0\) and any family of its local components,
LTP17's \(\mathrm{IsTypeWD}\) conclusion holds for these
\(\chi_{u,i}\), at every type place and for every original
\(\mathrm{IsRecT}\) map.
\begin{proof}
This \(\chi\) is the literal pushed-character family in the
definition of \(\mathrm{TypeForAlt}\): the inertia inclusion
into the original absolute Galois group is followed by
\(V.\mathrm{pr}\), then by \(t_{0,i}\), and finally by the
fixed coefficient inclusion and \(\iota\).
The asserted type condition for every such \(t_0\) is a field
of the already supplied \(\mathrm{TypeForAlt}\).
Apply LTP17.  The integral module, \(\tau_0\), \(O\), \(\iota\),
the vector space \(Y\), and the product action \(J\) are unchanged
as \(t_0\) and the local correspondence vary.
In particular the argument retains coefficient prime three,
repeated inertia characters, an empty type set, and every
original allowed rank.  It asserts the Weil--Deligne type
conclusion; descent of a Galois representation and its inertia
conjugating matrices are separate subsequent steps.
\end{proof}
% END LTP18

\subsection{Producing the local components required by the assembly}

% BEGIN LTP19
\paragraph{LTP19. An admissible local component inside the actual automorphic representation.}
\label{claim:LTP19}
For every original automorphic representation \(\sigma\) and
every split finite place \(w\), there is a smooth irreducible
admissible representation \(\tau_w\) of
\(\operatorname{GL}_n(L_w)\) and a nonzero local intertwiner
\(\tau_w\to\sigma\).  This proves the existence and admissibility
clauses of \(\mathrm{DefU.exists\_\allowbreak locComp\_\allowbreak split}\).
The argument uses only the identity-form group and therefore
applies to every CM field and every \(n\ge0\), as in that endpoint.
\begin{proof}
First let \(n\ge1\), and put \(J_w=\operatorname{GL}_n(L_w)\).
LTP2's coordinate construction is a topological product
decomposition
\[
                    G_f=j_w(J_w)\times G_f^{(w)} ,
\]
where the second factor has identity coordinates above
\(w|_{L^+}\).  Both factors commute with \(G_\infty\).

Choose a nonzero vector of \(\sigma\).  Its finite-dimensional
\(G_\infty\)-orbit span contains a nonzero irreducible
\(G_\infty\)-subrepresentation \(S\): the positive invariant form
of ADP1 allows finite-dimensional invariant orthogonal
complements, or one may choose a minimal nonzero invariant
subspace.  Thus the inclusion \(S\hookrightarrow\sigma\) is
an element of
\[
                \mathcal E=\operatorname{Hom}_{G_\infty}(S,\sigma).
\]
Choose a basis of \(S\).  The images of its finitely many basis
vectors have open finite stabilizers by the original automorphic
smoothness condition.  Their intersection contains
\(C_w\times K^{(w)}\), with \(C_w\) compact open in \(J_w\)
and \(K^{(w)}\) compact open in \(G_f^{(w)}\).
Here compact open subgroups are obtained by integral matrices
and principal congruence conditions in the actual finite
coordinates.  Hence
\[
                       M=\mathcal E^{K^{(w)}}\ne0
\]
and \(M\) is stable under \(J_w\).

Every vector of \(M\) is smooth under \(J_w\): intersect the
open finite stabilizers of its values on a basis of \(S\), and
pull back along the continuous \(j_w\).
For every compact open \(C\subset J_w\),
\[
    M^C\ \subset\
       \operatorname{Hom}_{G_\infty}(S,\mathcal A)^{
                                      j_w(C)K^{(w)}} .
\]
The group \(j_w(C)K^{(w)}\) is compact open in the actual \(G_f\).
ADP2 therefore makes the right side finite-dimensional.
This proves the exact local admissibility bound for \(M\);
it does not infer a finite multiplicity statement merely from
compactness of the quotient.

If \(s_1,\ldots,s_e\) is a basis of \(S\), give \(M\) the form
\[
            (T,T')=\sum_{b=1}^e
                         \langle T(s_b),T'(s_b)\rangle_{\mathcal A}.
\]
ADP1 makes it positive definite and \(J_w\)-invariant.
Compact subgroups of \(J_w\) have finite orbits on every smooth
vector of \(M\).  For a compact open \(C\), \(M^C\) is finite
dimensional as just proved.  The algebraic projection argument
PKT5--PKT6 consequently expresses every vector of \(M\) as a
finite sum of vectors in irreducible \(J_w\)-submodules.
All of its hypotheses have been verified here on \(M\).
For clarity, that argument projects inside the finite space
\(M^C\), uses finite compact orbits to make the residual
orthogonal to an entire invariant submodule, and inducts on the
dimension of the fixed space; no Hilbert completion or closed
subrepresentation assumption is used.

Choose a nonzero irreducible submodule \(U\subset M\).
It is smooth, and \(U^C\subset M^C\) is finite-dimensional for
every compact open \(C\), so it is admissible.
There are \(T\in U\) and \(s_0\in S\) with \(T(s_0)\ne0\).
Evaluation
\[
                       U\longrightarrow\sigma,\qquad
                                T\longmapsto T(s_0)
\]
is \(J_w\)-equivariant, since \(J_w\) acts only on the target
of the Hom space.  It is nonzero, so its kernel is zero by
irreducibility of \(U\).  Regard \(U\) with this action as the
smooth irreducible \(\tau_w\).  Since every
\(x\in\mathrm{locU}(n,w)\) equals \(j_w(\mathrm{cLoc}_w x)\),
the evaluation map is precisely a nonzero intertwiner in the
source's \(\mathrm{IsCompAlong}\).

The group, finite-class and multiplicity arguments used here
depend only on a CM field and its identity-form unitary group;
the auxiliary Hida data are absent from their statements and
proofs.  This supplies the asserted general CM-field scope.
For \(n=0\), both adelic and local matrix groups are trivial,
the nonzero irreducible automorphic space is the one-dimensional
space of functions on a point, and its one-dimensional local
representation is smooth and admissible.  This treats the
remaining rank in the literal \texttt{exists\_\allowbreak loc\allowbreak Comp\_\allowbreak split}
statement.
\end{proof}
% END LTP19

% BEGIN LTP20
\paragraph{LTP20. The exact unramified--spherical equivalence for every local component.}
\label{claim:LTP20}
For every smooth irreducible local component \(\tau\) of
\(\sigma\) at a split place \(w\),
\[
            \mathrm{IsUnramU}(\sigma,w)
                 \quad\Longleftrightarrow\quad
                   \tau^{KGL_w}\ne0 .
\]
Together with LTP19 this proves the full original
\(\mathrm{DefU.exists\_\allowbreak locComp\_\allowbreak split}\), including admissibility.
\begin{proof}
If \(\mathrm{IsUnramU}\) holds, its nonzero function is fixed
by \(j_w(KGL_w)\).  LTP3 makes \(\sigma\) restricted along
\(j_w\) \(\tau\)-isotypic.  The equality \(KGL_w=KO_w\) and
compactness from LTP5 therefore imply \(\tau^{KGL_w}\ne0\).

Conversely, let \(0\ne v\in\tau^{KGL_w}\).
The supplied component map \(f:\tau\to\sigma\) is nonzero.
Its kernel is invariant under the local group, so irreducibility
makes \(f\) injective.  Thus \(\phi=f(v)\ne0\).
For every \(x\in\mathrm{locU}(n,w)\) whose \(w\)-coordinate
belongs to \(KGL_w\), equivariance gives
\[
       \mathscr R_x\phi
                   =f\bigl(\tau(\mathrm{cLoc}_w x)v\bigr)
                   =f(v).
\]
This is exactly the witness demanded by \(\mathrm{IsUnramU}\).
The equivalence holds for every supplied component, not only
for the admissible component chosen in LTP19.
\end{proof}
% END LTP20

% BEGIN LTP21
\paragraph{LTP21. One family of components supplies all the original local fields.}
\label{claim:LTP21}
Fix \(Y\ne0\).  One may choose a single family
\((\tau_w)_{w\text{ split}}\) before using any of the following
local conclusions.  Each \(\tau_w\) is smooth irreducible and
admissible and is a component of this same \(\sigma\).  For this
one family the source's local requirements hold simultaneously:
\[
\begin{array}{c|l}
 \text{place or condition}&\text{conclusion on the original objects}\\ \hline
 w=s.w(v),\ v\mid\ell&\mathrm{IsoSum}(H.\mathrm{atV}(v),Y,\tau_w)\\
 w=Q.wq(q)&\mathrm{IsoSum}(H.\mathrm{atW}(w),Y,\tau_w)\\
 w|_{L^+}\notin V.S',\ w\text{ split}
       &\tau_w^{KGL_w}\ne0,\ \text{and the LTP11 spherical transfer}\\
 a\in I&\mathrm{HasWeightAt}(\sigma,\iota\circ a.\mathrm{res},
                                               \lambda(a.\mathrm{res})^\vee)\\
 w=s.w(v_0)&\tau_w^{K_c(L_w,n,1)}\ne0\\
 u\in s.St&\text{LTP17 for every original }\mathrm{IsRecT}
                 \text{ and every supplied type family.}
\end{array}
\]
In addition \(\mathrm{IsUnramU}(\sigma,w)\) holds at every
finite \(w\) outside \(V.S'\), whether or not \(w\) splits.
\begin{proof}
Apply choice once to LTP19, indexed by the set of split places.
Retain the resulting local representations and component maps
throughout.  FNL4 identifies the selected \(\ell\)-places,
auxiliary places, type places and \(v_0\) as members of this
index set.  LTP7, LTP9, LTP12, LTP14 and LTP15 give the
corresponding statements, and LTP20 changes LTP12 into spherical
fixedness at a split outside place.  LTP11 transfers every
given spherical scalar, on the original operator and every
coefficient vector in \(Y\).  Restrict the same family to the
type places in LTP17; LTP18 makes the conclusion universal over
the original Teichmuller families when the supplied integral
type is used.  FNL16 supplies the original \(H.\mathrm{Ax}\)
needed by LTP9 and LTP11.

This is one collection of local components of one original
automorphic representation, at the same \(E,\xi,\iota,\tau_0\)
and \(Y\).  No independent replacement of components in the
different rows is necessary.  It supplies the local-transport
inputs to the subsequent \texttt{Part\allowbreak U} construction; that
construction still has to provide its global Galois representation,
local--global and Frobenius comparisons, quotient descent,
pairing, and point-property record.
\end{proof}
% END LTP21


\section{The literal base-change-to-Speh bridge}
\label{sec:base-change-v44}

This chapter proves the complete finite consumer
\[
 \mathrm{BaseChangeIsobaricN}\ \Longrightarrow\
 \mathrm{DefiniteN}\ \Longrightarrow\ \mathrm{BcSpeh}.
 \tag{BCS.1}
\]
The first premise is the original named hypothesis
\texttt{H.\allowbreak lit\allowbreak N.\allowbreak aut.\allowbreak bc}.  It is retained with its full definite-inner-form,
all-rank and nongeneric scope.  The chapter does not supply a new
trace-formula proof of that premise.  The other named inputs retained
from the original literature record are
\texttt{H.\allowbreak lit\allowbreak N.\allowbreak loc.\allowbreak bz}, \texttt{bern}, \texttt{sq}, \texttt{cos},
\texttt{H.\allowbreak lit\allowbreak N.\allowbreak aut.\allowbreak fin}, and the boundedness clause of
\texttt{H.\allowbreak lit\allowbreak N.\allowbreak aut.\allowbreak cusp}.  A precisely stated published local
generic-unitary decomposition is matched in BCS8.  Basic Haar measure,
the finite-volume general-linear adelic quotient, Fourier uniqueness on
compact additive adelic quotients, and highest-weight theory for
\(U(n)\) are the original analytic and representation-theoretic
foundations.  Their uses are identified below.  None of these inputs
is silently counted as a newly reconstructed theorem.

Let \(E/E^+\) be the original CM number field.  It is unramified at
all finite places, as required by the base-change hypothesis.
Let \(n\ge1\), and
let \(\pi\) be an actual irreducible subspace of the original
automorphic functions on the definite unitary adelic group
\(U_{\mathbb A}\), with \(\mathrm{IsAutRepU}(\pi)\).
At a finite place \(w\), put \(F=E_w\),
\(K_r=\mathrm{GL}_r(\mathcal O_F)\), and
\(q_w=\#(\mathcal O_E/w)\).  All products \(\times\) below are
normalized induction from the upper standard block parabolic, with
the original increasing block order.  Thus
\[
 f(pg)=\delta_P(p)^{1/2}\rho(m_p)f(g),\qquad
 \nu^t(g)=|\det g|_F^t,\qquad
 s_{d,k}=\frac{d-1}{2}-k\quad(0\le k<d).
 \tag{BCS.2}
\]
For a multiset \(m\), write
\[
 \operatorname{spehInf}(d,m)
   =\sum_{k=0}^{d-1}\{z+s_{d,k}:z\in m\}.
 \tag{BCS.3}
\]
The definition of \(\mathrm{IsSpeh}(\rho,d,\sigma)\) is stronger
than being an unspecified irreducible quotient: \(\sigma\) is the
unique irreducible cosocle of
\(\rho\nu^{s_{d,0}}\times\cdots\times\rho\nu^{s_{d,d-1}}\),
and every maximal proper invariant subspace of that induction is the
kernel of its quotient map.  We prove this precise assertion.

\subsection{Replacing the selected subquotient by the full local product}

% BEGIN BCS1
\paragraph{BCS1. The actual discrete and normalized-cuspidal local forms.}
\label{claim:BCS1}
Every local component of an original discrete automorphic
\(\mathrm{GL}_r\)-space is unitarisable.  A normalized cuspidal
space is itself discrete, so the same conclusion holds for each of
its local components.
\begin{proof}
The original \(\mathrm{IsDiscrete}(V)\) includes square
integrability of every actual function in \(V\), invariance under
\(\mathrm{GL}_r(E)\) and the real split center, stability under all
finite, compact-infinity and Lie operators, and nonzero irreducibility.
Choose one Haar measure and one fundamental domain for the source
subgroup
\(\mathrm{GL}_r(E)\langle\operatorname{splitCentre}(1)\rangle\).
The definition permits separate choices for separate functions; Haar
uniqueness and invariance of the integral under changing a fundamental
domain put every function in \(L^2\) for this common choice.  Polarization
and Cauchy--Schwarz therefore define a single form
\[
 B(f,h)=\int_{\mathcal F}\overline{f(g)}h(g)\,dg.
 \tag{BCS.4}
\]
It is positive definite on the actual continuous functions.  Indeed a
nonzero value gives an open set of positive quotient Haar measure on
which the absolute value is bounded below.  Such a function cannot have
zero \(L^2\) norm.  The adelic general linear group is unimodular;
right translation preserves the quotient measure and hence (BCS.4).

An original local component is a nonzero local intertwining map
\(T:\sigma\to V\).  Its kernel is invariant and \(\sigma\) is
irreducible, so \(T\) is injective.  The pullback \(B(Tv,Tz)\) is
the required positive invariant Hermitian form on \(\sigma\).
No completion of its underlying algebraic vector space is necessary.

For a normalized cusp \(P\), the normalization says that every
function is fixed by the same real split center.  The retained
\texttt{Cusp\allowbreak Forms\allowbreak Bounded\allowbreak GL} clause bounds each such cusp function.
The original quotient just used has finite volume, so boundedness
implies its square integrability.  The other clauses of
\(\mathrm{IsDiscrete}(P.V)\) are already clauses of the cuspidal
subspace.  Apply the preceding pullback argument to \(P.V\).
This last argument uses only normalization, not an unproved assertion
that a cuspidal local component is tempered.
\end{proof}
% END BCS1

% BEGIN BCS2
\paragraph{BCS2. All finite unitary local products, with the original order.}
\label{claim:BCS2}
Let \(\sigma_i\) be smooth irreducible unitarisable representations
of \(\mathrm{GL}_{N_i}(F)\), with \(N_i>0\), and let
\(\sum_iN_i=n\).  Their original normalized product
\(\operatorname{locProd}(h,w,\sigma)\) is smooth, unitarisable and
irreducible.
\begin{proof}
Tensor the positive forms on the inducing spaces.  By Iwasawa
decomposition, normalized induction is its space of locally constant
sections over the compact quotient \((P\cap K_n)\backslash K_n\).
The formula
\[
             (f_1,f_2)=\int_{K_n}B(f_1(k),f_2(k))\,dk,
             \qquad\operatorname{vol}(K_n)=1,
 \tag{BCS.5}
\]
is positive definite.  It is invariant under right translation by
\(G\): the change of Iwasawa coordinates has Jacobian
\(\delta_P^{-1}\), while the two covariance factors contribute
\(\delta_P\).  Smoothness follows from the right open stabilizer
in the induced model.  Thus normalized induction preserves
unitarisability in precisely this normalization.

For two factors, irreducibility is the retained literal
\texttt{Bernstein\allowbreak Unitary\allowbreak Product\allowbreak Irreducible} hypothesis.  Induct on
the number of factors.  The product of the first \(c-1\) factors
is smooth irreducible and unitarisable by the induction hypothesis,
so the same two-factor theorem applies to it and the last factor.
Induction in stages identifies this result with the original full
product.  Explicitly, a section is sent to the function of two group
variables obtained by successive evaluation; the inverse evaluates
the inner section.  The modulus identity
\(\delta_P^{1/2}=\delta_Q^{1/2}\delta_{P\cap M_Q}^{1/2}\)
checks covariance.  Both maps use the original block embeddings and
their original order.  One factor gives the representation itself;
the empty product, when needed at rank zero, is the one-dimensional
representation of the trivial group.  No distinctness of the factors
and no bound on their number enters this proof.
\end{proof}
% END BCS2

% BEGIN BCS3
\paragraph{BCS3. The full original \texttt{Definite\allowbreak N} bridge.}
\label{claim:BCS3}
The original \texttt{Base\allowbreak Change\allowbreak Isobaric\allowbreak N} and the inputs of BCS1--BCS2
imply the entire original \texttt{Definite\allowbreak N}.
\begin{proof}
Take exactly the base-change witnesses
\(c,N_i,h,V_i,\sigma_{i,w},M_{i,\tau},Q_w\).
Here \(Q_w\) is the selected irreducible subquotient of the full
product \(I_w=\operatorname{locProd}(h,w,\sigma_{\bullet,w})\).
BCS1 makes every \(\sigma_{i,w}\) unitarisable, and BCS2 makes
\(I_w\) irreducible.  A nonzero irreducible subquotient of an
irreducible representation is isomorphic to it: the invariant
subspace from which the quotient is taken is nonzero, hence the
whole space; the quotient kernel is proper, hence zero.
Choose the resulting equivalence \(e_w:Q_w\simeq I_w\).

Every fixed vector and every finite Hecke sum transports along
\(e_w\).  Thus the supplied spherical clause at every
\(\mathrm{IsUnramU}(\pi,w)\) becomes the spherical clause for
\(I_w\).  For every permitted \(w,s,\varpi,q,\alpha\), the
supplied inert clause becomes
\[
 \mathrm{HasSatGL}\bigl(w,I_w,\operatorname{bcParam}(n,\alpha)\bigr).
 \tag{BCS.6}
\]
This keeps the original conditions that \(s\) extends CM conjugation,
that \(\mathrm{IsUnramInv}(\varpi,s,q)\) holds, and that
\(\mathrm{HasSatAt}(\pi,w,s,\varpi,q,\alpha)\) holds.

At a split place the original nonzero map from \(Q_w\) into the
actual unitary automorphic space is precomposed with \(e_w^{-1}\).
Its image is still in \(\pi\), and its intertwining identity is
still along
\((cLoc_w)\circ(\operatorname{locU}(n,w)\hookrightarrow U_{\mathbb A})\).
Hence it gives the full \(\mathrm{IsCompAlong}\) clause for
\(I_w\), with precisely that subgroup and map.  All other fields
---positive \(N_i\), discreteness, \(\mathrm{CsdAsc}\), actual
local components, infinitesimal characters, and the weight sum---are
copied without change.  These are exactly all fields of
\(\mathrm{IsBaseChangeU}\), proving \texttt{Definite\allowbreak N}.
\end{proof}
% END BCS3

\subsection{The bounded residual-spectrum consumer}

The input \(\mathrm{CsdAsc}(V)\) already supplies one normalized
cuspidal \(P\) of rank \(a\) and a length \(d\), with
\(da=N\), together with an injective map from \(V\) to the
original unnormalized induced automorphic functions.  The map
intertwines all finite, compact-infinity and Lie operators.  Its
\(k\)-th cuspidal coefficient is
\[
                P\,|\det|_{\mathbb A_E}^{(a-1)s_{d,k}}.
 \tag{BCS.7}
\]
The following argument consumes this actual embedding.  It does not
infer the embedding from a classification assertion.

% BEGIN BCS4
\paragraph{BCS4. Actual local components, isotypy and infinity characters.}
\label{claim:BCS4}
Every original irreducible cuspidal or discrete automorphic space
has an actual irreducible smooth local component at every finite
place and an infinitesimal-character multiset at every embedding.
Its restriction to a fixed finite local group is an algebraic sum
of copies of any one of its local components.  In particular all
such components are isomorphic.  The local existence assertion uses
the retained \texttt{Automorphic\allowbreak Forms\allowbreak Finite\allowbreak GL}.
\begin{proof}
First, every central Capelli operator preserves the finite-dimensional
space generated from a nonzero automorphic vector by the center.
It has an eigenvector.  Since it commutes with all three kinds of
automorphic operators, its eigenspace is a nonzero globally stable
subspace, hence the entire irreducible space.  The Capelli expression
\(C_r(u)=\operatorname{qdet}(E_{ij},u+(r-1)/2)\) is monic of
degree \(r\) in \(u\).  Evaluating on a fixed nonzero vector and
a linear functional taking value one gives a monic scalar polynomial.
Factoring it over \(\mathbb C\) gives a multiset \(m\) of size
\(r\) such that
\(C_r(u)=\prod_{z\in m}(u+z)\) on the entire space.  This is the
literal \(\mathrm{HasInfCharGL}\), or \(P.\mathrm{HasInfChar}\).
The commutation identity used here follows by expanding the Capelli
column determinant and using
\([E_{ij},E_{kl}]=\delta_{jk}E_{il}-\delta_{li}E_{kj}\);
the two terms associated with each off-diagonal replacement cancel,
including the consecutive diagonal shifts.  Equal scalar polynomials
have equal multisets, so these characters are unique on a nonzero
space and are preserved by all-operator equivalences.

For completeness, local existence can be obtained without assuming
a restricted tensor product in advance.  Fix a nonzero vector \(v\)
and a compact open local subgroup \(K\) fixing it.  Let \(e_K\)
be normalized compact averaging.  The local subspace
\[
                  Z=\bigcap_{g\in\mathrm{GL}_r(F)}
                       \ker(e_K\,g)
 \tag{BCS.8}
\]
is stable under the local group and every operator commuting with
it.  Finite adelic operators factor as a local operator and one
trivial at this place; infinity operators commute with it.  Thus
\(Z\) is globally stable.  Since \(e_Kv=v\ne0\), global
irreducibility gives \(Z=0\).  Consequently every nonzero local
subrepresentation has nonzero \(K\)-fixed vectors: otherwise
averaging every translate puts it inside \(Z\).

Take the nonzero local-stable space \(T\) of vectors with the
fixed compact-infinity type generated by \(v\), fixed by a chosen
open compact group away from the place.  It has the central
character just proved.  This finite-dimensional infinity type is
continuous: evaluations separate its finite span, so finitely many
evaluations give coordinates, and each coordinate of its orbit is
continuous because the original automorphic functions are smooth.
For each local compact open \(K'\), its
\(K'\)-fixed part is finite dimensional by the retained automorphic
finiteness theorem.  This is an application to an actual open finite
adelic level: intersect the original level with the inverse image of
\(K'\) under the local projection and with the local inclusion
condition.  Thus all original level and infinity-type conditions are
preserved.

Among the nonzero invariant subspaces \(W\le T\), choose one
minimizing the positive integer \(\dim W^K\).  Let \(X\) be the
local span of \(W^K\).  A nonzero invariant subspace \(Y\le X\)
has \(\dim Y^K\ge\dim W^K\); since \(Y^K\subseteq W^K\),
they are equal.  Hence \(Y\) contains \(W^K\) and so equals
\(X\).  This gives an actual irreducible smooth subspace and its
inclusion, not just a composition factor.  It is admissible, either
by the same finite-dimensional fixed spaces or by
\texttt{Bernstein\allowbreak Zelevinsky\allowbreak Finiteness}.

Now fix any irreducible local copy \(X\).  The sum of the images
of all local maps \(X\to V\) is stable under local operators and
under their commuting operators, hence is all of \(V\).  A sum of
simple submodules is semisimple: enlarge a direct sum of them until
it is maximal; any further simple submodule either meets it trivially
or is contained in it.  This proves algebraic isotypy.  Schur's
endomorphism field is \(\mathbb C\): an endomorphism restricts to
a nonzero finite-dimensional compact-fixed space, has an eigenvalue
there, and irreducibility makes its eigenspace the whole local
representation.  Thus a nonzero vector in any tensor of these
isotypic restrictions is detected by a local-equivariant projection
onto the tensor of the chosen components.  Such a tensor is simple,
as successive use of the density theorem for a simple module with
scalar endomorphisms separates the finitely many tensor coordinates.
\end{proof}
% END BCS4

% BEGIN BCS5
\paragraph{BCS5. Cuspidal local components are generic.}
\label{claim:BCS5}
Every actual finite local component of a positive-rank original cusp
admits a nonzero nondegenerate Whittaker functional.
\begin{proof}
Choose the standard continuous additive character of
\(\mathbb A_E/E\), nontrivial at each completion.  We recall why
cuspidality gives the needed nonzero coefficient.  Expand a nonzero
cusp function along the last-column additive radical.  The constant
coefficient is its proper-parabolic constant term and is zero.
Fourier uniqueness on the compact quotient
\((\mathbb A_E/E)^{a-1}\) leaves a nonzero nonconstant coefficient.
The rational group \(\mathrm{GL}_{a-1}(E)\) is transitive on
nonzero row characters, so rational left invariance identifies these
coefficients with translates of the one using the last coordinate.
That coordinate coefficient therefore cannot vanish identically.

Repeat on the upper-left \((a-1)\)-block.  At each stage the
previously imposed last-coordinate characters commute with the next
column integration.  The zero new coefficient is the constant term
along the corresponding larger proper parabolic, followed by the
previous character integrations; it vanishes by cuspidality and
Fubini on compact additive quotients.  The same rational transitivity
and Fourier uniqueness show that the next designated nonconstant
coefficient is nonzero as a function.  Induction ends with a
nonzero function
\[
 W_f(g)=\int_{N(E)\backslash N(\mathbb A_E)}
     f(ug)\,\psi\!\left(-\sum_{j=1}^{a-1}u_{j,j+1}\right)du.
 \tag{BCS.9}
\]
Successive integration in the column coordinates is a definition of
this integral; compactness makes all integrations finite.  This
argument also proves the transformation law under \(N(\mathbb A_E)\)
by translating the integration coordinates.  For rank one, the
unipotent group is trivial and \(W_f=f\).

Let \(T:\rho\to P.V\) be the given nonzero local map.  It is
injective, so choose \(v\) with \(Tv\ne0\).  Some
\(W_{Tv}(g_0)\) is nonzero.  Factor \(g_0=h g_w\), where \(h\)
has identity component at \(w\).  Then
\(\ell(x)=W_{Tx}(h)\) is nonzero on \(\rho(g_w)v\).
The element \(h\) commutes with the local group.  The transformation
law from (BCS.9) therefore gives
\[
 \ell(\rho(u)x)=\psi_w\!\left(\sum_j u_{j,j+1}\right)\ell(x),
                         \qquad u\in N(E_w).
\]
The local character is nontrivial.  This is exactly genericity of
the given local representation.
\end{proof}
% END BCS5

% BEGIN BCS6
\paragraph{BCS6. Smooth Hermitian duality turns the ascending map into a quotient.}
\label{claim:BCS6}
Let \(\rho,\sigma\) be smooth irreducible unitarisable
representations of \(\mathrm{GL}_a(F),\mathrm{GL}_{ad}(F)\), with
\(d\ge1\).  A nonzero map
\[
 \sigma\longrightarrow
       \rho\nu^{-s_{d,0}}\times\cdots\times\rho\nu^{-s_{d,d-1}}
 \tag{BCS.10}
\]
implies a surjective map from the descending product
\(\rho\nu^{s_{d,0}}\times\cdots\times\rho\nu^{s_{d,d-1}}\)
onto \(\sigma\).
\begin{proof}
For a compact open \(K\), normalized averaging is a finite average
on each smooth vector: choose an open normal subgroup of \(K\)
fixing that vector and average over the finite quotient.  It is an
idempotent onto fixed vectors.  In particular, the smooth dual of an
injection is surjective.  To extend a smooth functional on the smaller
space, first extend it algebraically, then average the extension over
a compact open subgroup fixing the original functional.  The
restriction remains the original functional.  The same argument
applies to the conjugate, or Hermitian, smooth dual \(V^h\).

For an admissible unitary representation with positive form \(B\),
the map \(v\mapsto B(v,\cdot)\), interpreted with the conjugate
scalar structure, identifies it with its Hermitian smooth dual.  To
see surjectivity, a \(K\)-fixed functional is determined by its
restriction to the finite-dimensional space \(V^K\); there the
positive form represents it.  Orthogonal averaging shows that the
same representing vector works on all of \(V\).  The needed
admissibility is the literal irreducible-admissibility clause of
\texttt{H.\allowbreak lit\allowbreak N.\allowbreak loc.\allowbreak bz}.

The normalized-induction duality in that same record, paired by the
integral over \(K_{ad}\), identifies the Hermitian dual of an induced
product with the product of the Hermitian duals in the \emph{same}
upper-parabolic order.  Complex conjugation followed by ordinary
contragredience sends \(\rho\nu^t\) to
\(\rho^h\nu^{-t}\) for real \(t\).
The map (BCS.10) is injective by irreducibility.  Take its Hermitian
smooth dual, use the preceding surjectivity, and then the unitary
identifications \(\rho^h\simeq\rho\) and
\(\sigma^h\simeq\sigma\).  This gives the asserted quotient,
with each exponent \(-s_{d,k}\) changed to \(s_{d,k}\).
There is no reordering theorem or local Gelfand--Kazhdan assumption
hidden in this step.
\end{proof}
% END BCS6

% BEGIN BCS7
\paragraph{BCS7. Extracting the exact ascending local induction.}
\label{claim:BCS7}
\leavevmode\par
Suppose \(V\) is discrete, \(\operatorname{AscEmb}(V,P,d)\),
\(N\ge1\), and \(\rho,\sigma\) are any actual local
components of \(P,V\) at \(w\).  Then \(da=N\), both
\(a,d\) are positive, and there is a nonzero map (BCS.10).
\begin{proof}
The rank identity is the block identity in \(\operatorname{AscEmb}\).
Its positive product gives \(a,d\ge1\).  Compose the actual
component map \(\sigma\to V\) with the injective ascending
embedding to get a nonzero local map \(F\) into the given space
of induced automorphic functions.

Choose \(v_0,g_0\) with \(F(v_0)(g_0)\ne0\), and write the
adelic Iwasawa decomposition
\(g_0=u m_0 k_\infty x_f\).  Split
\(x_f=x^{(w)}\operatorname{loc}_w(x_w)\), with
\(x^{(w)}\) equal to one at \(w\).  The map
\[
 L(v)(m)=F(v)(m k_\infty x^{(w)})
 \tag{BCS.11}
\]
has values in the original tensor space of the twisted cuspidal
coefficients.  It is nonzero, since evaluation at
\(v=\sigma(x_w)v_0,m=m_0\) recovers the chosen nonzero value;
left unipotent invariance removes \(u\).  The element
\(k_\infty x^{(w)}\) commutes with the local group.  Writing a
local parabolic element as its Levi part times its unipotent part,
and moving the latter to the left, proves
\[
 L(\sigma(p)v)(m)=L(v)(m\,m_p),\qquad p\in P(F),
 \tag{BCS.12}
\]
where \(m_p\) is the original local Levi inclusion.

BCS4 makes the restriction of each twisted cuspidal space an
algebraic sum of copies of the local twist
\[
                         \rho\nu^{(a-1)s_{d,k}}.
\]
The multiplication map from the algebraic tensor of coefficient
spaces to their functions on the product of Levi groups is injective:
for each finite sum of pure tensors, independent evaluations on each
factor separate the finite-dimensional spans of its coordinates.
Thus the original \(\mathrm{tensorForms}\) is this actual tensor
space, not a further quotient.  It is therefore
isotypic for the tensor of these components.  Project to one copy
detecting a nonzero value of \(L\).  The resulting nonzero map
\(\ell\) satisfies (BCS.12) with values in this tensor.
Frobenius reciprocity is explicit here:
\[
                   \Phi(v)(g)=\ell(\sigma(g)v).
 \tag{BCS.13}
\]
It is smooth on the right, intertwines right translations, and has
the required unnormalized left covariance by (BCS.12).  The quotient
\(P(F)\backslash G(F)\) is compact, so there is no further
support condition.

For \(d\) equal blocks of size \(a\), the exponent of
\(|\det m_k|\) in \(\delta_P^{1/2}\) is
\[
 \frac{a(d-1-2k)}2=a s_{d,k}.
 \tag{BCS.14}
\]
Hence changing (BCS.13) to normalized induction subtracts
\(a s_{d,k}\) from its coefficient exponent
\((a-1)s_{d,k}\), leaving \(-s_{d,k}\).  This is exactly
(BCS.10).  Evaluation at the identity proves the map is nonzero.
All component maps, blocks and coefficient twists used here are the
ones from the original ascending embedding.
\end{proof}
% END BCS7

% BEGIN BCS8
\paragraph{BCS8. The precisely matched published generic-unitary input.}
\label{claim:BCS8}
For a nonarchimedean characteristic-zero local field, a smooth
irreducible generic unitarisable \(\rho\) has a presentation
\[
       \rho\simeq\delta_1\nu^{t_1}\times\cdots\times
                   \delta_b\nu^{t_b},\qquad |t_j|<\tfrac12,
 \tag{BCS.15}
\]
where the \(\delta_j\) are unitary square-integrable representations.
This is an explicit published input: M\oe glin--Waldspurger,
\emph{Le spectre r\'esiduel de \(\mathrm{GL}(n)\)},
\S\S I.10--I.11, pp.639--640,
\url{https://www.numdam.org/item/10.24033/asens.1595.pdf}.
Their \(\pi,\delta_j,v_j\) are respectively
\(\rho,\delta_j,t_j\); their normalized induction and absolute
value are (BCS.2).  The Whittaker hypothesis eliminates their longer
Speh factors.  All our \(E_w\) have characteristic zero.  BCS1 and
BCS5 supply precisely unitarisability and genericity for the actual
cuspidal local components, so (BCS.15) applies to each of them.
This application retains the strict interval; it does not assume
temperedness or prove a global residual-spectrum classification.
% END BCS8

% BEGIN BCS9
\paragraph{BCS9. The full unique cosocle of the descending Speh induction.}
\label{claim:BCS9}
For every \(\rho\) in BCS8 and every \(d\ge1\), the descending
product in BCS6 has an irreducible cosocle with a unique maximal
proper invariant subspace, in the literal sense of
\(\mathrm{IsCosocle}\).
\begin{proof}
Use the retained \texttt{Square\allowbreak Integrable\allowbreak Segment} to write
\(\delta_j=L(\Delta_j)\) for a segment.  The real center of each
\(\Delta_j\) is zero: its central character is the product of the
central characters of its segment terms, while unitarisability makes
that product have absolute value one.  More explicitly, if the first
term has rank \(r\), segment length \(l\), and central character
\(\omega\), then at a uniformizer \(z\)
\[
 \log|\omega(z)|=-\frac{r(l-1)}2\log|z|_F.
 \tag{BCS.16}
\]
Indeed each of the \(l\) terms contributes
\(\log|\omega(z)|+rh\log|z|_F\), for \(0\le h<l\),
and their sum is zero.  This also handles any unitary supercuspidal
line implicit in the segment notation.

If \(\Delta\nu^x\) precedes \(\Delta'\nu^y\) in the original
segment order and both underlying segments have unitary
square-integrable quotients, then
\[
                             x+1\le y.
 \tag{BCS.17}
\]
To check this exactly, express both on their common supercuspidal
line.  Precedence means their first endpoints differ by an integer
\(j\ge1\), their last endpoints differ by a positive integer,
and the intervals overlap or are adjacent.  If their lengths are
\(l,l'\), the last condition needed here is \(j+l'-l\ge1\).
Subtract (BCS.16) for the two first terms.  It gives
\(y-x=j+(l'-l)/2\), the average of the two positive endpoint
differences.  Therefore \(y-x\ge1\), proving (BCS.17).

Expand each factor of the descending product using (BCS.15), and
apply induction in stages.  The resulting ordered factors are
\[
          L\bigl(\Delta_j\nu^{t_j+s_{d,k}}\bigr),
       \qquad (k,j)=(0,1),\ldots,(0,b),(1,1),\ldots,(d-1,b).
 \tag{BCS.18}
\]
If an earlier segment preceded a later one, (BCS.17) would imply
\[
        t_j+s_{d,k}+1\le t_{j'}+s_{d,k'},\qquad k\le k'.
\]
But \(t_{j'}-t_j<1\) and
\(s_{d,k'}-s_{d,k}=k-k'\le0\), a contradiction.  Thus this is
exactly an ordered list of \(L\)-segment factors as required by
the original \texttt{Standard\allowbreak Cosocle\allowbreak GL}.  That retained theorem
gives its irreducible unique cosocle.  Transporting through the
explicit induction-in-stages equivalence gives the asserted cosocle
of the original descending product.  Its uniqueness is the unique
maximal-kernel assertion of that theorem, so it applies to every
irreducible quotient, not merely one chosen Langlands quotient.
\end{proof}
% END BCS9

% BEGIN BCS10
\paragraph{BCS10. Every original discrete local component is the required Speh.}
\label{claim:BCS10}
Let \(V\) be discrete of positive rank and let
\(\operatorname{AscEmb}(V,P,d)\), with \(P\) normalized.
For every \(w\), every actual local component \(\rho\) of
\(P\), and every actual local component \(\sigma\) of \(V\),
\[
                         \mathrm{IsSpeh}(\rho,d,\sigma).
 \tag{BCS.19}
\]
\begin{proof}
BCS1 supplies the two positive invariant forms.  BCS7 supplies the
nonzero ascending map, and BCS6 gives a surjection from the
descending product onto \(\sigma\).  BCS5--BCS9 give the unique
cosocle of that product.  The kernel of a surjection onto an
irreducible representation is a maximal proper invariant subspace,
so it is the unique kernel of the cosocle map.  The induced quotient
equivalence identifies \(\sigma\) with this cosocle.  This proves
the literal \(\mathrm{IsSpeh}\), for every actual component map
in the statement and for every positive length.
\end{proof}
% END BCS10

% BEGIN BCS11
\paragraph{BCS11. The exact block Capelli shift.}
\label{claim:BCS11}
For the original unnormalized induced automorphic functions of block
ranks \(n_i\), with \(\sum_i n_i=N\), whose coefficient spaces
have infinitesimal characters \(m_i\), the infinitesimal character
on that induced space is
\[
       \sum_i\{z-\mu_i/2:z\in m_i\},\qquad
       \mu_i=\sum_{j>i}n_j-\sum_{j<i}n_j.
 \tag{BCS.20}
\]
\begin{proof}
Here is the algebraic block computation, including its sign.
For operators satisfying the \(\mathfrak{gl}_N\) relations, use
the source column determinant
\[
 q_N(X,u)=\sum_{\omega\in S_N}\operatorname{sgn}(\omega)
       \prod_{k=0}^{N-1}
       \bigl(X_{\omega(k),k}+\delta_{\omega(k),k}(u-k)\bigr),
 \tag{BCS.21}
\]
with factors in increasing column order.  Let \(J\) be the span
of \(X_{pq}A\), where the block of \(p\) is earlier than that of
\(q\).  It is a right ideal and is stable under left multiplication
by Levi entries.  The latter assertion follows immediately from
\([X_{ij},X_{pq}]=\delta_{jp}X_{iq}-\delta_{qi}X_{pj}\): the
commutator is again upper-block when \(i,j\) belong to one block.

For two consecutive blocks of sizes \(r,m\), the determinant
expansion gives
\[
 q_{r+m}(X,u)\equiv
       q_m(X_{22},u)q_r(X_{11},u-m)\pmod J.
 \tag{BCS.22}
\]
One way to see every discarded term is to move the last block before
the first simultaneously in the rows and columns.  The quantum
minor alternating identities preserve the determinant.  They follow
by interchanging adjacent columns in (BCS.21): the commutator of
their two entries is canceled by the difference of their consecutive
shifts; row interchange gives the ordinary alternating sign.  In the
new order, permutations preserving both blocks give exactly the
right side of (BCS.22).  Any other permutation has, in its initial
block of columns, an upper-block entry.  The entries preceding the
first such entry are Levi entries, and hence preserve \(J\) on the
left.  This term is in \(J\).  Entries in the two diagonal blocks
commute, so their two determinant factors may be interchanged.
Induction on blocks proves
\[
 q_N\!\left(X,u+\frac{N-1}{2}\right)
 \equiv\prod_i q_{n_i}\!\left(X_{ii},
           u+\frac{n_i-1}{2}-\frac{\mu_i}{2}\right)\pmod J.
 \tag{BCS.23}
\]
Indeed the shift for block \(i\) is also
\(u+(N-1)/2-\sum_{j>i}n_j\), which is the displayed expression.

Apply this to the right Lie operators on the source induced
functions.  At a Levi point, an upper-block operator kills a left
unipotent-invariant function: right infinitesimal multiplication
becomes left unipotent multiplication after conjugation by that
Levi point.  The same is true after any right differential operator,
since such operators preserve left unipotent invariance.  Thus every
element of \(J\) vanishes there.  On a pure tensor of coefficient
functions, each diagonal-block determinant acts by its specified
polynomial \(\prod_{z\in m_i}(u-\mu_i/2+z)\).  Linearity gives
the same conclusion on their algebraic tensor span.  Apply the
argument to every right translate by \(K_\infty G(\mathbb A_f)\);
adelic Iwasawa decomposition and left unipotent invariance then give
the identity at every point.  Right differentiation preserves these
induced coefficient spaces by the differentiated Iwasawa action,
so successive block determinants are legitimate.  This proves
(BCS.20) on the actual source functions.
\end{proof}
% END BCS11

% BEGIN BCS12
\paragraph{BCS12. The infinity character of the ascending embedding.}
\label{claim:BCS12}
Under the hypotheses of BCS10, for every embedding \(\tau\) there
is a multiset \(m_\tau\) of cardinality \(a\) such that
\[
 P.\mathrm{HasInfChar}(\tau,m_\tau),\qquad
 \mathrm{HasInfCharGL}\bigl(V,\tau,
                    \operatorname{spehInf}(d,m_\tau)\bigr).
 \tag{BCS.24}
\]
Any already supplied infinitesimal character of \(V\) equals this
one exactly.
\begin{proof}
Choose \(m_\tau\) by BCS4.  Twisting by
\(|\det|_{\mathbb A_E}^{t}\) adds \(t\) to each entry of
\(m_\tau\): differentiate the determinant character in each
diagonal Lie operator in (BCS.21).  For the equal blocks in (BCS.7),
\(\mu_k/2=a s_{d,k}\) by (BCS.14).  Therefore BCS11 gives on
the target of the ascending embedding the multiset
\[
             \sum_{k=0}^{d-1}\{z-s_{d,k}:z\in m_\tau\}.
\]
The involution \(k\mapsto d-1-k\) changes \(-s_{d,k}\) to
\(s_{d,k}\), so this is (BCS.3).  The ascending map intertwines
every Lie operator, hence every Capelli polynomial.  Its injectivity
transfers that scalar polynomial to \(V\).  The degree is
\(da=N\).  Uniqueness of the monic scalar polynomial in BCS4 proves
the final assertion; there is no unspecified common scalar twist.
\end{proof}
% END BCS12

% BEGIN BCS13
\paragraph{BCS13. The complete original \texttt{Asc\allowbreak Comp} consumer.}
\label{claim:BCS13}
The preceding assertions prove the literal \texttt{Asc\allowbreak Comp} of
\texttt{Base\allowbreak Change\allowbreak Identities/\allowbreak Narrow.\allowbreak lean} for every totally complex
number field and every \(N\ge1\).
\begin{proof}
For the first clause, a normalized cusp of rank \(a=N\) has
\(V=P.V\) discrete by BCS1.  At length one the unique exponent
is zero, the parabolic is the whole group, and the identity inclusion
is the original \(\operatorname{AscEmb}(V,P,1)\).

For the second clause, take any discrete \(V\), normalized cusp
\(P\), and actual \(\operatorname{AscEmb}(V,P,d)\).
BCS4 provides both local components at every finite place.
BCS10 gives the \texttt{loc} clause for \emph{all} choices of
those components.  BCS12 gives the infinity clause.  These are
exactly the three fields \texttt{ex\allowbreak Loc}, \texttt{loc}, \texttt{inf}
of \(\mathrm{HasSpehComp}(V,P,d)\).  This proves all of
\texttt{Asc\allowbreak Comp}; it does not assert the separate production of an
ascending embedding for an arbitrary discrete \(V\).
\end{proof}
% END BCS13

\subsection{The original spherical parameters and the final witnesses}

% BEGIN BCS14
\paragraph{BCS14. Exact fixed vectors of induction and quotients.}
\label{claim:BCS14}
For smooth local factors \(\tau_i\), evaluation at one identifies
\[
 (\tau_1\times\cdots\times\tau_c)^{K_n}
                  \simeq\bigotimes_i\tau_i^{K_{N_i}}.
 \tag{BCS.25}
\]
Compact fixed vectors are exact on smooth representations.  Hence a
spherical quotient of such an induction forces each nonzero-rank
factor to be spherical.
\begin{proof}
Iwasawa decomposition gives \(G=PK_n\).  A right \(K_n\)-fixed
section is determined by its value at one.  This value is fixed by
the integral Levi because \(\delta_P=1\) on \(P\cap K_n\).
Conversely, such a value defines
\(f(pk)=\delta_P(p)^{1/2}\tau(m_p)v\); two decompositions
differ by \(P\cap K_n\), so the definition is independent of
choices and is a smooth right-fixed section.

On the algebraic tensor space the average over the integral Levi is
the tensor product of the individual averages.  Their image is
exactly the tensor of the fixed spaces, proving (BCS.25).  An
algebraic tensor of nonzero vector spaces over a field is nonzero;
if some fixed space is zero, so is the tensor.

For a surjective smooth map \(T:X\to Y\) and \(y\in Y^K\),
lift \(y\) to \(x\) and average \(x\) over \(K\).  The finite
average has image \(y\).  Injection is automatic, proving exactness.
A nonzero fixed vector in a quotient therefore lifts to a nonzero
fixed vector in the induction, and (BCS.25) proves the assertion.
\end{proof}
% END BCS14

% BEGIN BCS15
\paragraph{BCS15. The normalized finite Satake calculation.}
\label{claim:BCS15}
For smooth irreducible spherical \(\tau_i\), choose their
original Satake multisets \(B_i\subset\mathbb C^\times\).
Then every nonzero fixed vector of their normalized product has
Satake multiset \(\sum_iB_i\).  The original normalization is
\[
 T_j\bigm|_{\tau^{K_r}}=
             q_w^{j(r-j)/2}e_j(B)\qquad(0\le j\le r),
 \tag{BCS.26}
\]
where \(T_j\) is the characteristic function of
\(K_r\operatorname{diag}(\varpi I_j,I_{r-j})K_r\), and
\(\operatorname{vol}(K_r)=1\).
\begin{proof}
We first record why these parameters exist for each irreducible
spherical factor.  The spherical Hecke algebra is commutative:
transpose is a Haar-preserving anti-involution fixing each double
coset, by elementary divisors.  It is generated by
\(T_1,\ldots,T_r,T_r^{-1}\).  To verify generation, write a
dominant integral elementary-divisor tuple as a sum of the tuples
\((1^j,0^{r-j})\) and a central tuple.  The corresponding product
of \(T_j\)'s has the indicated top double coset with nonzero
coefficient and only lower dominance tuples.  This follows by
successively taking the dimensions of the reductions of the
associated lattices.  There are finitely many lower tuples with
the fixed determinant and with minimum entry bounded below, so
induction eliminates them.  Inverses of \(T_r\) supply the central
negative tuples.

The fixed space of an irreducible spherical representation is a
simple module for this algebra.  If a fixed subspace generates the
representation under \(G\), averaging its translates back to
\(K\) gives its Hecke span; thus a nonzero Hecke submodule is the
whole fixed space.  Admissibility makes this space finite dimensional.
A simple finite-dimensional module over a commutative complex algebra
is one dimensional.  Its eigenvalues in (BCS.26) define the
elementary symmetric coefficients of a monic polynomial of degree
\(r\).  Its constant coefficient is nonzero because \(T_r\) is
invertible.  Its roots give a multiset in \(\mathbb C^\times\),
uniquely.  The same calculation works for every uniformizer, since
changing it by a unit does not change any of the double cosets.

Here is the block calculation that fixes the normalization.  For
\(\sum_i j_i=j\), decomposition of the finite flags for
\(K_n\operatorname{diag}(\varpi I_j,I_{n-j})K_n/K_n\) by their
intersections with the block flag gives the normalized constant term
\[
 T_j\longmapsto
 \sum_{\substack{0\le j_i\le N_i\\\sum_i j_i=j}}
 q_w^{\{j(n-j)-\sum_i j_i(N_i-j_i)\}/2}
                         \bigotimes_i T_{i,j_i}.
 \tag{BCS.27}
\]
For explicit counting, the off-diagonal graph coordinates contribute
\(q_w^{\sum_{i<l}j_i(N_l-j_l)}\).  The modulus correction in
normalized induction contributes
\(q_w^{-\frac12\sum_{i<l}(j_iN_l-j_lN_i)}\).
Adding these exponents gives exactly the exponent in (BCS.27).
The diagonal flags give the indicated \(T_{i,j_i}\); these cells
partition the finite quotient, so no multiplicity is omitted.

Under evaluation at one in BCS14, this constant term is the action
on the tensor fixed vector.  Substitute (BCS.26) for each factor.
All powers of \(q_w\) combine to \(q_w^{j(n-j)/2}\), and the
remaining sum is
\(\sum_{\sum j_i=j}\prod_i e_{j_i}(B_i)=e_j(\sum_iB_i)\).
The fixed space is one dimensional by (BCS.25); hence every nonzero
fixed vector has these eigenvalues.  This proves exactly the source
\(\mathrm{HasSatGL}\) condition for the full product.
\end{proof}
% END BCS15

% BEGIN BCS16
\paragraph{BCS16. The exact Speh strings in the inert parameter.}
\label{claim:BCS16}
If \(d\ge1\), \(\mathrm{IsSpeh}(\rho,d,\sigma)\), and
\(\mathrm{HasSatGL}(w,\sigma,B)\), then some multiset
\(S\subset\mathbb C^\times\) satisfies
\[
 \mathrm{HasSatGL}(w,\rho,S),\qquad
 B_{\mathbb C}=\sum_{k=0}^{d-1}
                       \{zq_w^{s_{d,k}}:z\in S\}.
 \tag{BCS.28}
\]
Consequently, if the full product of Speh factors has Satake multiset
\(T\), then
\[
             T_{\mathbb C}=\sum_i\sum_{k=0}^{d_i-1}
                       \{zq_w^{s_{d_i,k}}:z\in S_i\},
 \qquad\mathrm{HasSatGL}(w,\rho_i,S_i).
 \tag{BCS.29}
\]
Here subscript \(\mathbb C\) means the original map from units
to their values, retaining multiplicities.
\begin{proof}
Lift a nonzero fixed vector of \(\sigma\) through its actual
Speh quotient by BCS14.  Formula (BCS.25) makes each
\(\rho\nu^{s_{d,k}}\) spherical.  These unramified determinant
twists are trivial on \(K_a\), so \(\rho\) is spherical.
Choose its unique multiset \(S\) by BCS15.

Twisting by \(\nu^t\) multiplies the Satake roots by
\(q_w^{-t}\): on the \(j\)-th double coset the determinant has
absolute value \(q_w^{-j}\), so its Hecke operator is multiplied
by \(q_w^{-tj}\).  Apply BCS15 to the descending induction and
transport the eigenvalue equations through its quotient map.  Its
root multiset is therefore \(\sum_k q_w^{-s_{d,k}}S\).
The index involution \(k\mapsto d-1-k\) changes this to the right
side of (BCS.28).  The image vector is nonzero, so uniqueness of the
scalar polynomial proves equality with \(B\).  Apply BCS15 first
to the outer product and then (BCS.28) to each factor to get
(BCS.29).  In particular \(q_w\) is the absolute norm of \(w\),
not the separate parameter \(q\) appearing in
\(\mathrm{IsUnramInv}(\varpi,s,q)\).
\end{proof}
% END BCS16

% BEGIN BCS17
\paragraph{BCS17. Positive ranks and distinct equal-length cuspidal classes.}
\label{claim:BCS17}
For any \texttt{Definite\allowbreak N} witness and the ascending cusps it
supplies, write \(N_i=d_i a_i\).  Then \(a_i,d_i\ge1\),
\(c>0\), \(\sum_i d_i a_i=n\), and
\[
 d_i=d_j,\quad
 (a_i,[P_i])=(a_j,[P_j])\text{ in }\coprod_a\mathrm{CG}(E,a)
                         \quad\Longrightarrow\quad i=j.
 \tag{BCS.30}
\]
\begin{proof}
The rank identities follow directly from each ascending embedding;
positive \(N_i\) gives both positive factors.  If \(c=0\), the
rank sum would give \(n=0\), contradicting \(n\ge1\).

At any complex embedding \(\tau\), the original \(\pi\) has a
dominant integral weight \(\lambda\) satisfying
\[
                    \mathrm{HasWeightAt}(\pi,\tau,\lambda).
\]
To see existence
on these actual functions, take the nonzero finite-dimensional
infinity-orbit span of a vector in \(\pi\), restrict it to the
\(U(n)\) factor at \(\tau\), and average an inner product over
that compact group.  A simple constituent is one of the original
highest-weight modules by compact \(U(n)\) highest-weight theory.
Its inclusion is exactly the map in \(\mathrm{HasWeightAt}\),
because the factor inclusion is the original \(cEmb_\tau\)
on \(\operatorname{embU}(n,\tau)\).

The source multiset is
\[
 \operatorname{infOfWeight}(\lambda)=
       \{\lambda_k+(n-1)/2-k:0\le k<n\}.
 \tag{BCS.31}
\]
Its entries are pairwise distinct: for \(k<l\), their difference
is \(\lambda_k-\lambda_l+l-k>0\).  By BCS12 and the unchanged
base-change weight clause,
\[
 \sum_i\operatorname{spehInf}(d_i,m_{i,\tau})
                           =\operatorname{infOfWeight}(\lambda).
 \tag{BCS.32}
\]
Each summand has positive cardinality \(N_i\).  Equality of the
cuspidal classes in (BCS.30) identifies their infinitesimal
characters, by the all-operator equivalence and BCS4.  Equal lengths
then make the two summands in (BCS.32) equal.  If \(i\ne j\),
any entry of this nonempty multiset occurs at least twice in the
sum, contradicting (BCS.31).  This proves the original distinctness
clause without assuming a multiplicity-free base-change decomposition.
\end{proof}
% END BCS17

% BEGIN BCS18
\paragraph{BCS18. The entire original \texttt{Bc\allowbreak Speh}, with unchanged witnesses.}
\label{claim:BCS18}
Under the retained inputs stated at the start of this chapter, the
literal \texttt{Bc\allowbreak Speh} holds.  More precisely, for every original
\(E,n,\pi\) there are \(c>0\), \(a_i,d_i\ge1\),
\(\sum_i d_i a_i=n\), actual normalized unitary conjugate-self-dual
cusps \(P_i\), actual local cusp components \(\rho_{i,w}\),
multisets \(m_{i,\tau}\), and the original discrete local components
\(\sigma_{i,w}\), with all seven fields of \texttt{Parts}, both
local unitarisability assertions, the literal \(\mathrm{IsSpeh}\),
and the full split local-product component assertion.
\begin{proof}
Use BCS3 to obtain the original \texttt{Definite\allowbreak N} data.  From each
\(\mathrm{CsdAsc}(V_i)\), choose \(a_i,d_i,P_i\) and the
supplied ascending embedding.  The global normalization, unitarity
and conjugate self-duality of \(P_i\) are retained from that very
record.  BCS13 supplies its \(\mathrm{HasSpehComp}\).  Choose
\(\rho_{i,w}\) and \(m_{i,\tau}\) from its existence fields;
retain the \emph{original} \(\sigma_{i,w}\) from base change.
The universal local clause of BCS10 applies to those components.
BCS1 supplies the two local forms.  Replacing the equality
\(N_i=d_i a_i\) only relabels finite index sets by their canonical
rank equality; no new representation is substituted.

Here is every field of the target record.
\begin{center}
\begin{tabular}{p{0.16\linewidth}p{0.74\linewidth}}
\texttt{pos} & BCS17 gives \(a_i,d_i\ge1\) and the exact rank sum.\\
\texttt{csd} & The \(P_i\) retain their original conjugate self-duality.\\
\texttt{loc} & The chosen \(\rho_{i,w}\) have the actual cusp component maps.\\
\texttt{inf} & Each \(m_{i,\tau}\) is the actual cusp infinity character.\\
\texttt{weight} & BCS12 replaces every original \(M_{i,\tau}\) by
\(\operatorname{spehInf}(d_i,m_{i,\tau})\); the original weight
identity holds for every \(\tau,\lambda\) satisfying
\(\mathrm{HasWeightAt}\).\\
\texttt{distinct} & BCS17 proves precisely (BCS.30).\\
\texttt{inert} & Apply BCS16 to (BCS.6), obtaining (BCS.29) with
\(T=\operatorname{bcParam}(n,\alpha)\), for every original
\(w,s,\varpi,q,\alpha\) and every original guard.\\
\end{tabular}
\end{center}
The two unitarisability clauses and \(\mathrm{IsSpeh}\) were just
verified for every \(i,w\).  Finally, the split clause of BCS3 is
already \(\mathrm{IsCompAlong}\) for the full product of these
same \(\sigma_{i,w}\) along the same \(\operatorname{locU}\)
and \(cLoc_w\).  It is copied without alteration.  This is the
full existential statement \texttt{Bc\allowbreak Speh}, with all positive
lengths allowed and with the original rank \(n\).
\end{proof}
% END BCS18

% BEGIN BCS19
\paragraph{BCS19. All-unramified spherical cusp factors for the same witnesses.}
\label{claim:BCS19}
The witnesses in BCS18 also satisfy, simultaneously at every finite
place,
\[
 \mathrm{IsUnramU}(\pi,w)\quad\Longrightarrow\quad
           \rho_{i,w}^{K_{a_i}}\ne0\quad\text{for every }i.
 \tag{BCS.33}
\]
The assertion also holds for every alternative actual local component
of the same \(P_i\).
\begin{proof}
The stronger \texttt{Definite\allowbreak N} witness retained in BCS3 makes the
full product of the \(\sigma_{i,w}\) spherical at every such
place.  BCS14 makes each \(\sigma_{i,w}\) spherical.  Lift a
nonzero fixed vector through its actual descending Speh quotient
and apply BCS14 once more.  Each determinant twist of
\(\rho_{i,w}\) is then spherical, and the twists are trivial on
\(K_{a_i}\), proving (BCS.33).  BCS4 identifies any alternative
actual component with the chosen one, and an intertwining
equivalence transports its nonzero fixed vector.  This argument
requires neither a new choice of \texttt{Bc\allowbreak Speh} at different places
nor a split/inert case distinction.
\end{proof}
% END BCS19

\paragraph{Scope of the completed bridge.}
BCS3 proves the original \texttt{Definite\allowbreak N}; BCS13 proves the original
bounded \texttt{Asc\allowbreak Comp}; BCS18 proves the original \texttt{Bc\allowbreak Speh};
BCS19 supplies the additional all-unramified consequence for precisely
the same witnesses.  The named \texttt{Base\allowbreak Change\allowbreak Isobaric\allowbreak N} foundation
remains an input.  The separate global Galois-attachment record,
ordinary Hecke producers, arithmetic row counts, and terminal RH/GRH
obligations are not conclusions of this chapter.


\section{Assembling the original Galois representation of a definite unitary form}
\label{sec:galois-assembly-v44}

\status{This chapter reconstructs \texttt{Gal\allowbreak Unitary\allowbreak T} and its
\texttt{Gu\allowbreak Wrap} consumer.  It uses the original polarized cuspidal
attachment input \texttt{Row\allowbreak T}, the original local representation and
period foundations, and the base-change/Speh construction BCS18--BCS19.
These inputs are identified separately from the proofs below.  The
checked published generic-unitary input is BCS8; GUA4 supplies its
needed elementary ordering consequence.  We preserve every positive rank,
every Speh length, every odd coefficient prime including $3$, and
the universal quantifiers over local correspondences and period
bridges.  In particular, the local conclusion retains the source's
\texttt{Gd} condition.  This is a reconstruction of an installation
input, not an unconditional proof of RH or GRH.}

Let $E$ be a CM number field, $E^+$ its maximal totally real
subfield, and $\Omega/E$ the specified algebraic closure.  Suppose
$E/E^+$ is unramified at every finite place.  Fix $n\geq1$, a prime
$\ell\ne2$, and the given field isomorphism
$\iota:\overline{\mathbb Q}_\ell\simeq\mathbb C$.
Let $\sigma$ be the original irreducible automorphic subspace of
functions on the definite unitary adelic group $U(\mathbb A_{E^+})$,
with the right-translation action; this is \texttt{Is\allowbreak Aut\allowbreak Rep\allowbreak U}.
No continuity of $\iota$ for the complex topology is assumed.

Write $\nu=|\det|_F$ over a finite completion $F=E_w$, with
$|\varpi|_F=q_w^{-1}$.  The source conventions are
\[
 \operatorname{algNorm}(\tau)=\tau\nu^{(1-\dim\tau)/2},\qquad
 x(t)=x.\texttt{lp\allowbreak Tw}(t),\qquad
 \epsilon(\operatorname{Frob}^{\rm ar}_w)=q_w.        \tag{GUA.1}
\]
The twist $x(t)$ multiplies the Weil action on an arithmetic
Frobenius by $q_w^t$; it leaves monodromy unchanged.  Thus its
geometric Frobenius eigenvalues are multiplied by $q_w^{-t}$.
All real powers of positive numbers in this chapter are the
positive real powers, embedded in $\mathbb C$.

For a Weil--Deligne isomorphism class $y=(r_y,N_y)$, put
\[
 \operatorname{Gd}(y)\quad\Longleftrightarrow\quad
 N_y=0\quad\hbox{and}\quad
 \det(X-r_y([u,v]))=(X-1)^{\dim y}
                       \quad(u,v\in W_{E_w}).       \tag{GUA.2}
\]
This is the literal \texttt{Tor}$\wedge$\texttt{mono0} predicate.
It is not being strengthened to a specified diagonalization.
For every valuation $A$ of $\Omega$ over $w$ and every map
$\mathrm{rc}:\operatorname{Irr}\operatorname{GL}(E_w)
\longrightarrow\operatorname{WD}(A)$ satisfying the full original
\texttt{Is\allowbreak Rec\allowbreak T}, the rank-one law uses that correspondence's own
local Artin map $a:W_A\to E_w^\times$.  Its degree convention is
$\operatorname{ord}_w(a(g))=-\deg(g)$.

We use the following exact compatibility condition.  A triple
$(R,p,H)$ satisfies \texttt{Compat\allowbreak G} if $R$ is continuous and
semisimple and the following statements hold for every $w$.
First, for every $A,h,\mathrm{rc}$ as above, at $w\nmid\ell$
there is an actual Weil--Deligne parameter $x$ of $R|_{D_A}$
such that
\[
             \operatorname{Gd}(p_{w,A,\mathrm{rc}})
                    \ \Longrightarrow\ x=p_{w,A,\mathrm{rc}}.
                                                        \tag{GUA.3}
\]
Second, for every original coefficient ring $O$, local bridge
$B_w$ and witness \texttt{Induces}, goodness of that same
parameter implies potential semistability in the specified
period model $\texttt{Loc.\allowbreak PP}(\ell)$ and the exact equality
\[
 \operatorname{wdAt}(R|_{D_A},B_w,\iota)
                              =p_{w,A,\mathrm{rc}}.       \tag{GUA.4}
\]
Finally, if the parameters are good for \emph{all}
$A,h,\mathrm{rc}$ at $w$, the representation is de Rham for
every such bridge, and its filtration jumps at every induced
embedding are the multiset $H_\tau$; this is
\texttt{Has\allowbreak HTAt}.  The quantification includes every finite
coefficient extension and every inducing embedding appearing
in that definition.  No existence of a bridge at an unrelated
residue prime is asserted.

\subsection{Local parameter identities for every permitted correspondence}

% BEGIN GUA1
\paragraph{GUA1. Goodness under finite sums and unramified twists.}
\label{claim:GUA1}
For Weil--Deligne classes $x,y$ and $t\in\mathbb R$,
\[
 \operatorname{Gd}(x\oplus y)
       \Longleftrightarrow\operatorname{Gd}(x)\wedge\operatorname{Gd}(y),
 \qquad
 \operatorname{Gd}(x(t))\Longleftrightarrow\operatorname{Gd}(x).
\]
The corresponding statement holds for every finite list,
including the empty list.  An unramified class is good.
\begin{proof}
Choose representatives.  Monodromy on a sum is the block
matrix $N_x\oplus N_y$, which vanishes exactly when both
blocks vanish.  For each commutator the characteristic
polynomial is the product of the two block characteristic
polynomials.  A product of two monic complex polynomials is
$(X-1)^{a+b}$ exactly when its degree-$a$ and degree-$b$
factors are $(X-1)^a$ and $(X-1)^b$: all their complex roots
must be $1$.  This proves the first equivalence.  An
unramified twist factors through the degree map
$W_A\to\mathbb Z$, so it is $1$ on every commutator and
does not change $N$.  Induction proves the list statement;
the zero class satisfies both conditions.  Finally a
commutator has degree zero and hence lies in inertia.  An
unramified class is trivial there and has zero monodromy.
These arguments are invariant under change of representative.
\end{proof}
% END GUA1

% BEGIN GUA2
\paragraph{GUA2. The full \texttt{Plain\allowbreak Gd} statement.}
\label{claim:GUA2}
For every $A,h,\mathrm{rc}$ satisfying \texttt{Is\allowbreak Rec\allowbreak T}, every
rank $a\geq0$ and every smooth irreducible $\tau$ of
$\operatorname{GL}_a(E_w)$,
\[
 \texttt{Plain\allowbreak Tor}(\operatorname{cl}\tau)
            \Longrightarrow
       \operatorname{Gd}\bigl(\mathrm{rc}(\operatorname{algNorm}\tau)\bigr).
                                                        \tag{GUA.5}
\]
\begin{proof}
The original \texttt{Plain\allowbreak Tor} says that the cuspidal support
consists of rank-one characters $\xi_1,\ldots,\xi_a$ and is
unlinked.  Thus the singleton segments are pairwise unlinked;
the retained Bernstein--Zelevinsky irreducibility criterion
makes their normalized principal series irreducible.  The
support condition makes $\tau$ a subquotient of this principal
series, hence isomorphic to it.  Order the singleton segments
so none precedes a later one.  The standard-cosocle theorem
then identifies this same representation with their Langlands
quotient.  There is no requirement that the characters be
distinct: repeated singletons do not form linked segments.

Write $a_{\mathrm{rc}}$ for the local Artin map contained in
the given \texttt{Is\allowbreak Rec\allowbreak T}.  Its rank-one and Langlands laws
give the equality of actual isomorphism classes
\[
 \mathrm{rc}(\operatorname{cl}\tau)
       =\bigoplus_{j=1}^a (\xi_j\circ a_{\mathrm{rc}},0).
\]
Each character is trivial on commutators and has zero
monodromy.  GUA1 proves goodness.  The real-twist law of
that same correspondence identifies the parameter of
$\operatorname{algNorm}\tau$ with the twist by $(1-a)/2$;
GUA1 applies once more.  Empty induction and the dimension
law give the zero parameter when $a=0$.  The construction
used no choice of a preferred correspondence or Artin map,
so it proves the stated universal quantifier.
\end{proof}
% END GUA2

% BEGIN GUA3
\paragraph{GUA3. Irreducible normalized products have summed parameters.}
\label{claim:GUA3}
If an irreducible smooth representation $\tau$ is isomorphic
to a finite normalized parabolic product of irreducibles
$\tau_i$, then every \texttt{Is\allowbreak Rec\allowbreak T} satisfies
\[
          \mathrm{rc}(\operatorname{cl}\tau)
                      =\bigoplus_i\mathrm{rc}(\operatorname{cl}\tau_i).
                                                        \tag{GUA.6}
\]
In particular this proves the literal \texttt{Prod\allowbreak T}, with
the original block ordering and dimension equality.
\begin{proof}
For two factors choose their Langlands multisegments
$\mathfrak m,\mathfrak n$ using the retained exhaustion
theorem.  Order $\mathfrak m+\mathfrak n$ and form its
standard module.  The retained standard-cosocle theorem
gives its irreducible quotient $L(\mathfrak m+\mathfrak n)$.
The precise multisegment product theorem places this quotient
as a subquotient of $L(\mathfrak m)\times L(\mathfrak n)$.
If the latter is irreducible the two are isomorphic.
Now apply the Langlands law in \texttt{Is\allowbreak Rec\allowbreak T}.  Its
\texttt{seg\allowbreak Sum} attaches the segment parameter to each
member of the list and takes their direct sum.  Concatenating
the two lists concatenates these blocks, so it gives (GUA.6).
For several factors, the same product theorem repeatedly
puts the Langlands quotient of the total multisegment into
the full induction; irreducibility of that full induction
again identifies it with $\tau$.  This argument does not
require an unstated irreducibility of an intermediate product.
The empty case has dimension zero and the unique zero
Weil--Deligne class.  Reassociation and the source's
$\texttt{std\allowbreak Blocks}$ coordinate permutation give an actual
intertwining isomorphism, and pass to the same WD class.
\end{proof}
% END GUA3

% BEGIN GUA4
\paragraph{GUA4. The required generic unitary factorization and ordering bound.}
\label{claim:GUA4}
Let $F$ be any finite extension of $\mathbb Q_p$ and let
$\rho$ be a smooth irreducible, unitarisable, generic
representation of $\operatorname{GL}_a(F)$, $a\geq1$.
There is an actual normalized-product isomorphism
\[
       \rho\simeq\mathop{\times}_{j=1}^b\delta_j\nu^{s_j},
       \qquad -\tfrac12<s_j<\tfrac12,                 \tag{GUA.7}
\]
where every $\delta_j$ is unitary square integrable and is
represented by a segment $\Delta_j$.  If two unitary
segment representations correspond to $\Delta,\Delta'$,
then
\[
            \Delta\nu^x\text{ precedes }\Delta'\nu^y
                       \quad\Longrightarrow\quad x+1\leq y.
                                                        \tag{GUA.8}
\]
\begin{proof}
BCS8 is the single checked published input for (GUA.7):
M\oe glin--Waldspurger, \emph{Le spectre r\'esiduel de
$\mathrm{GL}(n)$}, \S\S I.10--I.11, pp.~639--640.
The field here is nonarchimedean of characteristic zero,
the representation is complex, smooth, irreducible,
generic and unitarisable, and these are precisely its
hypotheses.  No rank bound or exclusion of a residue
prime occurs.  Retain its actual normalized-product
isomorphism and its strict inequalities.  The retained
square-integrable segment theorem supplies each $\Delta_j$.

For completeness, the numerical ordering bound is elementary.
Let $r,r'$ be the lengths of $\Delta,\Delta'$ and $d$ their
common cuspidal rank when the shifted segments are linked.
If $\omega$ is the central character of the first term of
$\Delta$, unitarity of its segment representation gives,
on a uniformizer,
\[
 \log|\omega(\varpi)|
          =-d\frac{r-1}{2}\log|\varpi|_F.            \tag{GUA.9}
\]
Indeed its central character is the product of the $r$
cuspidal central characters at shifts $0,\ldots,r-1$;
its norm is $1$, and $\sum_{j<r}j=r(r-1)/2$.
Precedence means that the first term of
$\Delta'\nu^y$ is the $j$th term along $\Delta\nu^x$,
with $j\geq1$, and that $r<j+r'$, because the second
segment is not contained in the first.  Applying (GUA.9)
to both segments gives
$x+j-y=(r-r')/2$.  Since $j\geq1$ and
$j+r'-r\geq1$, it follows that
$y-x=(j+j+r'-r)/2\geq1$.  This proves (GUA.8).
\end{proof}
% END GUA4

% BEGIN GUA5
\paragraph{GUA5. The local Speh parameter, with every length.}
\label{claim:GUA5}
Let $P$ be an original cuspidal representation of
$\operatorname{GL}_a(\mathbb A_E)$ over a totally complex
number field, let $\rho=P_w$ be a specified unitary local
component, and let $d\geq1$.  If the actual smooth
irreducible $\eta$ is the source's \texttt{Is\allowbreak Speh}
quotient of
\[
           \rho\nu^{(d-1)/2}\times\rho\nu^{(d-3)/2}
                                  \times\cdots\times\rho\nu^{-(d-1)/2},
\]
then for every $A,h,\mathrm{rc}$ with \texttt{Is\allowbreak Rec\allowbreak T},
\[
 \mathrm{rc}(\operatorname{cl}\eta)
       =\bigoplus_{k=0}^{d-1}
          \mathrm{rc}(\operatorname{cl}\rho)((d-1)/2-k).
                                                        \tag{GUA.10}
\]
\begin{proof}
BCS5 proves that the original cusp $P$ is globally generic
and that its specified local component is generic.
One can see the local implication directly: restrict a
nonzero global Whittaker functional to a finite sum of
images of the local component; one image has nonzero
restriction, furnishing the local Whittaker functional.
Apply GUA4 and write $\rho$ as the product (GUA.7).
If $\Delta_j\nu^{s_j}$ preceded
$\Delta_{j'}\nu^{s_{j'}}$, (GUA.8) would give
$s_j+1\leq s_{j'}$, impossible since both exponents
lie in $(-1/2,1/2)$.  Consequently this product is a
standard module in any ordering of these factors; its
unique quotient is $\rho$ itself.

Use induction in stages to replace the displayed Speh
induction by the induction of the factors
\[
       \delta_j\nu^{s_j+(d-1)/2-k}
                  \qquad(0\leq k<d,\ 1\leq j\leq b),
\]
ordered first by increasing $k$.  If an earlier segment
preceded a later one, GUA4 would require
\[
 s_j+(d-1)/2-k+1
                  \leq s_{j'}+(d-1)/2-k',\qquad k\leq k'.
\]
This implies $1\leq s_{j'}-s_j-(k'-k)<1$, a
contradiction.  The standard-cosocle theorem therefore
identifies every irreducible quotient of this actual
induction, including the specified $\eta$, with its
Langlands quotient.  Apply the Langlands and real-twist
laws of the given \texttt{Is\allowbreak Rec\allowbreak T}.  The segment sum of
the flattened list is the sum over $k$ of the segment
sum for $\rho$, twisted by $(d-1)/2-k$.
Tensoring a segment parameter with a scalar unramified
character commutes with its special block and with finite
sums, so this is exactly (GUA.10), with its monodromy,
not just an equality of characteristic polynomials.
\end{proof}
% END GUA5

% BEGIN GUA6
\paragraph{GUA6. The total parameter identity \texttt{Rec\allowbreak Tot\allowbreak T}.}
\label{claim:GUA6}
Let $a_i,d_i\geq1$, $\sum_i d_i a_i=n$; let
$\rho_{i,w}$ and $\eta_{i,w}$ be the unitary cusp and
Speh local representations supplied by BCS.  Put
\[
 e_{ik}=d_i-2k-n,\qquad
 p_{w,A,\mathrm{rc}}
   =\bigoplus_i\bigoplus_{k=0}^{d_i-1}
                 \mathrm{rc}(\operatorname{cl}\rho_{i,w})(e_{ik}/2).
                                                        \tag{GUA.11}
\]
If an irreducible $\tau$ is isomorphic to the full original
normalized product $\mathop{\times}_i\eta_{i,w}$, then
for every permitted $A,h,\mathrm{rc}$,
\[
                  p_{w,A,\mathrm{rc}}
                       =\mathrm{rc}(\operatorname{algNorm}\tau).
                                                        \tag{GUA.12}
\]
\begin{proof}
GUA3 gives the parameter of the full product as the sum
of the parameters of its actual factors.  Substitute
GUA5 for each factor.  The \texttt{Is\allowbreak Rec\allowbreak T} twist law
then applies to the single normalization twist
$(1-n)/2$.  Scalar twists distribute over block sums,
and
\[
        \frac{d_i-1}{2}-k+\frac{1-n}{2}
                              =\frac{d_i-2k-n}{2}.
\]
This is (GUA.11), in the same order as the original
\texttt{tot\allowbreak Par}.  The dimension reindexing in
$\sum_i d_i a_i=n$ is a permutation of the explicitly
ordered block basis.  It cannot change the WD class.
\end{proof}
% END GUA6

\subsection{Polarized cuspidal attachment and half shifts}

% BEGIN GUA7
\paragraph{GUA7. Exactly the toral branch of \texttt{Row\allowbreak T}.}
\label{claim:GUA7}
Assume the original \texttt{Row\allowbreak T}.  Let $P$ be regular
algebraic cuspidal of rank $a\geq1$, polarized by the
original Hecke character $\chi$ of $E^+$, with
specified local components $\rho_w$.  There are $r,H$
with
\[
 (r,\,\mathrm{rc}(\operatorname{algNorm}\rho_w),\,H)
                                      \text{ satisfying \texttt{Compat\allowbreak G}},
\]
and
\[
   \iota H_\tau
       =\left\{\frac{a-1}{2}-z:
             z\in\operatorname{Inf}(P,\iota\tau)\right\}.
                                                        \tag{GUA.13}
\]
Here $\iota H$ means the integer multiset mapped to $\mathbb C$.
\begin{proof}
The source's \texttt{Is\allowbreak Polarized} first gives the sign
$\chi_v(-1)=(-1)^a$ and an isomorphism with the conjugate
inverse-transpose twist.  The retained
Gelfand--Kazhdan duality identifies inverse transpose
with the smooth contragredient.  Pair the resulting
contragredient with $P$ by evaluation.  This is
nondegenerate in each variable; finite adelic operators
act contragrediently, compact operators act by the
corresponding inverse, and differentiating gives the
sum-zero identity for Lie operators.  Thus it is
the literal \texttt{Polarized\allowbreak Pair} required by
\texttt{Row\allowbreak T}, including its sign.

Apply \texttt{Row\allowbreak T}, which is \texttt{Row0b(Is\allowbreak Rec\allowbreak T)},
and retain its actual $r$, $H$ and all fields of
\texttt{Props\allowbreak B}.  Continuity, semisimplicity and
(GUA.13) are three of those fields.  For $w\nmid\ell$
take the actual monodromy parameter $\operatorname{wdOf}(r)$.
The \texttt{away} field gives \texttt{Away\allowbreak B} against
the parameter of every permitted $\mathrm{rc}$.  Its
\texttt{tor} clause says these are equal when
\texttt{Tor} and $N=0$ hold, proving (GUA.3).

For every $O,B_w,\texttt{Induces}$ use the \texttt{above}
field with the original tied period model
$\texttt{Loc.\allowbreak PP}(\ell)$.  Its \texttt{nice} alternative
contains \texttt{Tor} as a separate disjunct, followed
by $N=0$.  That branch gives de Rhamness, potential
semistability, every prescribed filtration-jump statement,
and \texttt{Away\allowbreak B} for the actual $\operatorname{wdAt}$.
Applying the latter's toral clause proves (GUA.4).
If the parameters are good for all correspondences,
the retained \texttt{Rec\allowbreak TExists} supplies one at each
$A,h$ to apply this branch; its jump conclusion does
not depend on which one was chosen.  This proves
\texttt{Has\allowbreak HTAt}.  No use was made of the other
\texttt{nice} branch, which asks for $\mathbb Q_p$
and a rank-dependent lower bound on $p$.  In particular
this proof imposes no such bound and no exclusion of $\ell=3$.
\end{proof}
% END GUA7

% BEGIN GUA8
\paragraph{GUA8. The half-shift construction \texttt{Half\allowbreak G}.}
\label{claim:GUA8}
Let $P$ be a conjugate-self-dual cusp of rank $a\geq1$,
with actual local components $\rho_w$ and infinity
multisets $m_\phi$ satisfying $P.\texttt{Has\allowbreak Inf\allowbreak Char}(\phi,m_\phi)$.
Suppose $e\in\mathbb Z$, every $m_\phi$ has distinct
entries, and $z+e/2\in\mathbb Z$ for every entry.
There are $r,H$ satisfying \texttt{Compat\allowbreak G} with
\[
       p_{w,A,\mathrm{rc}}
                       =\mathrm{rc}(\operatorname{cl}\rho_w)(e/2),
       \qquad
       \iota H_\tau=\{-z-e/2:z\in m_{\iota\tau}\}.
                                                        \tag{GUA.14}
\]
\begin{proof}
Set $s=(e+a-1)/2$ and form the \emph{actual}
automorphic twist $P_s=P\otimes|\det|_{\mathbb A_E}^{s}$.
Multiplication of functions by the determinant character
is an invertible linear map onto this cusp subspace;
constant terms remain zero, and right translation is
twisted by $|\det|^s$.  Its restriction at $w$ therefore
carries the given local intertwiner for $\rho_w$ to
an intertwiner for $\rho_w\nu^s$.  At infinity the
infinitesimal character is $m_\phi+s$.  These entries
are distinct and lie in $\mathbb Z+(a-1)/2$, so this
is regular algebraic, with the source normalization
also in rank $1$.

Write $\xi=|\cdot|_{\mathbb A_E}^{s}$ and let
$\operatorname{res}_{E^+}\xi$ be its restriction to
the ideles of $E^+$.  Its norm pullback is
$\xi\xi^c=\xi^2$.  Let $\delta_{E/E^+}$ be the
quadratic Hecke character.  It is trivial on norms
and has value $-1$ at $-1$ at every real place.
Consequently
\[
        \chi=(\operatorname{res}_{E^+}\xi)
                                  \delta_{E/E^+}^{a}
\]
has $\chi_v(-1)=(-1)^a$ and norm pullback $\xi^2$.
Transport the given conjugate-self-duality intertwiner
of $P$ across multiplication by $\xi$.  The identity
$\xi(\xi\xi^c)^{-1}=(\xi^c)^{-1}$ verifies, for
each finite, compact and Lie operator, precisely the
conjugate-dual twist identity for $P_s$ and $\chi$.
Thus $P_s$ has the polarization used in GUA7.

Apply GUA7.  For each individual permitted $\mathrm{rc}$,
its real-twist law gives
\[
 \mathrm{rc}(\operatorname{algNorm}(\rho_w\nu^s))
  =\mathrm{rc}(\operatorname{cl}\rho_w)
                       \bigl(s+(1-a)/2\bigr)
  =\mathrm{rc}(\operatorname{cl}\rho_w)(e/2).
\]
Substituting this equality preserves every compatibility
and conditional-goodness quantifier.  Formula (GUA.13)
becomes $(a-1)/2-(z+s)=-z-e/2$, proving (GUA.14).
The representation was produced by the original
attachment theorem applied to $P_s$.  This construction
does not require a square root of the cyclotomic
character as a Galois character.
\end{proof}
% END GUA8

% BEGIN GUA9
\paragraph{GUA9. Infinity weights and the exact Speh strings.}
\label{claim:GUA9}
At every embedding $\phi:E\to\mathbb C$ there is a
dominant integral weight $\lambda$ with
$\texttt{Has\allowbreak Weight\allowbreak At}(\sigma,\phi,\lambda)$.  Suppose
the BCS infinity identity is
\[
 \coprod_i\coprod_{k=0}^{d_i-1}
       \{z+(d_i-1)/2-k:z\in m_{i,\phi}\}
     =\operatorname{infOfWeight}(\lambda)
     =\{\lambda_j+(n-1)/2-j:0\leq j<n\}.           \tag{GUA.15}
\]
Then each $m_{i,\phi}$ has distinct entries, and for
each $k<d_i$ all $z+e_{ik}/2$ are integers.  Moreover
\[
 \coprod_{i,k}\{-z-e_{ik}/2:z\in m_{i,\phi}\}
   =\{(n-1)/2-z:z\in\operatorname{infOfWeight}(\lambda)\}.
                                                        \tag{GUA.16}
\]
\begin{proof}
We first justify the existence of a weight on the actual
automorphic function space.  The infinity orbit of a
nonzero function spans a nonzero finite-dimensional space
by \texttt{Is\allowbreak Aut\allowbreak Rep\allowbreak U}.  Evaluation at group elements
separates this space, so finitely many evaluations form
a basis of its dual.  Continuity of the functions and
of right translation therefore proves continuity of
its matrix representation on the compact factor
$U(n)$ at $\phi$.  Averaging an inner product splits
off a finite-dimensional irreducible continuous module $S$.

Here the compact highest-weight classification can be
obtained from the already proved HP4 without introducing
a new automorphic theorem.  The restrictions of regular
functions on $\operatorname{GL}_n(\mathbb C)$ form a
self-conjugate algebra separating points of $U(n)$:
on $U(n)$, $\overline{g}_{ij}=(g^{-1})_{ji}$.
Stone--Weierstrass makes them dense in the continuous
functions.  A nonzero matrix coefficient of $S$ thus
has nonzero Haar pairing with some such regular function.
The finite-dimensional right-translation space generated
by that function is a rational $\operatorname{GL}_n$
module.  The Haar pairing, or its matrix-coefficient
projector, gives a nonzero $U(n)$-map from this space
to $S$; invariant orthogonal complements split it.
The group $U(n)$ is Zariski dense in
$\operatorname{GL}_n(\mathbb C)$: the Cayley map from
skew-Hermitian matrices has real coordinates whose
complexification is the full matrix space, and a complex
polynomial vanishing on all those real coordinates
vanishes identically.  Hence the resulting irreducible
$U(n)$ summand is an irreducible rational
$\operatorname{GL}_n$ submodule.  HP4 identifies it
with the literal polynomial module
$\texttt{sq\allowbreak Rep}(\lambda)$ for a dominant integral
$\lambda$.  Its nonzero intertwiner into the original
infinity orbit is exactly the definition of
\texttt{Has\allowbreak Weight\allowbreak At}.  This argument also preserves
the chosen embedding, including its conjugate orientation.

The right side of (GUA.15) has distinct entries:
if $j<j'$, dominance gives
$\lambda_j-j>\lambda_{j'}-j'$.  A submultiset of a
multiset with distinct entries has distinct entries,
and translation is injective; since $d_i\geq1$ this
proves the first assertion.  If an entry on the left
equals the $j$th entry on the right, subtract
$(n-1)/2$ to obtain
\[
          z+(d_i-2k-n)/2=\lambda_j-j\in\mathbb Z.
\]
Apply $u\mapsto(n-1)/2-u$ to both sides of
(GUA.15) to obtain (GUA.16), with multiplicities.
The argument applies to every weight witness in
(GUA.15), not only the one used to prove existence.
\end{proof}
% END GUA9

\subsection{One representation and its actual period modules}

% BEGIN GUA10
\paragraph{GUA10. Finite direct sums preserve the full compatibility record.}
\label{claim:GUA10}
Given finitely many triples $(r_j,p_j,H_j)$ satisfying
\texttt{Compat\allowbreak G}, their actual block sum satisfies
\texttt{Compat\allowbreak G} with parameter $\bigoplus_jp_j$
and jump multiset $\coprod_jH_j$.  If there is at least
one block and all block ranks are positive, irreducibility
of the sum forces exactly one block, which is irreducible.
This proves \texttt{Sum\allowbreak HTAt}, \texttt{Sum\allowbreak Two\allowbreak G} and
the full dimension-reindexed \texttt{Sum\allowbreak Fam\allowbreak G}.
\begin{proof}
Use the ordered direct sum of the given coordinate spaces;
on it define $R(g)=\operatorname{diag}(r_j(g))$.
Entries and inverse entries are continuous functions of
those of the constituents.  It is semisimple because a
finite direct sum of semisimple modules is semisimple.
For the remaining assertions, the coordinate inclusions
$u_j$ and projections $v_j$ satisfy
\[
 v_j u_k=\delta_{jk},\qquad \sum_j u_jv_j=1,
\]
and intertwine the actual Galois actions.

Here are the needed period compatibilities explicitly.
Choose a common finite coefficient extension for the
finitely many models in any period assertion.  It is
obtained by adjoining their coefficient fields; extension
of a finite coefficient field commutes with fixed spaces,
since a basis of the larger field expands the fixed
equations coefficient by coefficient and the coefficient
Galois action is trivial.  Tensor the displayed maps with
the identity of the specified period ring.  They commute
with the \emph{diagonal} Galois action: both the
representation coordinate and the period coordinate
transform on each side.  Thus the same two identities
hold on invariants, and give
\[
 D_{\rm dR}(R)=\bigoplus_jD_{\rm dR}(r_j),\qquad
 D_{\rm pst}(R)=\bigoplus_jD_{\rm pst}(r_j).        \tag{GUA.17}
\]
For $D_{\rm pst}$ this is first true at every common
finite extension of the local field.  In the union over
such extensions, finitely many components of an element
are all defined at one common extension, so (GUA.17)
continues to hold.  The same proof applies to the
maximal-unramified scalar version in \texttt{Is\allowbreak Pst}.
It works in the source's actual tied period model,
using its own rings and actions.

The maps also commute with $\varphi$ and $N$, which
act on the period coordinate, and preserve each
filtration level.  The de Rham and potentially
semistable dimension equalities therefore add.
Every graded filtration dimension adds at each
specified coefficient embedding, giving multiset union
of the jumps for \emph{all} embeddings in
\texttt{Has\allowbreak HTAt}.  The WD construction from
$D_{\rm pst}$, including Frobenius semisimplification,
is consequently the block sum of the constituent
constructions.  Semisimplification commutes with a
finite block sum, since the Jordan decomposition of a
block-diagonal matrix is block diagonal.

Away from $\ell$, choose a common finite coefficient
field (ADK1) and a common open inertia subgroup for
the finitely many monodromy descriptions.  On it the
representations are $\exp(t_\ell(g)N_j)$.
Their block sum is $\exp(t_\ell(g)\bigoplus_jN_j)$;
removing this exponential and choosing the same
Frobenius gives the sum of their Weil representations.
This proves the asserted sum identity for the actual
\texttt{Is\allowbreak WDOf} parameters.

To check the conditional statements, GUA1 says goodness
of the total parameter implies goodness of every
summand for that same $A,h,\mathrm{rc}$.  The preceding
period and WD identities therefore prove (GUA.3)--(GUA.4).
If the total is good for every such triple, each constituent
satisfies the premise of its own \texttt{Has\allowbreak HTAt};
(GUA.17) proves the total jump assertion.  This keeps
both universal quantifiers in their original order.

Finally, if two positive blocks occur, the first block's
coordinate space is a nonzero proper stable subspace.
If exactly one block occurs but is reducible, its
nonzero proper stable subspace is also one for the total.
Reindexing the block coordinates by the specified
equality of their rank sum with $n$ is a fixed linear
permutation and transports all these statements.
\end{proof}
% END GUA10

% BEGIN GUA11
\paragraph{GUA11. Construction of the one global $R$ and $H$.}
\label{claim:GUA11}
Choose the BCS witnesses for the original $\sigma$:
finitely many unitary conjugate-self-dual cusps $P_i$
of ranks $a_i\geq1$, lengths $d_i\geq1$, actual
local cusp and Speh components $\rho_{i,w},\eta_{i,w}$,
and actual infinity multisets $m_{i,\phi}$, with
$\sum_i a_i d_i=n$.  There are a continuous semisimple
representation
\[
           R:G_E\longrightarrow\operatorname{GL}_n(
                                  \overline{\mathbb Q}_\ell)
\]
and integer multisets $H_\tau$ satisfying
\texttt{Compat\allowbreak G} for (GUA.11).  Every infinity-weight
witness $\lambda$ at $\iota\tau$ satisfies
\[
 \iota H_\tau
       =\{(n-1)/2-z:z\in\operatorname{infOfWeight}(\lambda)\}.
                                                        \tag{GUA.18}
\]
If $R$ is irreducible, then there is just one $i$ and
its $d_i$ is $1$.
\begin{proof}
The number of cusps is positive since otherwise their
positive rank sum would be $0=n$.  For each pair $(i,k)$,
$0\leq k<d_i$, GUA9 proves that $m_{i,\phi}$ and
$e=e_{ik}$ satisfy GUA8 at every $\phi$; existence
of an infinity weight was proved for each embedding
before making this inference.  Apply GUA8 to this
particular cusp, its entire family of specified local
components, and this one integer $e_{ik}$.  It gives
an actual representation $r_{ik}$ of rank $a_i$ and
multisets $H_{ik,\tau}$, simultaneously at all places
and all correspondences.

Take the ordered block sum of these representations,
first over $k$, then over $i$, and transport the basis
by $\sum_i d_i a_i=n$.  Define
$H_\tau=\coprod_{i,k}H_{ik,\tau}$.  GUA10 proves
the entire compatibility record for this one $R$.
The two successive sums give exactly the two lists in
the source's \texttt{tot\allowbreak Par}, with no change of map
between different local conditions.  Substitution of
(GUA.14) and then (GUA.16) proves (GUA.18).
All blocks have positive rank.  Hence GUA10 first
forces one outer block, and then one inner block,
under irreducibility.  These are precisely $c=1$
and $d_0=1$; they were not premises of the construction.
\end{proof}
% END GUA11

\subsection{All unramified places and every split local component}

% BEGIN GUA12
\paragraph{GUA12. Spherical normalization for every \texttt{Is\allowbreak Rec\allowbreak T}.}
\label{claim:GUA12}
If a smooth irreducible $\rho$ of $\operatorname{GL}_a(E_w)$
has normalized Satake multiset $S=\{\alpha_1,\ldots,\alpha_a\}$,
then for every permitted $A,h,\mathrm{rc}$ the class
$\mathrm{rc}(\operatorname{cl}\rho)$ is unramified,
has $N=0$, and has arithmetic Frobenius eigenvalues
$\alpha_1^{-1},\ldots,\alpha_a^{-1}$.
\begin{proof}
We spell out the needed spherical specialization of the
retained local classification, rather than assuming that
a quotient of a spherical induction is automatically
spherical.  BCS14 proves exactness of maximal-compact
fixed vectors and their tensor formula for induction.
BCS15 proves that the spherical Hecke algebra is generated
by $T_1,\ldots,T_a,T_a^{-1}$, computes its normalized
characters, and proves that an irreducible spherical
fixed space is one dimensional.

We also need uniqueness for a given Hecke character,
which has the following algebraic proof.  Let
$X=\mathrm{c\!\operatorname{-}Ind}_K^G\mathbf1$ and let
$\mathcal H=\operatorname{End}_G(X)$, with the right
module convention on $X$.  For a Hecke character $\lambda$,
$M=X\otimes_{\mathcal H}\mathbb C_\lambda$ is generated
by its single $K$-fixed line.  Averaging by the idempotent
$e_K$ shows that $M^K=\mathbb C_\lambda$.  Every proper
submodule has zero $K$-fixed space, since otherwise it
contains this generating line.  The sum of all proper
submodules still has zero fixed space (averaging a
finite sum is the sum of the averages), hence is proper.
Its quotient is the unique irreducible spherical quotient.
Every irreducible spherical representation with character
$\lambda$ is a quotient of $M$ by the compact-induction
adjunction.  Consequently two such representations are
isomorphic.  The unramified principal series of the
characters $\xi_j(\varpi)=\alpha_j$ has this character
by BCS15; its fixed-generated submodule has this unique
quotient.  Thus $\rho$ is a subquotient of this actual
principal series.

Choose the Langlands multisegment of $\rho$ by the
retained exhaustion theorem.  Its standard-module
surjection and BCS14 force each segment representation
to be spherical.  A spherical segment representation is
a quotient of the induction of its cuspidal terms, so
BCS14 makes each such term spherical.  A spherical
supercuspidal is an unramified character: the preceding
principal-series realization gives its support as a
list of unramified rank-one characters, whereas its
supercuspidal support is the singleton consisting of
itself.  The retained support uniqueness theorem forces
rank one and that character.

It remains to rule out a spherical segment of length
$t>1$.  Write its first term as the unramified character
$\mu$, so its induction has characters
$\mu,\mu\nu,\ldots,\mu\nu^{t-1}$.  The determinant
character $\eta\circ\det$, where
$\eta=\mu\nu^{(t-1)/2}$, embeds in that induction:
the function $g\mapsto\eta(\det g)$ satisfies its
normalized covariance exactly.  Both this character
and the purported spherical segment have the same
spherical Hecke character, by BCS14--BCS15.  The
uniqueness just proved identifies them.  But the
Langlands multisegment of the determinant character
consists of the $t$ singleton characters in descending
order.  To verify this last fact, embed the inverse
determinant character in its ascending induction by
the same formula, take smooth contragredients, and use
the retained normalized-induction duality; this yields
a surjection of the descending singleton induction
onto $\eta\circ\det$.  That list has no earlier
segment preceding a later one, so the standard-cosocle
theorem gives its Langlands multisegment.  Uniqueness
compares its $t$ segments with the single length-$t$
segment, forcing $t=1$.

Thus all Langlands segments of $\rho$ are unramified
singletons.  BCS15 on their standard induction and
BCS14 on its quotient identify their uniformizer values
with $S$.  This establishes the required singleton
Langlands description, including its ordering and
normalization, for the given $\rho$.

The Langlands and rank-one laws of the given
\texttt{Is\allowbreak Rec\allowbreak T} therefore give
\[
          \mathrm{rc}(\operatorname{cl}\rho)
                     =\bigoplus_{j=1}^a(\xi_j\circ a_{\mathrm{rc}},0).
\]
The degree convention in $\texttt{Is\allowbreak Local\allowbreak Artin}$
gives valuation $-1$ on arithmetic Frobenius, so the
unramified characters have values $\alpha_j^{-1}$
there.  Their inertia action is $1$.  This proves
the statement at the level of the actual WD class.
For rank zero the sum is empty.  In particular the
calculation is independent of any choices in
$\mathrm{rc}$, even though its full predicate is
stronger than the usual unramified correspondence.
\end{proof}
% END GUA12

% BEGIN GUA13
\paragraph{GUA13. Local Satake data on the original cusp subspace.}
\label{claim:GUA13}
For each specified cusp component $\rho_{i,w}$, its
local Satake multiset is also the multiset in the literal
$P_i.\texttt{Has\allowbreak Sat}(w,-)$, and conversely.  Such
multisets exist for all $i$ at a cofinite set of places.
They also exist for every $w$ satisfying the original
$\texttt{Is\allowbreak Unram\allowbreak U}(\sigma,w)$.
\begin{proof}
The restriction of an irreducible cusp representation to
one local factor is isotypic of any specified local
component.  Indeed the complementary adelic factors
and the infinity operators commute with this factor.
Translates of one nonzero local intertwiner are again
images of that same local representation, and their
sum is a nonzero subspace stable under every global
operator.  Irreducibility makes their sum the whole
cusp space.  Equivalently this is the elementary
isotypy argument of ITA10.

Let $K=\operatorname{GL}_{a_i}(\mathcal O_{E_w})$.
A local $K$-fixed eigenvector embeds into the global
cusp space through a nonzero intertwiner, which is
injective by local irreducibility.  Conversely, write
a global $K$-fixed vector as a finite sum of images
of local vectors, and average those vectors over $K$.
The average is a finite sum on each smooth vector;
the sum of the images is still the original nonzero
vector, so at least one local average is nonzero.
The same argument shows that every global fixed
vector has the local Hecke eigencharacter.  Uniqueness
of the spherical Satake multiset gives the asserted
equivalence, including the source requirement on all
nonzero fixed vectors.

Choose a nonzero smooth finite-adelic vector in each
of the finitely many $P_i$.  Its open stabilizer contains
the standard maximal compact at all but finitely many
places, by the restricted-product topology.  The
previous paragraph makes all the $\rho_{i,w}$
spherical there; the spherical Hecke classification
then supplies their Satake multisets.  For every
$\texttt{Is\allowbreak Unram\allowbreak U}(\sigma,w)$, BCS19 instead supplies spherical $\rho_{i,w}$
at that very place, so the same argument gives the
last assertion without discarding an extra finite set.
\end{proof}
% END GUA13

% BEGIN GUA14
\paragraph{GUA14. \texttt{Frob\allowbreak Tot\allowbreak T} and the complete unramified assertion.}
\label{claim:GUA14}
Let $w\nmid\ell$ and suppose the specified
$\rho_{i,w}$ have Satake multisets $S_{i,w}$.
Define, with multiplicities,
\[
 B_w=\coprod_i\coprod_{k=0}^{d_i-1}
       \{\alpha q_w^{k-(d_i-1)/2}:\alpha\in S_{i,w}\}.
                                                        \tag{GUA.19}
\]
For any continuous $R$ satisfying
$\texttt{Par\allowbreak At\allowbreak T}(\iota,R,w,p_w)$ with (GUA.11),
all inertia groups over $w$ act trivially, and for
every valuation $A$ over $w$ and every arithmetic
Frobenius $g$ at $A$,
\[
 \iota\det(X-R(g^{-1}))
   =\prod_{\beta\in B_w}(X-q_w^{(n-1)/2}\beta).
                                                        \tag{GUA.20}
\]
This is the entire \texttt{Frob\allowbreak Eig} assertion, hence
\texttt{Frob\allowbreak Tot\allowbreak T}.  The $R$ of GUA11 is consequently
unramified at \emph{every} $w\nmid\ell$ with
$\texttt{Is\allowbreak Unram\allowbreak U}(\sigma,w)$.
\begin{proof}
Fix any $A,h$ and use the retained \texttt{Rec\allowbreak TExists}
to choose one permissible $\mathrm{rc}$ there.  GUA12
and (GUA.11) make $p_{w,A,\mathrm{rc}}$ unramified,
with arithmetic eigenvalues
$\alpha^{-1}q_w^{e_{ik}/2}$ and zero monodromy.
It is good by GUA1.  Thus the conditional equality
in \texttt{Par\allowbreak At\allowbreak T} identifies it with the actual
Weil--Deligne parameter of $R$ at $A$.

This equality implies trivial \emph{original} inertia
action.  In the literal \texttt{Is\allowbreak WDOf} witness, the
zero monodromy operator makes $R$ trivial on an open
subgroup of inertia.  Compactness of inertia then
makes its image finite, so each inertia matrix is
semisimple in characteristic zero.  The original
\texttt{Is\allowbreak Frob\allowbreak SS} witness commutes with that matrix
and differs from it by a nilpotent matrix; since
both commuting matrices are semisimple, simultaneous
diagonalization makes their nilpotent difference zero.
Thus the original inertia matrix equals the inertia
matrix of the unramified Weil--Deligne parameter,
which is $1$.  For general Weil elements the same
commuting nilpotent difference preserves the
characteristic polynomial.
Inverting the arithmetic eigenvalues gives
\[
 \alpha q_w^{-e_{ik}/2}
    =q_w^{(n-1)/2}\alpha q_w^{k-(d_i-1)/2},
\]
which proves (GUA.20).  Trivial inertia makes it true
for every arithmetic Frobenius lift, and $A,h$ were
arbitrary.  The proof uses only the hypotheses of
the literal \texttt{Frob\allowbreak Tot\allowbreak T}, so it is not confined
to the representation constructed in GUA11.
For that representation use GUA13 at every
$\texttt{Is\allowbreak Unram\allowbreak U}$ place.  No cofinite exception
is introduced in this last conclusion.
\end{proof}
% END GUA14

% BEGIN GUA15
\paragraph{GUA15. Every split local component and its conditional jumps.}
\label{claim:GUA15}
For every split $w$ and every original local component
$\tau$ of $\sigma$ along the specified map
$\texttt{c\allowbreak Loc}\circ\texttt{loc\allowbreak U}$, the representation
of GUA11 satisfies
\[
       \texttt{Loc\allowbreak Glob\allowbreak T}(\iota,R,w,\tau),\qquad
       \texttt{Plain\allowbreak Tor}(\operatorname{cl}\tau)
                        \Longrightarrow\texttt{Has\allowbreak HTAt}(R,H,w).
                                                        \tag{GUA.21}
\]
\begin{proof}
At a split place the local unitary group is the
specified copy of $\operatorname{GL}_n(E_w)$; in
the paired coordinate the unitary equation determines
the other matrix as the conjugate inverse transpose.
The adelic group is a product of this factor and
its centralizer.  If an irreducible global module
contains one local irreducible image $V$, its global
translates are again images of $V$: remove the local
factor of the translating element, and the remaining
factor commutes with the local group.  The images
span the whole irreducible module.  ITA10, or a
coordinate projection from a finite sum of these
simple images, now shows that any second nonzero
local irreducible image is isomorphic to $V$.
This proves local-component uniqueness on the actual
function subspace, with no tensor-product realization
of that subspace assumed.

BCS supplies the full local normalized product
$\mathop{\times}_i\eta_{i,w}$ as a local component
of $\sigma$.  It is irreducible by the retained
unitary-product theorem (each factor is unitary),
or directly by the irreducibility field of its
$\texttt{Is\allowbreak Comp\allowbreak Along}$ witness.  The uniqueness just
proved identifies the given $\tau$ with that exact
product.  GUA6 identifies (GUA.11) with
$\mathrm{rc}(\operatorname{algNorm}\tau)$ for
every $A,h,\mathrm{rc}$, so GUA11's full
\texttt{Par\allowbreak At\allowbreak T} is exactly \texttt{Loc\allowbreak Glob\allowbreak T}.
If $\tau$ is \texttt{Plain\allowbreak Tor}, GUA2 makes this
parameter good for every one of those triples.
The \texttt{ht} field of GUA11's \texttt{Compat\allowbreak G}
therefore gives \texttt{Has\allowbreak HTAt}, with all its
original coefficient-ring, bridge and embedding quantifiers.
\end{proof}
% END GUA15

\subsection{Polarization and the irreducible case}

% BEGIN GUA16
\paragraph{GUA16. The required pairing for every complex conjugation.}
\label{claim:GUA16}
For every $c\in G_{E^+}$ satisfying the source's
\texttt{Is\allowbreak Conj\allowbreak El}, the $R$ of GUA11 admits an
invertible matrix $P$ with
\[
       R(s)^t P R(csc^{-1})=\epsilon(s)^{1-n}P
                                      \qquad(s\in G_E). 
                                                        \tag{GUA.22}
\]
Equivalently this is the source's \texttt{Is\allowbreak Pairing}
for every pair $s,s'$ whose restrictions satisfy
$s'=csc^{-1}$ in $G_{E^+}$.
\begin{proof}
At a cofinite set of places, all the cusp Satake
multisets in GUA13 exist at both $w$ and $cw$, and
both places are away from $\ell$.  Complex
conjugation preserves their residue norms.  The
conjugate-self-duality intertwiner of $P_i$ takes
its local spherical Hecke character at $cw$ to the
inverse Satake character at $w$.  To check the
normalization, inverse transpose sends the diagonal
torus character to its inverse; the half-modulus
normalization of the Satake transform is invariant
under reversing the roots.  Thus
\[
                 S_{i,cw}=\{\alpha^{-1}:\alpha\in S_{i,w}\}.
\]
Reverse $k$ to $d_i-1-k$ in (GUA.19).  Its exponent
changes sign, so $B_{cw}=B_w^{-1}$, including every
Speh length.  Let $C_w=q_w^{(n-1)/2}$.
By GUA14, $R^c$ has geometric Frobenius eigenvalues
$C_wB_{cw}$ at $w$.  On the same lift
$R^\vee\otimes\epsilon^{1-n}$ has eigenvalues
\[
         q_w^{n-1}(C_wB_w)^{-1}
                  =C_wB_{cw}.                         \tag{GUA.23}
\]
Here $\epsilon$ on geometric Frobenius is $q_w^{-1}$,
so the exponent $1-n$ in (GUA.23) is forced.

Both representations are continuous and semisimple.
Conjugating the group preserves these properties;
duality preserves semisimplicity by annihilators
of submodules, and a character twist preserves
the stable-subspace lattice.  Chebotarev makes
the union of the Frobenius conjugacy classes
outside the finite exceptional set dense in $G_E$.
Their characteristic polynomials agree by (GUA.23).
Each coefficient is a continuous function into
the Hausdorff field $\overline{\mathbb Q}_\ell$,
so equality holds on all of $G_E$.  The
Brauer--Nesbitt argument then gives an isomorphism
of the two semisimple representations: in the
finite-dimensional algebra spanned by their
combined image, simple-module characters separate
the semisimple factors, and equality of traces
gives their multiplicities.  This needs no purity
statement or bound on Satake absolute values.

Choose the matrix $P$ of an intertwiner
$R^c\longrightarrow R^\vee\epsilon^{1-n}$.
Its identity is
$P R(csc^{-1})=\epsilon(s)^{1-n}R(s)^{-t}P$;
multiplication by $R(s)^t$ gives (GUA.22).
Restriction of automorphisms is injective here,
so the formulation with arbitrary $s'$ is the
same identity.  The construction applies to each
specified complex conjugation $c$.  No symmetric
or alternating sign for $P$ is required by this endpoint.
\end{proof}
% END GUA16

% BEGIN GUA17
\paragraph{GUA17. Irreducibility produces the actual cuspidal representation.}
\label{claim:GUA17}
If $R$ is irreducible, there is an original cusp
$\Pi$ of $\operatorname{GL}_n(\mathbb A_E)$ that is
regular algebraic and satisfies the source's
\texttt{Dual\allowbreak Is\allowbreak Conj} and \texttt{Att\allowbreak C}.  It has
all the infinity weights of $\sigma$, has a
Satake multiset at every
$\texttt{Is\allowbreak Unram\allowbreak U}(\sigma,w)$, and every one of
its local components at a split $w$ is a component
of $\sigma$ along the original local map.
\begin{proof}
GUA11 gives $c=1$, $d_0=1$ and consequently $a_0=n$.
Take the actual BCS cusp $\Pi=P_0$; no automorphy
of a new representation is inferred.  BCS's
infinity equality now reads
$m_{0,\phi}=\operatorname{infOfWeight}(\lambda)$
for every $\texttt{Has\allowbreak Weight\allowbreak At}$ witness.  Its
\texttt{Has\allowbreak Inf\allowbreak Char} field gives the desired exact
infinity character.  At each embedding, GUA9
supplies such a weight; its entries are distinct
and lie in $\mathbb Z+(n-1)/2$.  This proves
regular algebraicity, also for $n=1$.

The given conjugate-self-duality of $P_0$, combined
with the retained Gelfand--Kazhdan duality and
evaluation pairing as in GUA7, gives
\texttt{Dual\allowbreak Is\allowbreak Conj}, including the finite, compact
and Lie identities in its definition.  At a
cofinite set of places GUA13 supplies its global
Satake multiset $S_{0,w}$.  For length one (GUA.19)
is exactly $S_{0,w}$, and GUA14 supplies
$\texttt{Frob\allowbreak Eig}(\iota,R,w,
\operatorname{algSat}(n,q_w,S_{0,w}))$.
Together with semisimplicity this is precisely
\texttt{Att\allowbreak C}, not just equality of unspecified
Euler factors.  The every-unramified-place
assertion follows directly from the last clause
of GUA13, without imposing $w\nmid\ell$.

For length one the defining Speh induction is
$\rho_{0,w}$ itself, so its irreducible quotient
$\eta_{0,w}$ is isomorphic to $\rho_{0,w}$.
The one-factor normalized product is the same
representation.  The BCS split component map
therefore makes $\rho_{0,w}$ a component of
$\sigma$.  Any other local component of $P_0$
is isomorphic to it: apply the isotypy argument
of GUA13 to the cusp space.  Transporting the
actual nonzero local intertwiner along this
isomorphism proves the requested component
assertion for \emph{every} local component of $\Pi$.
\end{proof}
% END GUA17

% BEGIN GUA18
\paragraph{GUA18. The complete \texttt{Gal\allowbreak Unitary\allowbreak T} and \texttt{Gu\allowbreak Wrap} endpoint.}
\label{claim:GUA18}
Under the original hypotheses stated at the beginning,
BCS and the retained original \texttt{Row\allowbreak T} supply
one pair $(R,H)$ with every field of the literal
\texttt{Gal\allowbreak Unitary\allowbreak T}.  With the original
$\texttt{Hypotheses\allowbreak H}'$ literature record, this is
the original \texttt{Gu\allowbreak Wrap.\allowbreak gal\allowbreak Unitary\allowbreak T\_\allowbreak of\_\allowbreak prints\allowbreak H}.
\begin{proof}
Choose the actual witnesses of BCS18--BCS19 and construct
$R,H$ as in GUA11.  The following table identifies
each conjunct, including its quantifiers; each
proof concerns this same $R,H$.
\begin{center}
\begin{tabular}{p{.69\linewidth}p{.19\linewidth}}
\hline
Original conclusion & Proof \\
\hline
Continuity and semisimplicity on $G_E$ & GUA11 \\
Every $w\nmid\ell$ with $\texttt{Is\allowbreak Unram\allowbreak U}$,
every valuation above it, and every inertia element:
$R(s)=1$ & GUA13--14 \\
Every $\tau$ and every infinity-weight witness:
the exact multiset in (GUA.18) & GUA9, GUA11 \\
Every split place and every specified local component:
\texttt{Loc\allowbreak Glob\allowbreak T}, with every $A,h,\mathrm{rc},O,B_w$
and \texttt{Induces}; its \texttt{Plain\allowbreak Tor}
implication to \texttt{Has\allowbreak HTAt} & GUA15 \\
Every original complex conjugation:
\texttt{Is\allowbreak Pairing} with $\epsilon^{1-n}$ & GUA16 \\
If $R$ is irreducible: an actual regular algebraic
conjugate-self-dual cusp, \texttt{Att\allowbreak C}, all infinity
weights, all $\texttt{Is\allowbreak Unram\allowbreak U}$ Satake places and
every split local component & GUA17 \\
\hline
\end{tabular}
\end{center}
The hypotheses of BCS include exactly finite-place
unramifiedness of $E/E^+$ and the original
\texttt{Is\allowbreak Aut\allowbreak Rep\allowbreak U}; its discrete base-change input
is the retained $\texttt{lit\allowbreak N.\allowbreak aut.\allowbreak bc}$.  The
polarized attachment is the original
$\texttt{lit\allowbreak N.\allowbreak aut.\allowbreak gal}$, whose source wrapper
$\texttt{row\allowbreak Of}'$ is the identity on \texttt{Row\allowbreak T}.
GUA7 used its unrestricted toral branch.  All
local arguments used only the retained classical
foundations and the checked BCS local unitary
input, and kept the full \texttt{Is\allowbreak Rec\allowbreak T}
predicate.  Hence substitution of these original
inputs gives the claimed \texttt{Gu\allowbreak Wrap} theorem.
No step removes its foundational hypotheses, changes
$\ell\ne2$ to a larger-prime condition, or imposes
$d_i=1$ without the stated irreducibility premise.
\end{proof}
% END GUA18

\subsection{Scope and source boundary}

The new proofs establish the local parameter, normalization,
period-sum, block-construction, pairing and cuspidal-consumer
arguments.  The substantial inherited attachment theorem is
the full \texttt{Row\allowbreak T}/\texttt{Props\allowbreak B} record; it has not
been replaced by a theorem-name-only assertion about Galois
representations.  Likewise BCS retains the original
\texttt{Base\allowbreak Change\allowbreak Isobaric\allowbreak N} foundation.  The existence of
full \texttt{Is\allowbreak Rec\allowbreak T} correspondences and the actual tied
period model remain explicit local foundations.  The
source-pinned metadata records their precise statements
separately from the numbered-claim dependency graph.
This chapter does not assert the independent
\texttt{end\allowbreak Shape\allowbreak T} strict-norm-ratio endpoint or purity;
those conclusions have different hypotheses and consumers.


\section{The actual automorphic parts and density on the original ideal}
\label{sec:automorphic-installation-v44}

This chapter completes the passage from a nonzero coefficient-valued
automorphic subspace to the full arithmetic part used by the ordinary
Hecke argument.  The representation, its quotient descent, its
polarization, all local conditions, and its spherical polynomial
readings are constructed together.  The last two statements install
these parts in the original \texttt{Automorphic\allowbreak Hecke\allowbreak Set} and remove
the additional supplier retained in AFS22.

Throughout, the source data are fixed.  Write $F=L^+$ for the maximal
real subfield of the CM field $L$, let $\Omega$ be its specified
algebraic closure, and let $\ell$, $s$, $V$, $Q$, $\mathcal O$,
$\bar r$, and $\iota:\overline{\mathbb Q}_\ell\simeq\mathbb C$
have their original meanings.  Put $n\geq1$, and write
\[
 p=V.\mathrm{pr}:G_L\longrightarrow G_0,\qquad
 G_L=\operatorname{Gal}(\Omega/L),\qquad
 G_F=\operatorname{Gal}(\Omega/F).
\]
The original finite set of places of $F$ is $S'=V.S'$.
The chosen integral type is $\tau_0$ and satisfies the literal
\texttt{Type\allowbreak For\allowbreak Alt} condition.  All original Teichm\"uller families,
all correspondence maps satisfying \texttt{Is\allowbreak Rec\allowbreak T}, and all coefficient
and period bridges remain quantified as in the source.  The fixed
embedding $\iota$ is used algebraically; its continuity is never assumed.

The inherited hypotheses are the source's
\texttt{Hypotheses\allowbreak H'} and its stated literature and period foundations.
The local-transport claims LTP5--LTP21 and Galois-assembly claim GUA18 supply their indicated
consumers.  No additional \texttt{Part\allowbreak Alt\allowbreak Alt},
\texttt{Automorphic\allowbreak Hecke\allowbreak Set}, \texttt{At\allowbreak LT}, or density assumption is
introduced here.  In particular, the original base-change and
cuspidal Galois-attachment foundations remain identified at the scope
recorded in those chapters.

\subsection{The original quotient and its pairing}

% BEGIN UPA1
\paragraph{UPA1. The place and ramification consequences of the Hida setting.}
\label{claim:UPA1}
The prescribed extension $L/F$ is unramified at every finite place.
For every chosen place $s.w(v)$ one has
$\underline{s.w(v)}=v$, where an underline means restriction to $F$.
Moreover
\[
 \ell\in w\quad\Longrightarrow\quad \underline w\in S',
 \qquad
 \underline w\in S'\quad\Longrightarrow\quad w\text{ is split over }F.
 \tag{UPA.1}
\]
These conclusions apply simultaneously at the original $\ell$-adic,
type, auxiliary, and distinguished places.
\begin{proof}
The condition $s.\mathrm{unr}$ says that every absolute inertia subgroup
over a finite place of $F$ is contained in the image of
$G_L\hookrightarrow G_F$.  The image of absolute inertia in
$\operatorname{Gal}(L/F)$ is the finite inertia subgroup.  It is
therefore trivial.  Residue fields of number fields are perfect, so
the usual equality of the inertia order with the ramification index
gives ramification index one at every finite place.  This is exactly
the finite-place unramifiedness used by the Galois-assembly theorem.

The valuation ring $s.A(v)$ lies over both $v$ and $s.w(v)$.
Pulling its maximal ideal back through
$\mathcal O_F\longrightarrow\mathcal O_L$ identifies the restriction
of $s.w(v)$ with $v$.  The same pullback identity gives
$\ell\in w$ if and only if $\ell\in\underline w$.
All places of $F$ over $\ell$ belong to $s.S$, and $s.S\subset S'$.
This proves the first implication in (UPA.1).

For $v\in S'$, the original level condition puts the chosen
decomposition group at $v$ in the image of $G_L$.
The image of that group in the quadratic Galois group
$\operatorname{Gal}(L/F)$ is consequently trivial.  All prolongations
over $v$ have conjugate decomposition groups, and the image of $G_L$
is normal.  The decomposition group in the quadratic extension is
therefore trivial at every place above $v$.  Its orbit on the places
above $v$ has two elements; this is the literal \texttt{Is\allowbreak Split\allowbreak At}
condition.  The definitions of $s.S$ and $Q$ place the original type
places, the distinguished place, and the auxiliary places in $S'$.
The auxiliary places are away from $\ell$ by the already checked
original \texttt{ax\allowbreak U\_\allowbreak s} conditions.
\end{proof}
% END UPA1

% BEGIN UPA2
\paragraph{UPA2. Inertia outside the level kills the actual quotient kernel.}
\label{claim:UPA2}
Let $K\leq G_L$ be a closed subgroup.  Suppose that it contains every
inertia subgroup at every prolongation of every place $w$ of $L$
with $\underline w\notin S'$.  Then
\[
                         \ker p\subset K.                    \tag{UPA.2}
\]
\begin{proof}
First record the valuation-theoretic point used in passing between
the two number fields.  If a valuation ring $A\subset\Omega$ lies
over a finite place $v$ of $F$, then every element of $\mathcal O_L$
is integral over $A$ and hence lies in $A$.  The inverse image of
the maximal ideal of $A$ in $\mathcal O_L$ is a prime $w$ lying
over $v$.  It is nonzero because its contraction is $v\ne0$.
Thus it is a finite place of $L$, and $A$ lies over $w$.

Write $H$ for the image of $G_L$ in $G_F$, and $N_{S'}$ for the
closed normal subgroup generated by the absolute inertia groups
at the places of $F$ outside $S'$.  This is the original subgroup
used to define the level quotient.  The image of $K$ in $G_F$
is closed: $K$ is compact, restriction is continuous, and $G_F$
is Hausdorff.

Consider a generator $g$ from inertia at $A$ over
$v\notin S'$ and any conjugate $bgb^{-1}$ in $G_F$.
It belongs to inertia at the transported valuation ring
$A\circ b^{-1}$ over the same place $v$.  The condition
$s.\mathrm{unr}$ places this inertia group in $H$.
Choose the place $w$ of $L$ under $A\circ b^{-1}$.
For an automorphism fixing $L$, inertia over $F$ at that
valuation ring is the same inertia condition as inertia over $L$.
The lift of $bgb^{-1}$ to $G_L$ therefore belongs to $K$,
because $\underline w=v\notin S'$.
Every conjugate of every generator of $N_{S'}$ belongs to
the image of $K$.  Closedness now gives
$N_{S'}\subset\operatorname{res}(K)$.

If $t\in\ker p$, its restriction to $F$ belongs to $N_{S'}$,
hence equals the restriction of some $k\in K$.
Restriction $G_L\hookrightarrow G_F$ is injective on the
underlying automorphisms of $\Omega$, so $t=k\in K$.
This proves (UPA.2) for the original quotient, without replacing
its defining normal closure.
\end{proof}
% END UPA2

% BEGIN UPA3
\paragraph{UPA3. Continuous descent preserves the same polarization matrix.}
\label{claim:UPA3}
Let $R:G_L\to\operatorname{GL}_n(\overline{\mathbb Q}_\ell)$
be continuous and trivial on inertia outside $S'$.
There is a unique continuous representation
$\rho:G_0\to\operatorname{GL}_n(\overline{\mathbb Q}_\ell)$
with $\rho\circ p=R$.
For every matrix $P$ giving the original pairing
\[
 R(g)^{\mathsf t}P R(cgc^{-1})
                =\varepsilon_\ell(g)^{\,1-n}P,
 \tag{UPA.3}
\]
the same $P$ gives the original \texttt{Is\allowbreak Form} for $\rho$.
\begin{proof}
The map $p$ is surjective.  Indeed, lift an element of $G_0$
to an element of $G_F$ in the quotient by $N_{S'}$.
The level identity $V.h\eta$ identifies the inverse image of
$G_0$ with the image of $G_L$, so that lift comes from $G_L$.
This proves surjectivity for the actual map $p$.

The kernel of $R$ is closed.  It is the finite intersection
of the closed conditions $R(g)_{ij}=\delta_{ij}$, since
$\overline{\mathbb Q}_\ell$ is Hausdorff and each matrix entry
is continuous.  UPA2 gives $\ker p\subset\ker R$.
Define $\rho(p(g))=R(g)$.  The kernel inclusion proves
well-definedness, and surjectivity proves uniqueness and the
homomorphism law.

The group $G_L$ is compact.  The target $G_0$ is Hausdorff
because $N_{S'}$ is closed, and $p$ is continuous by the
original restriction and quotient construction.  Thus $p$
is a closed surjection and a quotient map.  The identity
$\rho\circ p=R$ proves continuity of $\rho$.

The frame's involution and cyclotomic character satisfy,
by their original defining identities,
\[
 \operatorname{conj}(p(g))=p(cgc^{-1}),\qquad
 \varepsilon_{G_0}(p(g))=\varepsilon_\ell(g).
 \tag{UPA.4}
\]
Here $cgc^{-1}$ denotes the unique automorphism fixing $L$
whose restriction to $F$ is the displayed conjugate; it
exists because $L/F$ is normal.  In coordinates over
$\mathcal O$, the second identity is precisely the
$V.h\varepsilon$ identity followed by the original scalar
map to $\overline{\mathbb Q}_\ell$.
Substitution of (UPA.4) and $\rho(p(g))=R(g)$ into
(UPA.3), followed by surjectivity of $p$, gives
\[
 \rho(h)^{\mathsf t}P\rho(\operatorname{conj}h)
                   =\varepsilon_{G_0}(h)^{\,1-n}P
 \quad(h\in G_0).
 \tag{UPA.5}
\]
Invertibility of $P$ is unchanged.  This is the literal
source \texttt{Is\allowbreak Form}, with no new pairing or scalar extension.
\end{proof}
% END UPA3

% BEGIN UPA4
\paragraph{UPA4. Semisimplicity and irreducibility pass through the quotient.}
\label{claim:UPA4}
For a surjection $p:G_L\to G_0$ and $R=\rho\circ p$,
$R$ and $\rho$ have exactly the same invariant subspaces.
Consequently $R$ is semisimple if and only if $\rho$ is
semisimple, and $R$ is irreducible if and only if $\rho$
is irreducible.
\begin{proof}
A $\rho$-invariant subspace is $R$-invariant by composition.
Conversely, for $h\in G_0$ choose $g$ with $p(g)=h$;
$R$-invariance then gives invariance under $\rho(h)$.
The same argument applies to any proposed invariant
complement.  The definitions of semisimplicity and
irreducibility now give the two equivalences.  The
underlying vector space and all of its subspaces are
unchanged.
\end{proof}
% END UPA4

\subsection{Spherical polynomials with the original Frobenius convention}

% BEGIN UPA5
\paragraph{UPA5. The spherical Satake theorem has the required normalization.}
\label{claim:UPA5}
For $k=L_w$, $K=\operatorname{GL}_n(\mathcal O_k)$,
and a smooth irreducible complex representation $\pi$ with
$\pi^K\ne0$, there is a multiset
$B=\{\beta_1,\ldots,\beta_n\}\subset\mathbb C^\times$
such that
\[
 {\bf1}_{K\operatorname{diag}(\varpi I_j,I_{n-j})K}
       \ \text{acts on }\pi^K\text{ by }\
 q^{j(n-j)/2}e_j(B)\quad(0\leq j\leq n).
 \tag{UPA.6}
\]
Here $\operatorname{vol}(K)=1$ and $q=\#(\mathcal O_k/\mathfrak m_k)$.
\begin{proof}[Published theorem and exact application]
Apply the spherical Satake isomorphism to the split reductive group
$\operatorname{GL}_n/k$.  A precise reference is B.~H.~Gross,
\emph{On the Satake isomorphism}, Proposition~3.6,
formula~(3.14), and Proposition~6.4.\footnote{
\url{https://people.math.harvard.edu/~gross/preprints/sat.pdf},
pp.~6, 8, 11; the Haar normalization is (2.5).}
The hyperspecial group is exactly $K$ and the Haar normalization
is the one in (UPA.6).  The cocharacter $(1^j,0^{n-j})$
is minuscule; its dual representation is $\bigwedge^j\mathbb C^n$,
whose character is $e_j$.  Its pairing with the half-sum of
positive roots is $j(n-j)/2$.  Thus the stated formula is
the normalized Satake formula.  The parameter lies in
$\operatorname{GL}_n(\mathbb C)$, so all $\beta_i$ are nonzero.
The result holds for every finite extension of $\mathbb Q_p$
and every $n\geq1$; the source places meet these hypotheses.
The fixed space is nonzero and one dimensional, the multiset has
cardinality $n$, and (UPA.6) holds on every fixed vector.
Changing the uniformizer by a unit leaves the double cosets unchanged.
Thus $B$ satisfies the literal source $\mathrm{HasSatGL}$ predicate,
with all its uniformizer and fixed-vector quantifiers; it is exactly
the normalized parameter used in GUA12.
\end{proof}
% END UPA5

% BEGIN UPA6
\paragraph{UPA6. The original spherical operators give the original Hecke polynomial.}
\label{claim:UPA6}
For the representation and multiset in UPA5, put
\[
 a_j=q^{h(n-h)/2}e_h(B),\qquad h=\min(j,n).
 \tag{UPA.7}
\]
Then $a_0=1$, every original operator $\operatorname{sphOp}_j$
acts on $\pi^K$ by $a_j$, and
\[
 \prod_{i=1}^n\bigl(X-q^{(n-1)/2}\beta_i\bigr)
 =\sum_{j=0}^n(-1)^j q^{j(j-1)/2}a_jX^{n-j}
 =\operatorname{heckePol}(n,q,a).
 \tag{UPA.8}
\]
\begin{proof}
The source's compact subgroup $K_O(k,n)$ equals $K$:
both consist of invertible integral matrices with integral inverse.
Its chosen uniformizer has valuation one for the normalized
discrete valuation.  Any two such choices differ by a unit,
and therefore determine the same diagonal double coset.

The source matrix $\operatorname{al}(k,n,j)$ has a uniformizer
in its first $\min(j,n)$ diagonal entries and one elsewhere.
For a $K$-fixed vector, its double-coset operator is the
sum of $\pi(x)$ over the right cosets $xK$ in that double
coset.  This is exactly convolution by its characteristic
function when $\operatorname{vol}(K)=1$; compactness makes
the sum finite.  UPA5 proves (UPA.7) for all $j$, including
$j>n$.  At $j=0$ the operator and its eigenvalue are one.

Expand the product on the left of (UPA.8) by elementary
symmetric functions.  Its coefficient of $X^{n-j}$ is
$(-1)^j q^{j(n-1)/2}e_j(B)$.  Since
\[
       \frac{j(n-1)}2
          =\frac{j(j-1)}2+\frac{j(n-j)}2,
\]
this is the coefficient on the right.  The number
$j(j-1)/2$ is an integer, and all real powers here use the
positive real $q$.  Thus (UPA.8) is an equality with the
literal source polynomial and its original normalizations.
\end{proof}
% END UPA6

% BEGIN UPA7
\paragraph{UPA7. Spherical local compatibility gives the geometric Frobenius polynomial.}
\label{claim:UPA7}
Let $w\nmid\ell$, let $\pi$ be spherical with Satake multiset
$B$, and let $R$ be a continuous global representation.
Suppose that, for every prolongation $A$ of $w$ and every
$\mathrm{rc}$ satisfying \texttt{Is\allowbreak Rec\allowbreak T},
\[
 \texttt{WDT}\bigl(\iota,R,\mathrm{rc}(\operatorname{algNorm}\pi)\bigr)
 \tag{UPA.9}
\]
holds.  Then $R$ is unramified at $w$ and, for every
arithmetic Frobenius lift $g$ at $A$,
\[
 \iota\!\left(\operatorname{charpoly}R(g^{-1})\right)
       =\prod_{\beta\in B}(X-q^{(n-1)/2}\beta)
       =\operatorname{heckePol}(n,q,a).
 \tag{UPA.10}
\]
\begin{proof}
Choose a correspondence using the retained nonemptiness theorem
\texttt{Rec\allowbreak TExists}.  The spherical correspondence calculation
GUA12, with the geometric reciprocity
normalization, gives
\[
 \mathrm{rc}(\operatorname{algNorm}\pi)
   =\left(\bigoplus_{\beta\in B}\operatorname{unr}
       \bigl(\beta^{-1}q^{(1-n)/2}\bigr),\,0\right).
 \tag{UPA.11}
\]
This calculation follows also by ordering the unramified
rank-one characters, taking their Langlands quotient, and
applying the \texttt{one}, \texttt{langlands}, and
\texttt{nu} clauses of \texttt{Is\allowbreak Rec\allowbreak T}; no global uniqueness
assertion for the weaker correspondence is used.

The parameter in (UPA.11) has zero monodromy and is trivial
on commutators of the Weil group, because every summand is
a character.  It therefore satisfies the source predicate
\texttt{Gd}.  Unpacking \texttt{WDT}, its witnessing
Weil--Deligne parameter must equal (UPA.11).
The zero-monodromy argument of UPA9, applied with trivial inertia
characters, makes the original inertia action trivial: an open
inertia subgroup is killed, every inertia matrix consequently has
finite order, and the commuting semisimple-plus-nilpotent comparison
then identifies it with the identity.  For every Weil element, the
same commuting semisimple-plus-nilpotent decomposition preserves
its characteristic polynomial.  This proves both unramifiedness
and the required polynomial comparison.  Arithmetic Frobenius
has degree one.  Its inverse has degree minus one, so its
eigenvalues in (UPA.11) are
\[
       \bigl(\beta^{-1}q^{(1-n)/2}\bigr)^{-1}
                          =q^{(n-1)/2}\beta.
\]
This proves the first equality of (UPA.10); UPA6 proves
the second.  In particular the inverse on $g$ is part of
the conclusion and has not been absorbed by renaming the
Satake parameters.
\end{proof}
% END UPA7

\subsection{Type and ordinary conditions for the same representation}

% BEGIN UPA8
\paragraph{UPA8. A global-character inertia type activates the exact WD comparison.}
\label{claim:UPA8}
Suppose a Weil--Deligne class $y$ has zero monodromy and
diagonal inertia characters that are restrictions of
characters $\chi_i:G_L\to\mathbb C^\times$.
Then $y$ satisfies \texttt{Gd}.  Consequently
\texttt{WDT}$(\iota,R,y)$ implies
\texttt{Is\allowbreak WDOf}$(\iota,R,y)$.
\begin{proof}
Choose the representative occurring in the literal
\texttt{Is\allowbreak Type\allowbreak WD} condition.  Every Weil commutator has
degree zero and hence belongs to inertia.  Each $\chi_i$
is trivial on it, because its target is commutative.
The representative therefore acts as the identity on
every commutator.  Its characteristic polynomial there
is $(X-1)^n$, which is precisely the source torality
condition \texttt{Tor}.  Its monodromy is zero by
hypothesis.  These two assertions are \texttt{Gd}.
The definition of \texttt{WDT} provides a parameter $x$
with \texttt{Is\allowbreak WDOf}$(\iota,R,x)$ and an implication
$\texttt{Gd}(y)\Rightarrow x=y$; substitution gives the
last assertion.
\end{proof}
% END UPA8

% BEGIN UPA9
\paragraph{UPA9. Zero-monodromy type gives literal diagonal inertia over the original field.}
\label{claim:UPA9}
Let $w\nmid\ell$, let $A$ lie over $w$, and suppose
\texttt{Is\allowbreak WDOf}$(\iota,R,y)$.
If $y$ has zero monodromy and inertia type
$\bigoplus_{i=0}^{n-1}\iota\chi_i$, with
$\chi_i:I_A\to\overline{\mathbb Q}_\ell^\times$,
then there is $C\in\operatorname{GL}_n(\overline{\mathbb Q}_\ell)$
such that, simultaneously for all $t\in I_A$,
\[
               C R(t) C^{-1}=\operatorname{diag}(\chi_i(t)).
 \tag{UPA.12}
\]
\begin{proof}
Unpack the source \texttt{Is\allowbreak WDOf} witness.  It contains a
monodromy matrix $N$ over $\overline{\mathbb Q}_\ell$, an
open inertia subgroup $U$ on which $R$ is the exponential
of a scalar multiple of $N$, and a Frobenius-semisimple
complex representative $(r_{\mathrm{WD}},N_{\mathrm{WD}})$.
The given type identifies this representative with one
having zero monodromy.  The monodromy intertwining
identity implies $N_{\mathrm{WD}}=0$; injectivity of the
matrix scalar map induced by $\iota$ gives $N=0$.
Thus $R|_U=1$.

For every $t\in I_A$, some positive power belongs to $U$.
Indeed, $U$ is open in the compact inertia group, hence
has finite index; two powers of $t$ have the same coset.
It follows that $R(t)$ has finite order.  Over a
characteristic-zero field its minimal polynomial divides
$X^m-1$ for some $m>0$, a separable polynomial.
Thus $\iota R(t)$ is semisimple.

The source Frobenius-semisimplification witness says that
$\iota R(t)$ commutes with $r_{\mathrm{WD}}(t)$ and that
their difference is nilpotent.  Here the usual monodromy
exponential is one because $N=0$.
The second operator is also semisimple.  Commuting
semisimple operators over $\mathbb C$ are simultaneously
diagonalizable, so their difference is semisimple.
A semisimple nilpotent operator is zero.  Hence
\[
                   \iota R(t)=r_{\mathrm{WD}}(t)
                   \quad(t\in I_A).                         \tag{UPA.13}
\]
The fixed intertwiner supplied by the type conjugates
the right-hand side of (UPA.13) to
$\operatorname{diag}(\iota\chi_i(t))$ for all $t$ at once.
Apply $\iota^{-1}$ entrywise to that intertwiner and its
inverse.  Their product identities remain valid and give
the matrix $C$ in (UPA.12).
This is algebraic transport of matrices; no topological
property of $\iota$ is required.  Repeated characters cause
no difficulty.
\end{proof}
% END UPA9

% BEGIN UPA10
\paragraph{UPA10. Distinct-unit principal series satisfy the exact plain-toral predicate.}
\label{claim:UPA10}
Let $\pi\simeq\operatorname{PS}(\omega)$ be smooth and
irreducible, where all $\omega_i:k^\times\to\mathbb C^\times$
are smooth and their restrictions to $\mathcal O_k^\times$
are pairwise distinct.  Then the irreducible class of
$\pi$ satisfies the source predicate \texttt{Plain\allowbreak Tor}.
\begin{proof}
By the definition of $\operatorname{PS}$, its normalized
induction characters are
\[
              \mu_i=\omega_i|\cdot|^{(n-1)/2}.
 \tag{UPA.14}
\]
The original cuspidal-support theorem identifies the
cuspidal support of the irreducible principal series with
the multiset of these rank-one characters.  Thus every
support member has rank one, the first clause of
\texttt{Plain\allowbreak Tor}.

No two support members differ by the source unit norm
twist.  For distinct indices, restrict a proposed equality
$\mu_i=\mu_j|\cdot|$ to units.  Both norm factors become
one, contradicting the distinct-unit hypothesis.
For the same index it would give $|\cdot|=1$, contradicted
at a uniformizer because $q>1$.  This is the second
clause of \texttt{Plain\allowbreak Tor}.  The support assertion uses
the original local classification foundation already
retained in the principal-series arguments; no assertion
about arbitrary nontoral principal series is needed.
\end{proof}
% END UPA10

% BEGIN UPA11
\paragraph{UPA11. The infinity multiset yields the original Hodge--Tate integers.}
\label{claim:UPA11}
Suppose the original multiset identity at a local embedding
$j:L_w\to\overline{\mathbb Q}_\ell$ is
\[
 \iota_{\mathbb Z}\bigl(\mathrm{HT}(j|_L)\bigr)
   =\left\{\,\frac{n-1}{2}-z:
                    z\in\operatorname{infOfWeight}(\lambda(j|_L)^\vee)
     \,\right\}.
 \tag{UPA.15}
\]
Then
\[
 \mathrm{HT}(j|_L)
       =\{\,\lambda(j|_L)_{n-1-i}+i:0\leq i<n\,\}.
 \tag{UPA.16}
\]
Every embedding arising from an original local period bridge
is a restriction $j|_L$ of this kind.
\begin{proof}
The source conventions give
$\lambda^\vee_i=-\lambda_{n-1-i}$ and
\[
 \operatorname{infOfWeight}(\lambda^\vee)
   =\left\{-\lambda_{n-1-i}+\frac{n-1}{2}-i:0\leq i<n\right\}.
 \tag{UPA.17}
\]
Subtracting each element from $(n-1)/2$ gives the integers
in (UPA.16), with their multiplicities.  The map
$\mathbb Z\hookrightarrow\mathbb C$ is injective, so its
map on finite multisets is injective.  This proves
(UPA.16).  The displayed integers are exactly the
source functions $\operatorname{htW}(\lambda,i,j|_L)$.

For the assertion about a period bridge, let $\tau:L\to
\overline{\mathbb Q}_\ell$ be induced by the bridge and
an embedding of its finite local coefficient field.
Embeddings between finite extensions of $\mathbb Q_\ell$
preserve the extended valuation: the absolute value is
the unique extension of the one on $\mathbb Q_\ell$,
equivalently the spectral norm determined by the minimal
polynomial.  The bridge's valuation identity therefore
shows that the finite place induced by $\tau$ is exactly
$w$.  The map $\tau$ extends continuously to the completion
$L_w$.  Its image lies in the finite complete extension
generated by $\tau(L)$ and $\mathbb Q_\ell$, so this
extension still takes values in $\overline{\mathbb Q}_\ell$.
It is the required $j$, and $j|_L=\tau$.
\end{proof}
% END UPA11

% BEGIN UPA12
\paragraph{UPA12. The original local-global and Hodge data imply the full AtLT condition.}
\label{claim:UPA12}
Suppose $R$, $\mathrm{HT}$, $w$, and $\pi$ satisfy
\texttt{Loc\allowbreak Glob\allowbreak T}; suppose
\texttt{Plain\allowbreak Tor}$(\pi)$ implies
\texttt{Has\allowbreak HTAt}$(R,\mathrm{HT},w)$; and suppose the
infinity-weight identity (UPA.15) holds at every local
embedding.  Then the original
\texttt{At\allowbreak LT}$(\iota,A,R|_{D_A},\lambda,\pi)$ holds.
\begin{proof}
Unpack an arbitrary test in the definition of
\texttt{At\allowbreak LT}.  It gives smooth characters $\omega_i$
with pairwise distinct unit restrictions and an
equivalence $\pi\simeq\operatorname{PS}(\omega)$.
By UPA10, $\pi$ is plain toral.  The stated Hodge input
therefore gives, for every original period bridge,
the de Rham assertion and its precise jump multiset.
UPA11 identifies that multiset with
$\{\operatorname{htW}(\lambda,i,\tau)\}_i$ for every
embedding induced by that bridge.

Fix also the test's arbitrary $\mathrm{rc}$ satisfying
\texttt{Is\allowbreak Rec\allowbreak T}.  The plain-toral parameter calculation in
GUA2 gives
\[
           \texttt{Gd}\bigl(\mathrm{rc}
                      (\operatorname{algNorm}\pi)\bigr).
 \tag{UPA.18}
\]
Applying the second clause of \texttt{Loc\allowbreak Glob\allowbreak T} to each
of the same period bridges now gives potential
semistability and equality of its Weil--Deligne class
with $\mathrm{rc}(\operatorname{algNorm}\pi)$.
In particular their characteristic polynomials agree
on every Weil element.

The de Rham assertion, the jump assertion, and this
characteristic-polynomial identity are exactly the
three clauses of the source \texttt{DRat} predicate.
They hold for every test and every bridge, so they
give \texttt{At\allowbreak LT}.  The additional \texttt{Ab\allowbreak W}
hypothesis in each test is retained and causes no
restriction in this argument.  No ordinary flag has
been silently substituted for the stated predicate.
\end{proof}
% END UPA12

\subsection{The full part and its point properties}

% BEGIN UPA13
\paragraph{UPA13. One representation and one family of local components suffice.}
\label{claim:UPA13}
Fix an original algebraic weight $\lambda$, its actual basis
$e$ and representation $\xi$, and an automorphic irreducible
subspace $\sigma$ for which
\[
                         Y_\sigma=YOf(Q,\iota,\tau_0,e,\xi,\sigma)
                                      \ne0.
 \tag{UPA.19}
\]
There are a continuous semisimple representation
$\rho:G_0\to\operatorname{GL}_n(\overline{\mathbb Q}_\ell)$,
one multiset function $\mathrm{HT}$, and a choice
$\pi_w$ of local component at every split place, such that:
the pullback $R=\rho\circ p$ has all the conclusions of
the literal \texttt{Gal\allowbreak Unitary\allowbreak T} theorem; every $\pi_w$
is a component of the same $\sigma$; and $\pi_w$ is spherical
whenever $\sigma$ is unramified at $w$.
\begin{proof}
Use the split-component existence and fixedness statement
from LTP19--LTP21, choosing $\pi_w$ for
each split $w$.  Its fixed-vector assertion applies to
these very components, so later unramifiedness does not
require a second component choice.

UPA1 verifies the finite-place unramifiedness hypothesis
for the CM extension.  The setting's original odd-prime
condition verifies the other prime hypothesis.  Apply
GUA18 to the original $\sigma$,
rank $n$, and fixed $\iota$.  It produces a single $R$
and $\mathrm{HT}$ with continuity, semisimplicity,
the stated inertia vanishing, the weight identity,
local compatibility at every split place and for every
component, the pairing, and the conditional cuspidal
conclusion.

LTP12 proves that $\sigma$ is
unramified at every $w$ with $\underline w\notin S'$.
By UPA1 these places are away from $\ell$.  Thus the
inertia-vanishing clause of GUA14 and GUA18
applies at every prolongation of each such $w$.
UPA3 descends $R$ to $\rho$ through the actual $p$,
and UPA4 gives semisimplicity of $\rho$.
All the stated local and weight identities are
identities for the same pullback $R=\rho\circ p$.
\end{proof}
% END UPA13

% BEGIN UPA14
\paragraph{UPA14. The sixteen fields of the original arithmetic part are filled.}
\label{claim:UPA14}
Under the data of UPA13 there exists an original
\texttt{Part\allowbreak Alt\allowbreak Alt} $Y$ with $Y.Y=Y_\sigma$.
Its representation is the $\rho$ of UPA13, and its local
components are the restrictions of the one chosen family
$\pi_w$ to the original $\ell$-adic and auxiliary places.
\begin{proof}
We specify every field.
Set $Y.Y=Y_\sigma$.  The definition of $YOf$ is an
intersection with the original \texttt{Big} space, so
$Y.Y\leq\texttt{Big}$.  Nonzeroness is (UPA.19).
At an $\ell$-adic index $v$, set
$Y.\sigma_l(v)=\pi_{s.w(v)}$; at an auxiliary index $q$,
set $Y.\sigma_q(q)=\pi_{Q.wq(q)}$.
UPA1 proves that these places are split, and
LTP7 and LTP9 give the exact
\texttt{iso\allowbreak L} and \texttt{iso\allowbreak Q} fields.

For each original spherical index $w\in H.Wp$, its
underlying place is split and lies outside $S'$.
The component $\pi_w$ chosen in UPA13 has nonzero
maximal-compact fixed vectors.  Apply UPA5 and UPA6
to obtain $B_w$ and the full sequence $a(w,j)$.
This gives $a(w,0)=1$.  LTP11
transfers the scalar action of every
$\operatorname{sphOp}_j$ from $\pi_w$ to every element
of $Y_\sigma$.  These are the \texttt{a},
\texttt{a\_\allowbreak zero}, and \texttt{sph} fields.

Set $Y.r=\rho$.  Its continuity and semisimplicity
are UPA13.  At each spherical index, the away-from-$\ell$
clause of \texttt{Loc\allowbreak Glob\allowbreak T} gives (UPA.9) for every
prolongation and every permitted correspondence.
UPA7 therefore gives, for every original Frobenius lift $g$,
\[
 \iota\!\left(\operatorname{charpoly}
                  \rho\bigl(p(g^{-1})\bigr)\right)
       =\operatorname{heckePol}(n,\operatorname{Norm}w,a(w,-)).
 \tag{UPA.20}
\]
This is \texttt{frob}, including its specified inverse.

At each $\ell$-adic place LTP14
identifies the infinity weight of $\sigma$
with the original dual of $\lambda$.  The weight clause
of GUA15 and GUA18 is therefore (UPA.15).
Its local clause gives \texttt{Loc\allowbreak Glob\allowbreak T} and the
plain-toral implication to \texttt{Has\allowbreak HTAt}.
UPA12 supplies the entire \texttt{at\allowbreak L} field for
the original local representation
$\rho\circ H.\mathrm{dec}(v)=R|_{D_{s.A(v)}}$.

Each auxiliary place is away from $\ell$ by the original
$Q$-conditions.  Taking the first clause of
\texttt{Loc\allowbreak Glob\allowbreak T} gives \texttt{WOT} for every
prolongation and every \texttt{Is\allowbreak Rec\allowbreak T} map, exactly
the \texttt{at\allowbreak Q} field.
We have supplied
\[
\begin{gathered}
Y,\ \mathrm{le},\ \mathrm{ne},\
\sigma_l,\ \mathrm{isoL},\
\sigma_q,\ \mathrm{isoQ},\
a,\ a_{\mathrm{zero}},\ \mathrm{sph},\\
r,\ \mathrm{cont},\ \mathrm{ss},\
\mathrm{frob},\ \mathrm{atL},\ \mathrm{atQ}.
\end{gathered}
\]
Every entry uses the original function space and the same $\rho$.
\end{proof}
% END UPA14

% BEGIN UPA15
\paragraph{UPA15. The same part satisfies all four unitary point properties.}
\label{claim:UPA15}
The part $Y$ of UPA14 satisfies the original
\texttt{Unitary\allowbreak Point\allowbreak Properties}$(\bar r,\iota,\lambda,Y.r)$:
its \texttt{form}, \texttt{type}, \texttt{cusp}, and
\texttt{aut} fields all hold with the original quantifiers.
\begin{proof}
For \texttt{form}, apply the pairing clause of
GUA16 to the setting's chosen complex
conjugation.  UPA3 transports its same invertible matrix
to the original frame, with multiplier
$\varepsilon_{G_0}^{1-n}$.

For \texttt{type}, fix an arbitrary Teichm\"uller family
$t$ for $\bar r$ and an arbitrary original type index $u$.
The condition \texttt{Type\allowbreak For\allowbreak Alt}, together with the
nonzero tensor-intertwiner constructed in
LTP17--LTP18, proves the literal
\texttt{Is\allowbreak Type\allowbreak WD} condition for $\pi_{s.w(u)}$ and
every allowed correspondence.  Choose one such
correspondence using \texttt{Rec\allowbreak TExists}.  The characters
on inertia are the restrictions of the actual global
characters
\[
 G_L\xrightarrow{p}G_0\xrightarrow{t_i}\mathcal O^\times
       \longrightarrow\overline{\mathbb Q}_\ell^\times.
 \tag{UPA.21}
\]
The type place is away from $\ell$.  Local compatibility
there gives \texttt{WDT}, and UPA8 upgrades it to
\texttt{Is\allowbreak WDOf}.  UPA9 supplies one matrix $C$ conjugating
the full inertia action to the ordered diagonal of (UPA.21).
The frame subgroup $I_u$ is the image of that actual
inertia group under $p$.  Taking preimages proves
\[
 C\rho(x)C^{-1}
       =\operatorname{diag}\bigl(\mathcal O\to
                       \overline{\mathbb Q}_\ell(t_i(x))\bigr)
       \quad(x\in I_u).
 \tag{UPA.22}
\]
The choices of $t$ and $u$ were arbitrary.  In particular
the assertion retains every Teichm\"uller family, not
only a family used earlier to construct $\tau_0$.

For \texttt{cusp}, suppose $\rho$ is irreducible.
UPA4 makes $R=\rho\circ p$ irreducible.
GUA17 gives a regular algebraic,
conjugate-self-dual cuspidal $\pi$ attached to $R$.
Its infinity-character clause, applied to the dual-weight
identities of LTP14, gives the original
$\operatorname{HasWt}(\iota,s,\pi,\lambda)$.
This proves precisely the conditional cusp assertion.

For \texttt{aut}, use the original $\sigma$ as witness.
It is automorphic by hypothesis.  Its required
all-component \texttt{Loc\allowbreak Glob\allowbreak T} is supplied by
GUA15 and GUA18 for the same $R$.  LTP12 and LTP14
give its dual infinity weights and unramifiedness
outside $S'$.  The family in UPA13 gives a local component
at every split place.  At $s.w(s.v_0)$ use that very
component; LTP15
gives its nonzero $K_c(1)$-fixed vector.
These are every clause in the source \texttt{aut} field.
\end{proof}
% END UPA15

% BEGIN UPA16
\paragraph{UPA16. The strengthened PartX statement and the original PartU follow.}
\label{claim:UPA16}
The constructions above prove the full original
\texttt{part\allowbreak X\_\allowbreak alt}, and hence \texttt{part\allowbreak U\_\allowbreak alt}, on
their original inputs.
In addition to UPA14--UPA15, the chosen $\ell$-adic components
are actual components of $\sigma$ and
\texttt{Loc\allowbreak Glob\allowbreak T} holds for every split place and every
component of $\sigma$.  If $Y.r\circ p$ is irreducible,
the resulting cusp is regular algebraic and conjugate
self-dual, is attached to $Y.r\circ p$, is spherical
wherever $\sigma$ is unramified, and has the required
split local-component agreement.
\begin{proof}
The two local-component assertions were retained
explicitly in UPA13 and were used unchanged in UPA14.
For the strengthened final implication, apply the
full conditional cuspidal clause of the same
GUA17--GUA18 to $R=Y.r\circ p$.
It gives all the listed assertions for one cusp,
including its Satake multiset at every place where
$\sigma$ is unramified and the implication from each
of its split local components to a component of
$\sigma$.  UPA15 supplies the four point properties,
and UPA14 supplies $Y.Y=Y_\sigma$.  These form the
entire tuple required by \texttt{part\allowbreak X\_\allowbreak alt}.
Projecting away its strengthened final implication
gives exactly \texttt{part\allowbreak U\_\allowbreak alt}.
\end{proof}
% END UPA16

\subsection{Installation and the original same-ideal density theorem}

% BEGIN UPA17
\paragraph{UPA17. The original AutomorphicHeckeSet is constructed.}
\label{claim:UPA17}
For every original $\tau_0$ satisfying \texttt{Type\allowbreak For\allowbreak Alt},
\[
 (hset(Q,\tau_0)).\operatorname{AutomorphicHeckeSet}
       \bigl(\iota,\operatorname{UnitaryPointProperties}
                              (\bar r,\iota)\bigr)
 \tag{UPA.23}
\]
holds under the original \texttt{Hypotheses\allowbreak H'}.
\begin{proof}
Fix an arbitrary admitted weight, an arbitrary algebraic
coefficient representation $\xi$, and a function $F$
meeting the source's \texttt{Big}, away-level, and
$I_1(e)$ fixedness conditions.  Choose the actual
algebraic-weight basis $e$ provided by
\texttt{Is\allowbreak Alg\allowbreak Wt}; the choices of level and depth
are kept as supplied.

ITA2 gives a Teichm\"uller family.  Apply
\texttt{Type\allowbreak For\allowbreak Alt} to that family to see that the
complexified type has open kernel.  The scalar map
$\mathcal O\to\overline{\mathbb Q}_\ell\xrightarrow{\iota}
\mathbb C$ is injective.  Entrywise injectivity on
general linear groups shows that the complexified
kernel is exactly $\ker\tau_0$.  Thus the integral
type has open kernel, as also recorded in ITA12.

ADP12 now places $F$ in the supremum of the original
$Y_\sigma$, over automorphic irreducible $\sigma$.
For each index with $Y_\sigma=0$, its contribution
is the zero subspace.  For every other index, UPA16
produces an original part $Y$ with $Y.Y=Y_\sigma$
and all the point properties.  Therefore each
$Y_\sigma$ is contained in the supremum over the
parts in (UPA.23), and so is their supremum.
This proves the required membership for $F$.
All initially fixed weight, basis, level, depth,
and function data were arbitrary.  This is the
literal \texttt{Au\allowbreak Teich\allowbreak Ex} endpoint.
\end{proof}
% END UPA17

% BEGIN UPA18
\paragraph{UPA18. The bare TJDense endpoint holds for the same original ideal.}
\label{claim:UPA18}
Let $B$ be the original universal ordinary determinant
ring, $\Lambda$ the original weight object, $R_H$ an
original \texttt{Unitary\allowbreak Hecke\allowbreak Data} record, and $K$ an
original permitted level.  Assume the source equality
$V.S'=s.S$.  For every original ideal $J\subset B$,
\[
                  R_H.\operatorname{IsJ}(B,K,J)
                         \quad\Longrightarrow\quad
                  B.\operatorname{Dense}(J).
 \tag{UPA.24}
\]
The original \texttt{Hypotheses\allowbreak H'} and all original
coefficient and rank hypotheses are retained; there
is no additional automorphic-Hecke-set premise.
\begin{proof}
The record $R_H$ supplies its literal type
$\tau_0$ and \texttt{Type\allowbreak For\allowbreak Alt} field.  UPA17
therefore supplies the precise automorphic-Hecke-set
input retained by AFS22.  Apply AFS22 to the original
$B,\Lambda,R_H,K,J$.

For clarity, the conclusion concerns the same ideal
throughout.  In the zero-module branch of
\texttt{Is\allowbreak J}, one has $J=B$, and the required membership
conclusion is immediate.  In the other branch,
\texttt{Is\allowbreak J} identifies $J$ with the kernel of the
specified map $f:B\to T$ to the original Hecke algebra.
For an arbitrary increasing weight, AFS1--AFS8 give
its dense good-weight set and the original arithmetic
point dictionary.  AFS20 separates elements of $T$;
AFS21 applies that separation to $f(b)$.
If $b$ vanishes at every original good point killing
$J$, then every separating point of $T$ kills $f(b)$,
so $f(b)=0$ and $b\in\ker f=J$.
This is (UPA.24), with no change of the ideal or
the original ambient ring.
\end{proof}
% END UPA18

\subsection{Scope after installation}

UPA16 supplies the actual Galois-decorated parts and
UPA17 installs them in the original ordinary Hecke setting.
UPA18 completes the previously conditional same-ideal density
endpoint.  The ordinary point output and its flags are developed
in the OFL and HEP chapters, and the full same-ideal
\texttt{Ring\allowbreak Props} record in the JRP chapter.  The later all-level
automorphy theorem and both entire rank-four support applications
retain their exact remaining inputs in the current audit.
The conditional irreducible-to-cuspidal assertions above
do not assert that all automorphic representations are
cuspidal or that every Speh multiplicity is one.


\section{Ordinary flags from the original local period data}
\label{sec:ordinary-flags-v44}

The purpose of this chapter is the full original implication
\(\mathrm{LiteratureRankNHypothesis}\Rightarrow\mathrm{OrdFlag}\).
Its conclusion is an actual ordered invariant flag over
\(\overline{\mathbb Q}_\ell\), together with potential crystallinity
for every original local bridge.  Equality of characteristic
polynomials alone is not substituted for that conclusion.

Fix a number field \(F\), its original algebraic closure \(\Omega\),
a place \(w\mid\ell\), a prolongation \(A\), and the original
field isomorphism \(\iota:\overline{\mathbb Q}_\ell\simeq\mathbb C\).
Write \(D\), \(I\), \(W\) for the decomposition, inertia, and Weil
groups.  The local Artin map is the source map
\(a:W\to F_w^\times\), with
\(\operatorname{ord}_w(a(s))=-\deg(s)\).
Every period assertion below is made using the original tied model
\(P=\mathrm{PP}(\ell)\) and an arbitrary inducing bridge
\(B_w\); finite coefficient extensions may be enlarged inside
\(\overline{\mathbb Q}_\ell\).  The isomorphism \(\iota\) is used
only algebraically and is never assumed continuous.

The retained published inputs are precisely the original global Artin
reciprocity and algebraic Hecke-character theorems
\texttt{N.\allowbreak aut.\allowbreak art}, \texttt{N.\allowbreak aut.\allowbreak ch\allowbreak A}, \texttt{N.\allowbreak aut.\allowbreak ch\allowbreak W}, the
original Chevalley unit-congruence input \texttt{N.\allowbreak loc.\allowbreak chev}, and
Berger's de Rham--potentially semistable theorem
\texttt{N.\allowbreak aut.\allowbreak berger}.  We also retain the original period-ring
foundations: finite coefficient models; comparison isomorphisms and
their tensor, dual, exactness, and finite-extension descent laws;
the fundamental identity
\[
 B_{\mathrm{cris}}^{\varphi=1}\cap
             \operatorname{Fil}^0B_{\mathrm{dR}}=\mathbb Q_\ell;
 \tag{OFL.1}
\]
and the crystalline comparison for an unramified rank-one character,
with Hodge--Tate weight zero and its actual Weil character.
These are the same foundations of the original \texttt{Local\allowbreak Published\allowbreak Theorems},
\texttt{Pst\allowbreak Laws}, and \texttt{Twist\allowbreak Above} period models.  They are
listed as retained foundations, rather than counted as new proofs.
The intermediate invariant-vector, weight-order, and flag assertions
used by the original source are all proved below.

\subsection{Rank-one characters, with their exact local normalization}

For an integer-valued function \(h\) on embeddings of \(F\), put
\[
 A_h(x)=\prod_{\tau:F_w\hookrightarrow\overline{\mathbb Q}_\ell}
       \tau(x)^{h(\tau|_F)}\quad(x\in F_w^\times).
 \tag{OFL.2}
\]
There are finitely many continuous embeddings.  Consequently the
product is nonzero, is multiplicative, and satisfies
\(A_{h+k}=A_hA_k\) and \(A_{-h}=A_h^{-1}\).
The literal \(\mathrm{OrdChar}(\iota,a,\omega,h,\psi)\)
says that \(\omega:F_w^\times\to\mathbb C^\times\) is smooth,
\(\psi:D\to\overline{\mathbb Q}_\ell^\times\) is continuous, and
\[
       \psi(s)=\iota^{-1}\bigl(\omega(a(s))\bigr)A_{-h}(a(s))
                         \quad(s\in W).
 \tag{OFL.3}
\]

% BEGIN OFL1
\paragraph{OFL1. The inertia algebraic part and inversion.}
\label{claim:OFL1}
An ordinary character is locally algebraic on inertia, with weight
\(h\).  Its inverse is an ordinary character with data
\((\omega^{-1},-h,\psi^{-1})\).
\begin{proof}
On inertia, multiply (OFL.3) by the character
\(\iota^{-1}(\omega\circ a)^{-1}\).  This factor has open
kernel: its kernel is the inverse image of \(\ker\omega\)
under the continuous map \(a|_I\).  Thus the product is exactly
the algebraic character \(A_{-h}\), which is the original
\(\mathrm{IsWtPart}\) convention.  If another allowed Artin map
is used, their original \(\mathrm{glob}\) clauses identify the
images of the local ideles after global Artin.  We verify the
needed injectivity in this profinite target.  Suppose the local
idele of \(x\in F_w^\times\) lies in the Artin kernel.  The
original global reciprocity law identifies that kernel with the
closure of the diagonal \(F^\times\) times the connected
infinity subgroup.  The latter has finite components one.
Consequently, for every positive integer \(a\), we can approximate
\(x\) arbitrarily closely at \(w\) by \(k\in F^\times\)
such that at every other finite place both \(k\) and \(k^{-1}\)
are integral and \((k-1)/a\) is integral.  Indeed these are open
conditions in the restricted-product topology, also imposed on
the inverse idele, and the given local idele satisfies them.

First \(x\) must be a local unit.  If \(|x|_w<1\), sufficiently
close approximants have a fixed positive \(w\)-valuation \(d\)
and valuation zero elsewhere.  Their principal ideal is
\(w^d\), so their integral norm has fixed absolute value
\(R=q^d\), where \(q=\operatorname{Norm}(w)\).
Choose \(a=q(R+1)+1\).  It is prime to \(w\) and greater
than \(R+1\).  A close approximant is integral everywhere,
and \((k-1)/a\) is integral also at \(w\), because \(a\)
is a local unit there.  Thus \(k=1+ay\) with \(y\in\mathcal O_F\).
The determinant defining its norm gives
\(\operatorname{Norm}_{F/\mathbb Q}(k)\equiv1\pmod a\).
An integer of absolute value \(R\) with this congruence must
equal one, contradicting \(d>0\).  Applying the same argument
to \(x^{-1}\) excludes \(|x|_w>1\).

Now fix \(N\ge1\), and let \(m\ge1\) annihilate the finite
unit group \((\mathcal O_F/q^N\mathcal O_F)^\times\).
The original \texttt{N.\allowbreak loc.\allowbreak chev} gives a positive integer
\(a\), prime to \(q\), such that every global unit congruent
to one modulo \(a\) is an \(m\)-th power of a global unit.
Choose an approximant with
\(|k-x|_w\le |q|_w^N<1\).  It is a unit at every finite
place, and \((k-1)/a\) is integral everywhere, by the same
prime-to-\(w\) argument.  Hence \(k=u^m\) for a global unit
\(u\), so \(k\equiv1\pmod{q^N}\).  The ultrametric
inequality yields \(|x-1|_w\le |q|_w^N\).
As \(N\) is arbitrary and \(|q|_w<1\), we obtain \(x=1\).
This proves injectivity after global Artin itself.  Applying it
to the quotient of the two local Artin values makes the two maps
equal, and supplies the universal Artin-map quantifier in
\(\mathrm{IsWtPart}\).  Taking inverses in (OFL.3)
and using (OFL.2) proves the second assertion.  The inverse has the
same smooth kernel and a continuous coefficient-field character.
\end{proof}
% END OFL1

% BEGIN OFL2
\paragraph{OFL2. A CM field realizing the fixed local field.}
\label{claim:OFL2}
For a finite extension \(K/\mathbb Q_\ell\) in the original
coefficient closure and the embedding \(j:\Omega\hookrightarrow
\overline{\mathbb Q}_\ell\) of a bridge, there is a CM subfield
\(E\subset\Omega\) whose image generates \(K\) over
\(\mathbb Q_\ell\).  It can be chosen with
\(0\ne\delta\in E\) such that every complex embedding sends
\(\delta\) to a purely imaginary number and \(j(\delta)\in
\mathbb Q_\ell\).
\begin{proof}
Choose a primitive element \(\alpha\) of \(K\).  The separability
of its minimum polynomial gives an open \(\ell\)-adic neighborhood
of its coefficients in which a root remains sufficiently close to
\(\alpha\) to generate \(K\).  To check the last assertion,
take disjoint balls around its conjugates; Hensel's lemma supplies
the nearby root, and Krasner's argument says that an automorphism
fixing that root must fix the nearest conjugate \(\alpha\).
Equality of the degrees then gives equality of the generated fields.
Independently choose a real monic polynomial of the same degree
having distinct real roots.  Having that many distinct real roots is
an open real condition, by disjoint intervals and sign changes.
Weak approximation for rational coefficients supplies one monic
polynomial satisfying both open conditions.  Its nearby root
\(\beta\in K\) is algebraic over \(\mathbb Q\), all its
conjugates are real, and \(\mathbb Q_\ell(\beta)=K\).
Lift \(\beta\) to \(b\in\Omega\) through \(j\): the
rational polynomial splits in \(\Omega\), so its image has every
root there.

The polynomial \(X^2-X+\ell\) has a simple root near zero in
\(\mathbb Z_\ell\), since its derivative there is \(-1\).
If that root is \(z\), then \(d=2z-1\) satisfies
\(d^2=1-4\ell\).  Choose a square root \(\delta\) of
\(1-4\ell\) in \(\Omega\), changing its sign so that
\(j(\delta)=d\).  Since \(1-4\ell<0\), every complex image
of \(\delta\) is nonzero and purely imaginary.
Now \(E=\mathbb Q(b,\delta)\) is the quadratic totally imaginary
extension of the totally real field \(\mathbb Q(b)\).  It is CM,
and its local image generates exactly \(K\).  This construction
also works at \(\ell=2\); no exclusion of the original prime
\(3\) occurs.
\end{proof}
% END OFL2

% BEGIN OFL3
\paragraph{OFL3. Globalizing the prescribed local algebraic weight.}
\label{claim:OFL3}
For every bridge \(B_w\) and every weight \(h\), a global
algebraic Hecke character on a field as in OFL2 supplies a continuous
local Galois character \(\nu\) with that same locally algebraic
weight.  The bridge can retain exactly its local field, embedding
\(j\), and coefficient splitting field.
\begin{proof}
Use OFL2 with the local field of \(B_w\).  Restrict the valuation
to \(E\) to obtain a place \(u\); the completion \(E_u\)
and \(F_w\) identify with the same finite field \(K\) inside
the coefficient closure.  Restriction through these identifications
identifies their decomposition groups.  Thus changing the global
field does not change the local field or period model.

Let \(S\) be the embeddings \(E\to\overline{\mathbb Q}_\ell\)
induced by embeddings of \(K\).  Every such embedding sends
\(\delta\) to the same nonzero \(d\in\mathbb Q_\ell\).
After composition with \(\iota\), complex conjugation changes
that image to its negative.  Hence \(\iota S\) and its conjugate
are disjoint.  An embedding in \(S\) determines a unique embedding
of \(F_w\): agreement on the generating field \(E\) determines
the map on \(K\).  Assign infinity exponent \(-h(\tau|_F)\)
on \(\iota S\), the opposite exponent on its conjugate, and
zero on the remaining conjugate pairs.  This gives an integral
infinity type \(n\) with \(n(\phi)+n(\bar\phi)=0\).

Here is the usual existence argument for a Hecke character of this
type, so that globalization is not a new local supplier assumption.
On units of a CM field, the algebraic expression prescribed by
\(n\) has absolute value one under every conjugate: complex
conjugation pairs the exponents with their negatives.  It is an
algebraic unit, so its values are roots of unity.  They lie in one
fixed number field and therefore form a finite group.  Choose a
finite congruence level on which all these unit values are one;
this can be seen by embedding that finite group in a local field
and choosing a sufficiently small neighborhood of one.  The
infinity character and the trivial character on this finite-level
compact subgroup then agree on their intersection with diagonal
\(E^\times\).  They define a character on the subgroup they
generate together with \(E^\times\), trivial on \(E^\times\).
Extend it to the idele group using divisibility of \(\mathbb C^\times\).
Its restriction to the stated open subgroup is continuous; hence
the extension is continuous.  This is an algebraic Hecke character
with exactly the prescribed infinity type.  No conductor restriction
is needed for the present endpoint.

The retained algebraic Hecke-character theorem supplies its
\(\ell\)-adic Galois character \(\mu\).  Pull it to \(D\)
through the common local field to define \(\nu\).  The Artin
compatibility in that theorem gives inertia weight \(h\), since
the chosen exponent was \(-h\).  To check compatibility of the
two global fields, take their compositum in \(\Omega\): its
chosen completion is still \(K\), and both restriction maps
induce the identity there.  Local reciprocity on the compositum
therefore identifies the two inertia formulas.  The induced
embeddings have the same period multiplicities because the
underlying local representation and local field are literally the
same.  This proves all assertions, including the original bridge
quantifiers.
\end{proof}
% END OFL3

% BEGIN OFL4
\paragraph{OFL4. The full rank-one \texttt{DRat} and potential-crystalline law.}
\label{claim:OFL4}
If \(\mathrm{OrdChar}(\iota,a,\omega,h,\psi)\), then
\[
 \mathrm{DRat}\!\left(\iota,w,\operatorname{rk}_1\psi,
                 \tau\mapsto\{h(\tau)\},
                 s\mapsto X-\omega(a(s))\right),
 \qquad \mathrm{PotCrys}(\operatorname{rk}_1\psi).
 \tag{OFL.4}
\]
These assertions hold for every original inducing bridge.
\begin{proof}
Fix such a bridge.  OFL1 and OFL3 give a globalized local character
\(\nu\) with the same inertia algebraic weight as \(\psi\).
Their quotient \(\eta=\psi\nu^{-1}\) has open inertia kernel:
intersect the two open kernels of their smooth inertia parts.
There is a finite extension of the local field whose inertia is
in that kernel.  Indeed an open identity neighborhood in inertia
contains the intersection with an open normal subgroup of the
absolute Galois group; the latter corresponds to a finite extension.
On that extension \(\eta\) is unramified.  The original
rank-one crystalline comparison makes it de Rham with jump zero.
Finite-extension descent preserves being de Rham and its labeled
jump multiplicities.  The retained Hecke-character theorem
\texttt{ch\allowbreak A}, applied to \(\nu\), gives jump \(h(\tau)\).
Tensoring with \(\eta\) gives the same jumps for \(\psi\).
Induced embeddings of \(F\) extend uniquely to \(F_w\), as
in UPA11, so this proves the exact \(\mathrm{IsJumpsAt}\)
quantifier, not only a list indexed by a chosen embedding family.

The retained algebraic-character WD theorem \texttt{ch\allowbreak W}, together
with the unramified rank-one comparison and its finite-extension
descent, identifies the parameter of \(\psi\) with the smooth
character
\[
 \xi(s)=\iota\bigl(\psi(s)A_h(a(s))\bigr).
 \tag{OFL.5}
\]
More explicitly, this identity holds first for each globalized
weight character: it is the Artin formula in \texttt{ch\allowbreak W} on
the common local field.  Tensor products and inverse characters
give arbitrary integral weights.  For the quotient \(\eta\)
the Weil factor is its actual character \(\iota\eta|_W\),
since the period comparison for a potentially unramified character
uses the original Weil action.  Multiplying these identities yields
(OFL.5); finite descent recovers all inertia elements and not just
a Frobenius power.  This also verifies the source degree sign.
By (OFL.3), (OFL.5) is exactly \(\omega(a(s))\), so its
characteristic polynomial is the polynomial in (OFL.4).

Finally apply \texttt{N.\allowbreak aut.\allowbreak berger} to this de Rham representation.
Its WD space has dimension one, and a nilpotent endomorphism of a
one-dimensional space is zero.  The retained equivalence between
zero monodromy and potential crystallinity therefore proves the
second assertion.  The bridge was arbitrary throughout.
\end{proof}
% END OFL4

\subsection{Distinct inertia and a fixed period}

% BEGIN OFL5
\paragraph{OFL5. The finite-inertia character projectors.}
\label{claim:OFL5}
Let \(I\to\Gamma\) be normal in the sense that its image is
preserved by conjugation.  Let \(r\) be a complex representation
with finite inertia image, and suppose
\[
 \det(X-r(g))=\prod_{j=0}^{n-1}(X-\mu_j(g))
                                 \quad(g\in\Gamma),
 \tag{OFL.6}
\]
where the characters \(\mu_j|_I\) are pairwise distinct.
Then inertia is the direct sum of its one-dimensional
\(\mu_j\)-spaces.  Every nilpotent endomorphism commuting with
inertia is zero.
\begin{proof}
If \(r(s)=1\) for \(s\in I\), (OFL.6) gives
\(\mu_j(s)=1\) for all \(j\).  Thus all these characters
factor through the finite group \(H=r(I)\).  For their induced
characters \(\chi_j\), form
\[
        e_j=|H|^{-1}\sum_{h\in H}\chi_j(h)^{-1}r(h).
 \tag{OFL.7}
\]
Here \(r(h)\) denotes the faithful action of the image group.
The finite character sum of a nontrivial character is zero:
translate the sum by an element on which that character is not
one.  It follows that \(e_j e_k=\delta_{jk}e_j\) and that
the image of \(e_j\) is the \(\chi_j\)-eigenspace.
The trace identity obtained from (OFL.6) shows
\(\operatorname{tr}(e_j)=1\).  Since an idempotent has
eigenvalues zero and one, its image has dimension one.
The degree in (OFL.6) gives \(\dim r=n\); hence these \(n\)
lines sum to the entire space.  A commuting nilpotent preserves
each line, acts there by a nilpotent scalar, and is zero on all
lines.  This proves the assertion, including the empty rank-zero
case.
\end{proof}
% END OFL5

% BEGIN OFL6
\paragraph{OFL6. The inertia eigenline has its prescribed full Weil character.}
\label{claim:OFL6}
Under OFL5, \(r(g)\) acts on the \(j\)-th inertia line by
\(\mu_j(g)\), for every \(g\in\Gamma\).
\begin{proof}
Normality of the inertia image and the fact that \(\mu_j\)
is a character of all of \(\Gamma\) show that its inertia
eigenspace is preserved by every \(g\).  Let its scalar there
be \(c_j(g)\).  Taking the trace after applying (OFL.7) gives
\[
 c_j(g)=\operatorname{tr}(r(g)e_j)
  =|H|^{-1}\sum_{h\in H}\chi_j(h)^{-1}
                        \sum_k\mu_k(g)\chi_k(h)
  =\mu_j(g).
\]
Consequently every vector fixed by inertia is fixed by the whole
Weil group when one of the \(\mu_j\) is the trivial character:
that character has the only trivial inertia restriction, by
distinctness.  In the WD application, normality is checked directly
from \(\deg(gsg^{-1})=0\).
\end{proof}
% END OFL6

% BEGIN OFL7
\paragraph{OFL7. Descending the invariant vector through WD scalarization.}
\label{claim:OFL7}
In the original potentially semistable period model, use the raw
complex Weil action \(\mathrm{wdRho}\) on \(\mathrm{WDSpace}\),
before Frobenius semisimplification.  A nonzero
inertia-fixed vector of this complex space yields a nonzero
inertia-fixed vector of \(D_{\mathrm{pst}}\).  If all the former
inertia-fixed vectors are Weil fixed, so are all the latter.
\begin{proof}
The original scalarization map \(\mathrm{wdMk}:D_{\mathrm{pst}}
\to\mathrm{WDSpace}\) is additive, intertwines the period Weil
action with this raw \(\mathrm{wdRho}\) action,
and has complex spanning image.  The inertia action on the period
space has open kernel, hence finite image.  Average over this finite
image.  If that average were zero on every period vector, its
complexification would vanish on the spanning image and therefore
on the whole WD space.  On a nonzero inertia-fixed WD vector,
however, the unnormalized average is multiplication by the positive
integer equal to the group order.  This contradiction gives a
nonzero inertia-fixed period vector.

For the second assertion, let \(x\) be any such period vector.
The period Frobenius \(\varphi\) commutes with the corrected
Weil action.  Hence \(\varphi^m x\) is inertia fixed for each
\(m\ge0\), and its scalarized image is fixed by every Weil
element.  For \(z=\mathrm{weilD}(s)x-x\) we obtain
\[
             \mathrm{wdMk}(\varphi^m z)=0\quad(m\ge0).
 \tag{OFL.8}
\]
The original faithfulness of all these Frobenius translates implies
\(z=0\).  To recall that faithfulness, scalarization selects one
embedding component of the maximal unramified coefficient algebra;
successive Frobenius maps visit every component over any finite
unramified subextension containing the coordinates of \(z\).
Vanishing on all components is vanishing of \(z\).  Thus one
does not incorrectly assume that a single scalarization map is
injective.  This proves the assertion on the original period space.
\end{proof}
% END OFL7

% BEGIN OFL8
\paragraph{OFL8. Turning a Weil-fixed period into a Frobenius-fixed period.}
\label{claim:OFL8}
Suppose \(D_{\mathrm{pst}}\) contains a nonzero inertia-fixed
vector and every inertia-fixed vector is Weil fixed.  Then it
contains \(z\ne0\) fixed both by \(\varphi\) and by the
entire local Galois action.
\begin{proof}
Let \(x\ne0\) be inertia fixed.  Every period vector is fixed
by an open local Galois subgroup, so some positive power of inverse
Frobenius fixes \(x\) under the Galois action.  It also fixes
\(x\) under the corrected Weil action.  Since that corrected
action is the Galois action times
\(\varphi^{-\alpha(s)}\), with
\(\alpha(s)=f_w\deg(s)\), it follows that
\(\varphi^m x=x\) for some integer \(m\ge2\).
Replacing the initial exponent by twice itself ensures this bound.

Choose \(c\) in the maximal unramified extension of
\(\mathbb Q_\ell\) with Frobenius orbit of length exactly
\(m\).  For example take a primitive
\((\ell^m-1)\)-st root of unity.  Frobenius sends it to
its \(\ell\)-th power, and
\(\ell^i<\ell^m-1\) for \(0\le i<m\), \(m\ge2\),
so the indicated orbit has that length.  For integers \(j\ge0\),
set
\[
                  z_j=\sum_{i=0}^{m-1}\varphi^i(c^j x).
 \tag{OFL.9}
\]
The identity \(\varphi^m(c)=c\) and
\(\varphi^m x=x\) makes \(\varphi z_j=z_j\).
At least one of \(z_0,\ldots,z_{m-1}\) is nonzero: otherwise
the Vandermonde matrix in the distinct scalars
\(c,\varphi(c),\ldots,\varphi^{m-1}(c)\) annihilates the
vector \((x,\varphi x,\ldots,\varphi^{m-1}x)\), forcing
\(x=0\).  This argument works in the vector space over the
maximal unramified field and does not require a choice of a
Frobenius eigenvalue in \(\mathbb C\).

Each \(z_j\) is inertia fixed: inertia fixes the unramified
scalars, commutes with \(\varphi\), and fixes \(x\).
The chosen nonzero \(z_j\) is therefore Weil fixed as well.
Since it is already Frobenius fixed, the corrected and uncorrected
Weil actions agree on it.  Its stabilizer for the original Galois
action is open and closed, and the Weil group is dense in the
decomposition group.  It is thus fixed by every Galois element.
\end{proof}
% END OFL8

% BEGIN OFL9
\paragraph{OFL9. Nonnegative Hodge--Tate weights recover a Galois invariant.}
\label{claim:OFL9}
Let \(V\) be a de Rham representation over a finite coefficient
field, with all its labeled Hodge--Tate weights nonnegative in the
source convention.  If \(0\ne z\in D_{\mathrm{pst}}(V)\)
is killed by monodromy and fixed by both Frobenius and Galois, then
\(V\) has a nonzero Galois-fixed vector.
\begin{proof}
Galois invariance puts \(z\) in \(D_{\mathrm{st}}(V)\).
The kernel of monodromy there is the image of
\(D_{\mathrm{cris}}(V)\), so write its unique antecedent as
\(y\).  The crystalline-to-semistable inclusion is injective and
commutes with Frobenius and Galois.  Thus \(y\) is fixed by both.
Its de Rham image lies in \(D_{\mathrm{dR}}(V)\).

Here is the filtration step with its sign.  Enlarge coefficients
to contain every embedding of the local field.  The source labeled
multiplicity is the dimension of the corresponding graded
filtration piece.  Every negative graded piece is zero by the
nonnegative-weight assumption.  A finite-dimensional, separated,
exhaustive decreasing filtration therefore satisfies
\(D_{\mathrm{dR}}=\operatorname{Fil}^0D_{\mathrm{dR}}\):
start at a sufficiently negative index where the filtration is the
whole space and use equality of successive pieces up to zero.
Descent through the enlarged coefficient field preserves this
identity.

Choose a \(\mathbb Q_\ell\)-basis of \(V\), and write
\(y=\sum_i v_i\otimes b_i\).  Its Frobenius invariance says
\(\varphi(b_i)=b_i\), and its membership in filtration zero
says \(b_i\in\operatorname{Fil}^0B_{\mathrm{dR}}\).
The fundamental identity (OFL.1) gives \(b_i\in\mathbb Q_\ell\).
Hence \(y=v\otimes1\) for a unique \(v\in V\).
The nonzero and invariant assertions descend through this injective
tensor map.  This gives \(v\ne0\) and \(g(v)=v\) for every
local Galois element, as required.
\end{proof}
% END OFL9

% BEGIN OFL10
\paragraph{OFL10. The complete local core of the slope argument.}
\label{claim:OFL10}
Let \(r:D\to\mathrm{GL}_n(\overline{\mathbb Q}_\ell)\)
be continuous, and assume the original \(\mathrm{DRat}\) with
weights \(H(\tau)\) and polynomials
\(\prod_j(X-\mu_j(s))\), where the smooth Weil characters
\(\mu_j\) have pairwise distinct inertia restrictions.  Then
\(r\) is potentially crystalline for every original bridge.
If some \(\mu_j=1\) and all \(H(\tau)\) are nonnegative,
there is a nonzero vector fixed by every element of \(D\).
\begin{proof}
For an arbitrary bridge the \(\mathrm{DRat}\) datum gives
de Rhamness and the asserted WD polynomials.  Berger's theorem
gives a potentially semistable finite coefficient model.  We
apply OFL5--OFL8 to its raw complex Weil action
\(\mathrm{wdRho}\), which is the action actually intertwined
by \(\mathrm{wdMk}\).  The semisimplified action
\(\mathrm{wdRep}\) defining \(\mathrm{wdPst}\) differs from
it by a commuting unipotent factor raised to \(\deg(s)\).
Commuting matrices are simultaneously triangularizable over
\(\mathbb C\); this factor has diagonal entries one, so the
two actions have the same characteristic polynomial on every
Weil element.  On inertia the factor is the identity.  Thus the
\(\mathrm{DRat}\) polynomials in the hypothesis are exactly
the raw-action polynomials required in OFL5--OFL6, and its
inertia image is finite.  No invariant vector is being inferred
merely from a Frobenius-semisimplified action.
The monodromy operator is
nilpotent and commutes with inertia, since inertia has degree zero
in the WD commutation identity.  OFL5 therefore makes that
operator zero.  The original monodromy-zero comparison gives
potential crystallinity, including the specified finite model.

When \(\mu_j=1\), OFL5--OFL6 give the nonzero trivial
inertia line and trivial Weil action on it.  OFL7--OFL8 produce
the nonzero period fixed by Frobenius and Galois.  OFL9 gives a
nonzero invariant vector in the finite coefficient model.
Embedding its coordinates in \(\overline{\mathbb Q}_\ell\)
preserves nonzeroness and the matrix identities.  This yields
the invariant of \(r\) itself.  For this existence assertion
one may choose any inducing bridge, whose existence is part of
the original local model; potential crystallinity was proved for
each bridge separately.  Rank zero is harmless: the latter
assertion is then vacuous because there is no index \(j\).
\end{proof}
% END OFL10

\subsection{Twisting, quotienting, and the ordered flag}

% BEGIN OFL11
\paragraph{OFL11. The exact \texttt{DRat} twist formula.}
\label{claim:OFL11}
Suppose \(r\) has \(\mathrm{DRat}\) weights \(H(\tau)\)
and polynomials \(\prod_j(X-\alpha_j(s))\), and a continuous
rank-one character \(\eta\) has weights \(c(\tau)\) and
polynomials \(X-\beta(s)\).  Then
\(r'(g)=\eta(g)r(g)\) is continuous and has weights
\(\{h+c(\tau):h\in H(\tau)\}\) and polynomials
\[
                    \prod_j(X-\alpha_j(s)\beta(s)).
 \tag{OFL.10}
\]
\begin{proof}
Continuity follows entry by entry, and the inverse matrix is
\(\eta(g)^{-1}r(g)^{-1}\), so it is continuity into the
original general linear group.  Work over a common finite
coefficient extension containing the two models.  Tensoring a de
Rham representation with a de Rham line preserves its comparison
isomorphism.  A filtered line with jump \(c\) shifts every
filtration index by \(c\), giving exactly the displayed multiset.
Berger makes the representations potentially semistable.  On
their period modules tensoring with the line multiplies the Weil
matrix by its scalar \(\beta(s)\); the monodromy on a line
is zero.  A scalar multiple of a matrix with roots
\(\alpha_j(s)\) has roots \(\alpha_j(s)\beta(s)\), with
multiplicities, so its monic characteristic polynomial is
(OFL.10).  These identifications commute with all inducing
embeddings and with coefficient enlargement.  They therefore
prove the full original \(\mathrm{DRat}\), for every bridge.
\end{proof}
% END OFL11

% BEGIN OFL12
\paragraph{OFL12. Removing a character line preserves the remaining data.}
\label{claim:OFL12}
Suppose \(r\) contains an invariant line with character
\(\psi_0\).  Assume that \(r\) has \(\mathrm{DRat}\)
weights \(\{h_0(\tau),\ldots,h_{n-1}(\tau)\}\) and
polynomials \(\prod_{j=0}^{n-1}(X-\alpha_j(s))\), and the
line has weights \(h_0\) and polynomial \(X-\alpha_0(s)\).
Then the quotient by this line is continuous and has the
\(\mathrm{DRat}\) data with index zero deleted.
\begin{proof}
Put a nonzero line vector first in a basis.  The matrix of \(r\)
is upper block triangular, with top left entry \(\psi_0\),
zero lower left block, and quotient block \(r'\).
The latter entries are continuous matrix coordinate functions;
its inverse is the same block of \(r(g^{-1})\), so it gives
a continuous quotient representation.

Fix any original bridge and enlarge coefficients to contain the
basis, the character, and the representation models.  Exactness
of de Rham comparison makes the quotient de Rham, and strictness
of the filtered exact sequence makes every labeled jump
multiplicity additive.  Cancellation of the singleton
\(\{h_0(\tau)\}\) gives the required quotient multiset,
even when weights repeat.  Choose a finite extension where all
three representations are semistable, using Berger for the de
Rham representations.  The exact WD sequence is upper block
triangular, hence its characteristic polynomial is the product
of those of the line and quotient.  Cancel the nonzero monic
polynomial \(X-\alpha_0(s)\) in \(\mathbb C[X]\).
This gives the required quotient polynomial for every Weil
element.  The comparison and finite descent laws identify this
calculation with the original bridge's parameter.  Thus every
clause of \(\mathrm{DRat}\) descends, not only its multiset
of Frobenius eigenvalues.
\end{proof}
% END OFL12

% BEGIN OFL13
\paragraph{OFL13. The first ordinary eigenline.}
\label{claim:OFL13}
Let \(n\ge1\), and suppose the ordinary characters
\((\omega_j,h_j,\psi_j)\) satisfy
\(h_0(\tau)<\cdots<h_{n-1}(\tau)\) for every embedding
\(\tau\), with pairwise distinct inertia restrictions of
\(\omega_j\circ a\).  If \(r\) has \(\mathrm{DRat}\)
weights \(\{h_j(\tau)\}_j\) and polynomials
\(\prod_j(X-\omega_j(a(s)))\), then \(r\) has a nonzero
eigenvector with character \(\psi_0\).
\begin{proof}
Twist \(r\) by \(\psi_0^{-1}\).  OFL1 and OFL4 give
the exact rank-one inverse data, and OFL11 gives the twisted
weights \(h_j(\tau)-h_0(\tau)\) and Weil characters
\(\mu_j=(\omega_j\omega_0^{-1})\circ a\).
They are pairwise distinct on inertia, because multiplying by the
same nonzero character preserves inequalities.  The zeroth
character is trivial.  Every shifted weight is nonnegative by
the stated order.  OFL10 gives a nonzero invariant vector \(v\)
of the twist.  The equality
\(\psi_0(g)^{-1}r(g)v=v\) is precisely
\(r(g)v=\psi_0(g)v\) for every \(g\in D\).
This constructs the first line in its prescribed order; it does
not choose an arbitrary inertia eigenline.
\end{proof}
% END OFL13

% BEGIN OFL14
\paragraph{OFL14. The full ordered invariant flag.}
\label{claim:OFL14}
Under the hypotheses of OFL13, allowing also \(n=0\), there is
a complete \(D\)-stable flag
\[
 0=V_0\subset V_1\subset\cdots\subset V_n
             =\overline{\mathbb Q}_\ell^n,
 \qquad D\text{ acts on }V_{j+1}/V_j\text{ by }\psi_j.
 \tag{OFL.11}
\]
\begin{proof}
Induct on \(n\).  In rank zero the empty flag has the required
property.  For positive rank, OFL13 gives the first eigenline
\(V_1\).  Complete its nonzero vector to a basis, so that the
quotient is the continuous representation in OFL12.  OFL4 provides
the rank-one \(\mathrm{DRat}\) for \(\psi_0\), and OFL12
gives the quotient data with index zero removed.  The remaining
weights are still strictly increasing, and their smooth inertia
characters are still pairwise distinct.  Induction gives their
ordered flag in the quotient.  Taking its inverse images and
preceding them by \(0\subset V_1\) gives (OFL.11).

The concatenated basis consists of the first eigenvector followed
by lifts of the quotient basis.  In it every matrix is upper
triangular with diagonal
\((\psi_0(g),\ldots,\psi_{n-1}(g))\).  Its change-of-basis
matrix is one invertible matrix over the original coefficient
closure, independent of \(g\).  This is exactly the source
\(\mathrm{HasFlag}(r,\psi)\), with the original order and
without a completion of the coefficient field.
\end{proof}
% END OFL14

% BEGIN OFL15
\paragraph{OFL15. Distinct unit characters give the original inertia hypothesis.}
\label{claim:OFL15}
The original pairwise distinct restrictions of \(\omega_j\) to
\(\mathcal O_{F_w}^\times\) imply the inertia distinctness used
in OFL5--OFL14.
\begin{proof}
For \(i\ne j\), choose a unit \(u\) with
\(\omega_i(u)\ne\omega_j(u)\).  Surjectivity of the
original local Artin map gives \(s\in W\) with \(a(s)=u\).
Its degree is zero, because
\(0=\operatorname{ord}_w(u)=-\deg(s)\).  Thus \(s\)
is the image of an inertia element, which supplies the required
inequality.  This argument uses the exact source reciprocity
sign and does not require an additional choice of uniformizer.
\end{proof}
% END OFL15

% BEGIN OFL16
\paragraph{OFL16. The literal \texttt{ord\allowbreak Flag\_\allowbreak alt} endpoint.}
\label{claim:OFL16}
The original \(\mathrm{LiteratureRankNHypothesis}\) implies
the complete original \(\mathrm{OrdFlag}\), at every allowed
prime and rank, on every original field, place, Artin map, and
coefficient isomorphism.
\begin{proof}
Unpack an arbitrary test datum in \(\mathrm{OrdFlag}\).
Its characters satisfy (OFL.3), its labeled weights are strictly
increasing in the stated order, its unit characters are distinct,
and its continuous representation satisfies the exact
\(\mathrm{DRat}\) polynomial identity.  OFL15 converts the
unit condition to inertia distinctness.  OFL14 gives
\(\mathrm{HasFlag}\) with exactly the supplied characters.
For every inducing bridge, OFL10 applied to the original,
untwisted representation gives potential crystallinity.  These
are the two conclusions of the original \(\mathrm{OrdFlag}\).
The original odd-prime hypothesis, when present in the consumer,
is retained; none of the steps discards the prime \(3\) or
requires a new lower bound on it.  No semisimplicity hypothesis
on \(r\) is used or added.
\end{proof}
% END OFL16

\paragraph{Boundary of this chapter.}
The original Hecke eigencharacter calculation is still needed to
produce the \(\mathrm{OrdChar}\) inputs at arithmetic Hecke
points.  That calculation is supplied separately by the HEP
chapter.  This chapter proves their complete ordered-flag consumer,
while retaining the named published and period-ring foundations
listed at the start.


\section{Hecke eigenpoints, their full determinants, and the original ordinary output}
\label{sec:hecke-eigenpoints-v44}

This chapter reconstructs the full original \texttt{TCore\allowbreak Eig} theorem
and the entire \texttt{Ord\allowbreak Core\allowbreak H.\allowbreak Out} record.  An arbitrary arithmetic
point of the specified Hecke algebra gives one continuous semisimple
representation, its complete determinant law, and the prescribed point
property.  At a good point it also gives the actual ordered ordinary
flags and potential crystallinity at every original local bridge.
The algebra, character labels, coefficient maps, and level are unchanged.

Fix the original data $L,\Omega,\ell,n,c,H$ and $H.\mathrm{Ax}$.
Write $O=c.O$, $\overline E=\overline{\mathbb Q}_\ell$, and fix the
given algebraic isomorphism $\iota:\overline E\simeq\mathbb C$.
Let $A$ have the original class $\mathrm{IsCO}$, let
$H.\mathrm{IsWt}(\theta)$ hold, and fix the residual operator labels
$a$ and the original level $K$.  Put
\[
 N=H.MK(\theta,a,K),\qquad T=H.TK(\theta,a,K)
       \subset\operatorname{End}_A(N).
 \tag{HEP.1}
\]
We retain the source assumptions that the specified operators commute
on $N$ and that $T$ is of class $\mathrm{IsCO}$ over $O$.
For an original operator index $\alpha$, write $t_\alpha$ for its
actual generator in $T$.  Suppose given
\[
 D\in\mathrm{GDet}_n(T,H.G_q),\qquad
 \psi_{v,i}:D_v\longrightarrow T^\times,
 \qquad H.\mathrm{IsUniv}(\theta,a,K,\mathfrak m_T,D,\psi).
 \tag{HEP.2}
\]
Here $D_v=\operatorname{Dec}(H.\mathrm{pl}(v))$, and all local maps
are the source maps.  No domain assumption or flatness of a point
map is imposed on $A$ or $T$.

The original $\mathrm{LiteratureRankNHypothesis}$ is retained.
The specific arithmetic foundations used below are its global Artin
reciprocity and local idele injectivity clauses, with the local--global
and norm normalizations of class field theory; its local rank-one,
unlinked principal-series, and Langlands laws; and the original
$\mathrm{RecTExists}$.  These are also the explicit foundations of
the earlier CIP, AFS, GUA, and OFL packets.  OFL16 carries precisely
its stated period foundations.  None is relabeled as a newly proved
class field theory, local Langlands, or period comparison theorem here.

\subsection{Continuity and equality of complete determinant laws}

% BEGIN HEP1
\paragraph{HEP1. Every original algebra point preserves adic continuity.}
\label{claim:HEP1}
Let $R$ be any $O$-algebra of class $\mathrm{IsCO}$ and
$y:R\to_O\overline E$ any algebra point.  If
$f:\Gamma\to R$ is adically continuous for $\mathfrak m_R$,
then $y\circ f$ is continuous for the original topology of
$\overline E$.  In particular every evaluated determinant
coefficient and every evaluated $\psi_{v,i}$ in (HEP.2) is continuous.
The evaluation of $H.\mathrm{epsV}(R,v)$ is the original
cyclotomic character $\varepsilon_\ell|_{D_v}$.
\begin{proof}
ASP2, applied to the very point $y$, gives $0<\eta<1$ with
\[
             |y(b)|_\ell\leq\eta^r
                  \qquad(b\in\mathfrak m_R^r).
 \tag{HEP.3}
\]
Fix $g\in\Gamma$ and $\epsilon>0$.  Choose $r$ with
$\eta^r<\epsilon$.  Adic continuity says that, on a neighborhood
of $g$, $f(g')-f(g)\in\mathfrak m_R^r$.  Thus
$|y(f(g'))-y(f(g))|_\ell<\epsilon$ there.  This proves continuity
at every $g$ without a continuity assumption on the point itself.

The \texttt{cont} and \texttt{\(\psi\)\_\allowbreak cont} fields of
$\mathrm{IsUniv}$ supply these adic hypotheses.  The resulting
character into $\overline E^\times$ is continuous in the topology
on units: its value is continuous by the argument just given, and
its inverse value is the value at $g^{-1}$, hence is continuous too.
Finally $\mathrm{epsV}$ is defined by the structural maps
$\mathbb Z_\ell\to O\to R$ applied to $\mathrm{cycZ}$.
Since $y$ is an $O$-algebra map, evaluation is the same structural
map to $\overline E$.  This proves the stated cyclotomic identity.
\end{proof}
% END HEP1

% BEGIN HEP2
\paragraph{HEP2. Dense characteristic polynomials determine the whole evaluated law.}
\label{claim:HEP2}
Let $\Gamma$ be a topological group and let $D_1,D_2$ be
degree-$n$ determinants over a characteristic-zero field $E$,
with its Hausdorff topology.  Suppose their characteristic
coefficients on every group element are continuous.  If their
characteristic polynomials agree on a dense subset of $\Gamma$,
then $D_1=D_2$ as polynomial laws, after every commutative
$E$-algebra extension.
\begin{proof}
Each difference of corresponding coefficient functions is continuous.
Its zero set is closed and contains a dense set, so it is all of
$\Gamma$.  In particular the two trace functions agree on every
group element.

We give the algebra argument that upgrades this assertion to full
polynomial laws.  For any multiplicative determinant $\Delta$,
write
\[
 \Delta(1-tz)=\sum_{k=0}^n(-1)^k\lambda_k(z)t^k,
 \qquad \lambda_0=1,\quad \operatorname{tr}_\Delta=\lambda_1.
 \tag{HEP.4}
\]
These expressions are natural in every scalar algebra.  In a ring
of dual numbers,
$(1+\epsilon x)(1+\epsilon z)=1+\epsilon(x+z)$.
Multiplicativity therefore makes the coefficient of $\epsilon$
additive in $x,z$; substitution $\epsilon\mapsto b\epsilon$
makes it linear over the scalar algebra.  Thus
$\operatorname{tr}_\Delta$ is a natural linear map.  Agreement on
$\Gamma$ implies agreement on every element of $B[\Gamma]$, for
every $E$-algebra $B$, by finite linear expansion in group elements.

For completeness, multiplicativity also gives the Newton identities
in the form needed here.  In formal power series, or in their
arbitrary finite truncations, factor
\[
 1-(t+\epsilon)z=(1-tz)
       \left(1-\epsilon z(1-tz)^{-1}\right).
\]
Apply the natural law to this identity and compare the coefficient
of $\epsilon$.  Since $(1-tz)^{-1}=\sum_{r\geq0}t^rz^r$, it gives
\[
 \frac{d}{dt}\Delta(1-tz)
   =-\Delta(1-tz)\sum_{r\geq0}
                        \operatorname{tr}_\Delta(z^{r+1})t^r.
\]
Each coefficient is justified in a finite truncation, so no
convergence or continuity of a polynomial law is involved.
Equating coefficients yields
\[
 k\lambda_k(z)=\sum_{j=1}^k(-1)^{j-1}
                   \lambda_{k-j}(z)\operatorname{tr}_\Delta(z^j).
 \tag{HEP.5}
\]
Every positive integer is invertible in $B$, because $B$ is an
$E$-algebra.  Induction on $k$ now identifies all the $\lambda_k$
of $D_1$ and $D_2$ on every $B[\Gamma]$.  Homogeneity identifies
$\lambda_n(z)$ with the value of the degree-$n$ law on $z$.
Hence the laws are equal on every scalar algebra, as claimed.
For $n=0$ this last assertion is the constant unital law and is
immediate.  In its application here, division by integers occurs
only after passage to $\overline E$; no integer is inverted in
$O,A$, or $T$.
\end{proof}
% END HEP2

% BEGIN HEP3
\paragraph{HEP3. A spherical eigenvector identifies the full specialized determinant.}
\label{claim:HEP3}
Let $y:T\to_O\overline E$, let $\lambda,\xi$ be any original
weight and coefficient representation, and let
$Y:H.\mathrm{PartAltAlt}(\iota,\lambda,\xi)$.
Suppose $0\ne F\in Y.Y$ satisfies
\[
 H.\operatorname{opC}(\iota,\lambda,d,(w,j))F
       =\iota(y(t_{(w,j)}))F
       \quad(w\in H.W_p,\ 0\leq j<n).
 \tag{HEP.6}
\]
Then, as full determinant laws,
\[
                    D\otimes_{T,y}\overline E
                             =\det(Y.r).
 \tag{HEP.7}
\]
No arithmetic or good-point assumption is required here.
\begin{proof}
The original index $(w,j)$ denotes the spherical operator of
degree $j+1$.  Its action on $Y.Y$ is also the scalar
$Y.a(w,j+1)$ by the part's \texttt{sph} field.
Since $F\ne0$, equality of the two scalar multiples implies
\[
       \iota\bigl(y(H.\operatorname{genW}(w,k))\bigr)=Y.a(w,k)
                                      \quad(1\leq k\leq n).
\]
At $k=0$ both sides equal one by the literal definition of
$\mathrm{genW}$ and the part's \texttt{a\_\allowbreak zero} field.

Scalar extension of a determinant sends its characteristic
polynomial to the coefficientwise image of that polynomial:
apply naturality to the group-algebra element $X\cdot1-g$.
The same is immediate for the finite sum defining
$\mathrm{heckePol}$.  Thus the \texttt{frob} fields of
$\mathrm{IsUniv}$ and $Y$ give equality of the characteristic
polynomials in (HEP.7) at every original element
$H.\mathrm{prj}(\sigma^{-1})$, for every allowed spherical place,
prolongation, and arithmetic Frobenius $\sigma$.
The fixed $\iota$ is injective, so equality after its coefficient
map gives equality over $\overline E$ itself.

The actual set of these inverse Frobenius elements is dense in
$H.G_q$ by the original $H.\mathrm{Ax.dense}$ field.
The evaluated coefficients of $D$ are continuous by HEP1.
The coefficients of the characteristic polynomial of $Y.r(g)$
are continuous finite polynomials in its matrix entries, because
$Y.\mathrm{cont}$ holds.  HEP2 therefore proves (HEP.7), with
its full scalar-algebra quantifier.
\end{proof}
% END HEP3

\subsection{Recovering each ordered local character}

At $v\in H.V$ put $k=L_{H.S_l(v)}$, $q=\#\kappa_k$, and
$b=H.\mathrm{art}(v):W_v\to k^\times$.
For any integer-valued function $h$ on the global embeddings put
\[
 A_h(z)=\prod_{\tau:k\hookrightarrow\overline E,\ \tau\ \text{continuous}}
                       \tau(z)^{h(\tau|_L)}.
 \tag{HEP.8}
\]
This is the literal $\mathrm{algPow}$ on nonzero arguments.
Write $l_i=\operatorname{lowW}(\lambda,i)$ and
$h_i=\operatorname{htW}(\lambda,i)=l_i+i$.

% BEGIN HEP4
\paragraph{HEP4. Algebraic factors and the full Weil cyclotomic formula.}
\label{claim:HEP4}
The products (HEP.8) are finite, nonzero, and multiplicative, with
$A_{h+h'}=A_hA_{h'}$ and $A_{-h}=A_h^{-1}$.
For every original $s\in W_v$,
\[
 \varepsilon_\ell(s)=q^{\deg(s)}A_1(b(s)),\qquad
 A_{-h_i}(z)A_{l_i}(z)A_1(z)^i=1.
 \tag{HEP.9}
\]
Both statements use the given original Artin map and embedding set.
\begin{proof}
WRP1 identifies the continuous embeddings with the embeddings of
the finite extension $k/\mathbb Q_\ell$.  There are finitely many,
and every field image of $z\ne0$ is nonzero.  The elementary
integer-power rules give all assertions about $A_h$, including
the second identity in (HEP.9), since $-h_i+l_i+i=0$.
The product $A_1$ is exactly $N_{k/\mathbb Q_\ell}$.

Here is the cyclotomic normalization also for nonzero Weil degree.
Over $\mathbb Q_\ell$, arithmetic reciprocity on a unit $u$
acts on $\ell$-power roots of unity by $u^{-1}$, as verified by
the principal-idele argument in CIP5.  Arithmetic reciprocity
on the uniformizer $\ell$ acts trivially on these roots: global
reciprocity kills the positive principal idele $\ell$, its positive
real component acts trivially, and every other finite component
is a unit at a place where the cyclotomic extension is unramified.
Consequently geometric reciprocity on any
$z\in\mathbb Q_\ell^\times$ acts by
$z\ell^{-\operatorname{ord}_\ell z}$.

Local norm compatibility therefore makes geometric reciprocity on
$x\in k^\times$ act by
\[
       N_{k/\mathbb Q_\ell}(x)\,
             \ell^{-f\operatorname{ord}_k(x)}
            =A_1(x)q^{-\operatorname{ord}_k(x)},
\]
where $f=[\kappa_k:\mathbb F_\ell]$.
The original $\mathrm{IsLocalArtin}$ map is the geometric map:
its \texttt{glob} field agrees with that map under the original
global reciprocity map, and the retained local idele injectivity
makes the two values equal.  This comparison works for every
original number field and place; it does not require a split CM
place.  Equivalently it is the normalization used in the source's
\texttt{eps\allowbreak Art}.  Since
$\operatorname{ord}_k(b(s))=-\deg(s)$, the displayed formula is
the first identity of (HEP.9).  This retains every admissible
Artin map and extends the inertia identity of CIP5 by the required
factor $q^{\deg(s)}$.
\end{proof}
% END HEP4

% BEGIN HEP5
\paragraph{HEP5. Unit values and ordered partial products determine the local characters.}
\label{claim:HEP5}
Let $\phi_i:W_v\to\overline E^\times$ be characters and
$\omega_i:k^\times\to\mathbb C^\times$ smooth characters.
Suppose, whenever $b(s)$ is a unit,
\[
 \phi_i(s)\varepsilon_\ell(s)^i A_{l_i}(b(s))
                             =\iota^{-1}(\omega_i(b(s))),
 \tag{HEP.10}
\]
and, whenever $b(s)=\varpi$ for the source uniformizer,
\[
 \iota\left(\prod_{i\leq j}\phi_i(s)\varepsilon_\ell(s)^i\right)
  =\iota\left(\prod_{i\leq j}A_{l_i}(\varpi)\right)^{-1}
                                  \mathrm{psEig}(\omega,j+1)
                      \quad(0\leq j<n).
 \tag{HEP.11}
\]
Then for every $i$ and $s\in W_v$,
\[
       \phi_i(s)=\iota^{-1}(\omega_i(b(s)))A_{-h_i}(b(s)).
 \tag{HEP.12}
\]
\begin{proof}
If $b(s)$ is a unit, the degree is zero.  Substitute
$\varepsilon_\ell(s)=A_1(b(s))$ from HEP4 in (HEP.10),
and divide by its nonzero algebraic factors.  This gives
(HEP.12) on all such $s$.

If $b(s)=\varpi$, its degree is $-1$.  By definition,
\[
       \mathrm{psEig}(\omega,j+1)
                  =\prod_{i\leq j}\omega_i(\varpi)q^{-i}.
 \tag{HEP.13}
\]
Thus (HEP.11), pulled back through $\iota$, equates the ordered
partial products of $\phi_i(s)\varepsilon_\ell(s)^i$ and
$A_{l_i}(\varpi)^{-1}\iota^{-1}(\omega_i(\varpi))q^{-i}$.
All factors are nonzero.  The equality for $j=0$ identifies the
first factors; dividing the equality for $j$ by that for $j-1$
identifies each subsequent factor.  HEP4 says
$\varepsilon_\ell(s)=q^{-1}A_1(\varpi)$, so cancellation gives
(HEP.12) also on every preimage of the uniformizer.

The right side of (HEP.12) is a multiplicative character of $W_v$.
Choose $s_0$ with $b(s_0)=\varpi$, using the actual Artin
surjectivity.  For any $s$, let $m=\operatorname{ord}_k(b(s))$.
Then $b(ss_0^{-m})$ is a unit and
$s=(ss_0^{-m})s_0^m$.  Agreement on the first factor and on $s_0$
proves agreement on $s$.  No commutativity of $W_v$ is required.
The ordered initial products have recovered the individually
labeled characters, not merely their multiset.
\end{proof}
% END HEP5

% BEGIN HEP6
\paragraph{HEP6. The universal unit and uniformizer relations give the exact OrdChar.}
\label{claim:HEP6}
Fix $y:T\to_O\overline E$ and put $x=y|_A$.
Suppose smooth characters $\omega_i$ at $v$ satisfy
\begin{align}
 \omega_i(u)\,\iota\bigl(H.\mathrm{finPart}(\theta,x,\lambda,v,i,u)\bigr)
       &=1 &&(u\in\mathcal O_k^\times), \tag{HEP.14}\\
 \iota(y(t_{(v,j)}))
       &=\operatorname{lowVal}(\iota,\lambda,v,j+1)^{-1}
                                      \mathrm{psEig}(\omega,j+1).
                                                        \tag{HEP.15}
\end{align}
For every $i$, the literal
\[
 \mathrm{OrdChar}\bigl(\iota,b,\omega_i,h_i,
                                      y_*\psi_{v,i}\bigr)
 \tag{HEP.16}
\]
holds on the original decomposition group.
\begin{proof}
Put $t_i(u)=\operatorname{image}_T\theta(e_{v,i}(u))$ and
$\phi_i=y_*\psi_{v,i}$.  HEP1 proves continuity of $\phi_i$
and identifies the evaluated cyclotomic character.
The \texttt{\(\psi\)\_\allowbreak unit} field in (HEP.2) gives, for an
original Weil element mapping to $u$,
\[
                 \phi_i(s)\varepsilon_\ell(s)^i y(t_i(u))=1.
\]
The definition of the finite part and (HEP.14) give
\[
                  y(t_i(u))=A_{l_i}(u)\,
                                   \iota^{-1}(\omega_i(u))^{-1}.
\]
Combining these identities gives (HEP.10).

The \texttt{\(\psi\)\_\allowbreak unif} field says that the partial product
$\prod_{i\leq j}\phi_i(s)\varepsilon_\ell(s)^i$ equals
$y(t_{(v,j)})$ whenever $b(s)=\varpi$.
The source definition
\[
 \operatorname{lowVal}(\iota,\lambda,v,j+1)
            =\iota\left(\prod_{i\leq j}A_{l_i}(\varpi)\right)
\]
turns (HEP.15) into (HEP.11).  HEP5 now proves the entire
Weil-value identity of $\mathrm{OrdChar}$.
Together with smoothness of $\omega_i$ and the continuity already
proved, these are its three exact fields.  Its \texttt{On\allowbreak Art}
formulation allows any decomposition-group element whose underlying
automorphism equals the given Weil element; such elements are equal
in the subgroup, so the same identity verifies that full quantifier.
\end{proof}
% END HEP6

% BEGIN HEP7
\paragraph{HEP7. The full original core eigenvalue theorem.}
\label{claim:HEP7}
Under precisely the original hypotheses of
\texttt{core\_\allowbreak eigen\_\allowbreak of\_\allowbreak weight\_\allowbreak char\_\allowbreak hecke\_\allowbreak level}, for every
$y,\lambda,\xi,Y,F$ satisfying (HEP.6) with $0\ne F\in Y.Y$,
the determinant equality (HEP.7) holds.  Moreover, at every $v$,
every smooth family $\omega_i$ satisfying (HEP.14)--(HEP.15)
gives all the original characters (HEP.16).
\begin{proof}
HEP3 gives the determinant conclusion using only the spherical
eigenvector assumptions in the source.  For an arbitrary $v$ and
an arbitrary family in the theorem's second conclusion, HEP6
gives each $\mathrm{OrdChar}$.  The original
$H.\mathrm{Ax}$ supplies $H.S_l(v)\mid\ell$, the validity and
surjectivity of $H.\mathrm{art}(v)$, and the dense inverse-Frobenius
set.  The two complete local coefficient-ring assumptions supply
HEP1 at $T$; $H.\mathrm{IsWt}$ and all remaining hypotheses are
kept exactly as in the endpoint.  This proves its entire
conjunction, without adding arithmetic, good-point, domain,
nonvanishing-module, or flatness hypotheses.
\end{proof}
% END HEP7

\subsection{The original representation at every arithmetic point}

% BEGIN HEP8
\paragraph{HEP8. Each admitted arithmetic point is realized in one original part.}
\label{claim:HEP8}
Let $P(\lambda,r)$ be any original point predicate and assume
$H.\mathrm{AutomorphicHeckeSet}(\iota,P)$.
For every admitted $\lambda$ and every $y:T\to_O\overline E$
whose restriction $x=y|_A$ is arithmetic of weight $\lambda$,
there are the original specialization data, one original part $Y$,
and a nonzero vector
\[
 F\in Y.Y\cap H.\mathrm{OrdAt}
                   (\theta,\iota,x,\lambda,\xi,K,c_0)
 \tag{HEP.17}
\]
with $P(\lambda,Y.r)$ such that every original generator acts on
$F$ by $\iota(y(t_\alpha))$.  In particular $r=Y.r$ is continuous,
semisimple, satisfies $P(\lambda,r)$, and has determinant (HEP.7).
\begin{proof}
ASP13 constructs the full original specialization datum at the
specified arithmetic $x$ and weight $\lambda$, using the unchanged
coefficient isomorphism and original operators.  AFS19(3), applied
to this datum and the given point $y$ extending $x$, places a
nonzero simultaneous eigenvector in one original part with the
property $P$.  That assertion permits the literal depth $c_0$;
equivalently one may keep its existential original depth $e\geq c_0$.
It gives (HEP.17) and the eigenvalues for both spherical and local
positive generators.  The part fields give continuity and
semisimplicity of $Y.r$.  Its spherical eigenvalues satisfy (HEP.6),
so HEP7 gives the complete determinant equality.  Goodness was
not used in choosing this representation.
\end{proof}
% END HEP8

% BEGIN HEP9
\paragraph{HEP9. The same part gives the exact de Rham polynomial at a good point.}
\label{claim:HEP9}
Suppose additionally that the point in HEP8 is good.  At every $v$
there is a smooth family $\omega_{v,i}$ satisfying
(HEP.14)--(HEP.15), pairwise distinct on local units, such that
\[
 Y.\sigma_l(v)\simeq\mathrm{PS}(\omega_v),
 \qquad
 \mathrm{DRat}\left(\iota,H.\mathrm{ovr}(v),r|_{D_v},
     \tau\mapsto\{h_i(\tau)\}_i,
     s\mapsto\prod_i(X-\omega_{v,i}(b(s)))\right).
 \tag{HEP.18}
\]
The polynomial identity has the original Artin map $b$.
\begin{proof}
The nonzero $F$ in (HEP.17) makes the part's intersection with
the corresponding ordinary space nonzero.  AFS17 applies to
that same part and supplies smooth $\omega_v$, its principal-series
equivalence, the inverse finite-part equations (HEP.14), and the
scalar action of each local positive operator on the entire
intersection.  Compare that action on $F$ with its given
generator eigenvalue.  Since $F\ne0$, the two scalars are equal,
which is exactly (HEP.15), with its unchanged $j+1$ index.

If two unit characters $\omega_{v,i},\omega_{v,j}$ were equal,
their nonzero values and (HEP.14) would make the corresponding
finite parts equal at every unit.  This contradicts the distinctness
field of $H.\mathrm{Good}$.  Thus the unit characters are pairwise
distinct.

Choose an original correspondence $\mathrm{rc}$ using
$\mathrm{RecTExists}$.  We check the local polynomial used in
the \texttt{At\allowbreak LT} field.  By the original definition,
$\mathrm{PS}(\omega_v)$ is normalized induction of
$\omega_{v,i}|\cdot|^{(n-1)/2}$.  Its irreducible representation
has these singleton cuspidal support terms.  Distinct unit
restrictions exclude every linking by a unit norm twist.
The singleton Langlands calculation in GUA2, with the inverse
twist $(1-n)/2$ of $\mathrm{algNorm}$, therefore gives the actual
parameter
\[
 \mathrm{rc}(\mathrm{algNorm}(Y.\sigma_l(v)))
                   =\bigoplus_i(\omega_{v,i}\circ b_{\mathrm{rc}},0).
 \tag{HEP.19}
\]
The Artin map $b_{\mathrm{rc}}$ belonging to this correspondence
and the source $b$ both satisfy the full original
$\mathrm{IsLocalArtin}$ predicate.  Their global reciprocity
identities and the retained local idele injectivity give
$b_{\mathrm{rc}}=b$, just as in HEP4.
This verifies the needed forward principal-series calculation on
the exact source map; no arbitrary subquotient converse is invoked.

The characteristic polynomial of (HEP.19) is the product in
(HEP.18).  Every character kills commutators, so at every Weil
commutator that polynomial is $(X-1)^n$.  The rank is $n$,
hence this is exactly the additional $\mathrm{AbW}$ test in
the original $\mathrm{AtLT}$ field.  Apply $Y.\mathrm{atL}(v)$
to the chosen smooth family, its distinct unit restrictions,
the principal-series equivalence, and this allowed correspondence.
It gives $\mathrm{DRat}$ for its parameter polynomial; (HEP.19)
is precisely the replacement giving (HEP.18).  All three
$\mathrm{DRat}$ clauses, and every original inducing bridge,
remain in that assertion.
\end{proof}
% END HEP9

% BEGIN HEP10
\paragraph{HEP10. The complete original arithmetic-point field Out.pt.}
\label{claim:HEP10}
For every admitted $\lambda$ and every arithmetic point
$y:T\to_O\overline E$, the representation $r$ of HEP8 satisfies
\[
 \mathrm{Continuous}(r),\quad \mathrm{IsSSRep}(r),\quad
 P(\lambda,r),\quad D\otimes_{T,y}\overline E=\det r.
 \tag{HEP.20}
\]
If that same point is good, then for every $v$,
\[
 \mathrm{HasFlag}\bigl(r\circ H.\mathrm{dec}(v),
                                  (y_*\psi_{v,i})_i\bigr)
       \quad\text{and}\quad
 \mathrm{PotCrys}\bigl(H.\mathrm{pl}(v),r\circ H.\mathrm{dec}(v)\bigr).
 \tag{HEP.21}
\]
This is the entire literal $\mathrm{Out.pt}$ field.
\begin{proof}
HEP8 proves (HEP.20) for every arithmetic point and chooses $r$
before a good-point hypothesis is imposed.  If the point is good,
HEP9 supplies its smooth local families, distinctness, and exact
$\mathrm{DRat}$ polynomial.  HEP7 applies (HEP.14)--(HEP.15)
to supply $\mathrm{OrdChar}$ for the original evaluated characters
$y_*\psi_{v,i}$.

For each embedding $\tau$ and $i<j$, dominance of $\lambda(\tau)$
gives
\[
 \lambda(\tau)_{n-1-i}\leq\lambda(\tau)_{n-1-j}.
\]
Consequently
$h_j(\tau)-h_i(\tau)
 =\lambda(\tau)_{n-1-j}-\lambda(\tau)_{n-1-i}+j-i>0$.
The weights are strictly increasing in the precise source order.
The local representation is continuous because $r$ is continuous
and $H.\mathrm{dec}(v)$ is the original continuous local map.
The original $H.\mathrm{Ax}$ gives the odd-prime condition,
the place-over-$\ell$ condition, and the exact Artin validity.

All hypotheses of OFL16 are therefore satisfied on the unchanged
local representation with these individually labeled characters.
It gives its actual ordered flag and potential crystallinity for
every original bridge, which is (HEP.21).  In particular no
semisimplification of the local representation replaces it.
The prime $3$ is included.  The implication to (HEP.21) retains
the original good-point guard, whereas (HEP.20) holds on every
admitted arithmetic point as required.
\end{proof}
% END HEP10

% BEGIN HEP11
\paragraph{HEP11. The whole literal OrdCoreH.Out record.}
\label{claim:HEP11}
Retain every original hypothesis of
\texttt{core\_\allowbreak hecke\_\allowbreak of\_\allowbreak automorphic\_\allowbreak hecke\_\allowbreak dense\_\allowbreak weight},
including the supplied original automorphic-Hecke-set predicate
and an admitted $\lambda_0$ with $H.\mathrm{DenseWt}(\theta,\lambda_0)$.
Then the original
\[
           H.\mathrm{Out}(\theta,a,K,P,\mathfrak m_T,D,\psi)
 \tag{HEP.22}
\]
holds, with all of its \texttt{univ}, \texttt{sep}, and \texttt{pt}
fields and their original quantifiers.
\begin{proof}
The \texttt{univ} field is exactly the given (HEP.2), with no
replacement determinant or local character.
For \texttt{sep}, fix any admitted weight $\lambda$ satisfying
$\mathrm{DenseWt}$, not merely the initially supplied $\lambda_0$.
AFS20 applied to the same $T$ says that
\[
 \left[\forall y:T\to_O\overline E,\quad
 H.\mathrm{Good}(\theta,y|_A,\lambda)\Rightarrow y(t)=0\right]
                      \Longrightarrow t=0
 \tag{HEP.23}
\]
for every $t\in T$.  This is exactly the full \texttt{sep} field.
Finally HEP10 proves \texttt{pt} for every admitted weight and
every arithmetic point.  Its representation is chosen before its
conditional good-point conclusion, as the target record requires.
These three fields prove (HEP.22).  The initially supplied
$\lambda_0$ is retained as in the source theorem; no argument
restricts its later universally quantified weights to $\lambda_0$.
\end{proof}
% END HEP11

% BEGIN HEP12
\paragraph{HEP12. Installation for the original unitary Hecke data.}
\label{claim:HEP12}
Let the original $\mathrm{HypothesesH'}$, Hida setting, residual
datum $\bar r$, and weight object $\Lambda$ be given, together with
\[
                 R_H:\mathrm{UnitaryHeckeData}(Q,\iota,\Lambda).
\]
Keep any original level $K$, set
$H=hset(Q,R_H.\tau_0)$, $A=\Lambda.\Lambda$,
$\theta=R_H.\theta$, $a=R_H.a$, and let $T$ be exactly (HEP.1).
Under the original commutativity and $\mathrm{IsCO}(O,T)$ hypotheses,
every supplied (HEP.2) has the full output
\[
 H.\mathrm{Out}\bigl(\theta,a,K,
       \mathrm{UnitaryPointProperties}(\bar r,\iota),
       \mathfrak m_T,D,\psi\bigr).
 \tag{HEP.24}
\]
Thus its representation at each admitted arithmetic point carries
the original polarization, every Teichm\"uller-family inertia type,
conditional cuspidal conclusion, and full automorphic witness;
at a good point it also has exactly the prescribed flags and
potential crystallinity.  Equation (HEP.23) holds at every admitted
dense weight, including the zero weight.
\begin{proof}
FNL16 supplies every original $H.\mathrm{Ax}$ field at this
$hset$.  The original data record supplies
$R_H.\mathrm{isWt}$ and the literal $\mathrm{TypeForAlt}$ field.
UPA17 therefore supplies precisely
$H.\mathrm{AutomorphicHeckeSet}
(\iota,\mathrm{UnitaryPointProperties}(\bar r,\iota))$.
AFS5, applied to its actual $\mathrm{IsTheta}$ field, proves
$\mathrm{DenseWt}$ for every weight.  The source admitted set
for this $hset$ is the entire set of weights, so the constant zero
dominant weight is an allowed initial $\lambda_0$.

Apply HEP11 to these original data and the given $D,\psi$.
This proves (HEP.24).  The four assertions about the point
representation are exactly the four fields of the supplied
$\mathrm{UnitaryPointProperties}$, already installed by UPA15--UPA17.
The determinant, local characters, level, and Hecke algebra remain
the ones in (HEP.1)--(HEP.2); there is no new coefficient algebra
or kernel ideal in this construction.  The statement makes no
claim about a level-to-level comparison, the complete
$\mathrm{RingProps}$ record, or an all-level automorphy endpoint.
Those are separate consumers of this full output.
\end{proof}
% END HEP12


\section{The full ring properties of the original ordinary Hecke ideal}
\label{sec:ring-properties-v44}

\status{This chapter proves the literal \texttt{TJRing} consumer, including
the zero-module branch and all eight fields of \texttt{Ring\allowbreak Props}.  The
ordinary Hecke algebra, its scalar map, the universal determinant and the
ideal are the original ones.  The arithmetic inputs are the already
constructed ordinary algebra OHA11, density AFS5, and the actual
\texttt{Out} supplied by HEP12.  We prove the auxiliary-place inertia
assertion needed by this consumer, rather than adding it as a new input.
No finite-rank bound, semisimplicity assumption at that auxiliary place,
or replacement of the coefficient ring is introduced.}

Fix the original prime \(\ell\), coefficient object \(c\), and
\(O=c.O\subset\overline{\mathbb Q}_\ell\).  Thus \(O\) is a DVR
finite over \(\mathbb Z_\ell\).  Write \(k=O/\mathfrak m_O\).
For a frame \(F\), put \(G=F.G_0\), with its given involution
\(g\mapsto g^c\) and cyclotomic character \(\epsilon\).  Let
\(r\) be the original residual datum of rank \(n\),
\(B\) its supplied universal ordinary determinant ring, and
\(\Lambda\) its supplied weight object.  We write
\[
 D=B.u.D,\qquad \psi_{v,i}=B.u.\psi(v,i),\qquad
 \lambda=\Lambda.\mathrm{toPair}(B.u):\Lambda\longrightarrow B.
 \tag{JRP.1}
\]
For an ideal \(J\subset B\), put \(A=B/J\),
\(D_J=D\otimes_B A\), and \(\psi^J_{v,i}=\psi_{v,i}\otimes_B A\).
We use \([g]\) for the basis element of a group algebra.  The original
anti-involution and ordered ordinary element are
\[
 [g]^*=\epsilon(g)^{1-n}[(g^c)^{-1}],\qquad
 o_v(\psi^J;g_0,\ldots,g_{n-1})=
       \prod_{i=0}^{n-1}\bigl([\phi_v(g_i)]-\psi^J_{v,i}(g_i)[1]\bigr),
 \tag{JRP.2}
\]
where the factors are written in increasing \(i\), without commuting
them.  A degree-\(n\) determinant \(D_A\) is the original natural
multiplicative homogeneous polynomial law on the group algebra.
Its kernel predicate has the full meaning
\[
 \begin{gathered}
 D_A.\mathrm{Kills}(y)\quad\Longleftrightarrow\quad
 D_{A,C}(1+h y_C)=1\\
 \hbox{for every commutative \(A\)-algebra \(C\) and every \(h\in C[G]\)}.
 \end{gathered}
 \tag{JRP.3}
\]
Thus equality of traces at group elements alone is not what we descend.

The residual characters are \(\bar\chi_i:G\to k^\times\), and
\[
 I^p_u=I_u\cap\bigcap_i\ker\bar\chi_i.
\]
The original \(\mathrm{V0ok}\) is precisely
\[
 \bar\chi_i|_{I_0}=1\ (\hbox{all }i),\qquad
 \bar\chi_i|_{D_0}=\bar\chi_j|_{D_0}\ (\hbox{all }i,j),\qquad
 \exists g\in D_0:\ \bar\epsilon(g)\ne1.
 \tag{JRP.4}
\]
All algebra homomorphisms below preserve the given \(O\)-structure.

\subsection{Integral algebraic points and the actual auxiliary-place inertia}

% BEGIN JRP1
\paragraph{JRP1. The original coefficient ring and its Teichmuller lifts.}
\label{claim:JRP1}
The given \(O\) is the ring of integers of its fraction field
\(E_0\subset\overline{\mathbb Q}_\ell\), which is finite over
\(\mathbb Q_\ell\).  It is complete, its residue field has order
\(q=\ell^f\), and there is a multiplicative section
\(\tau:k\to O\) of reduction, with \(\tau(1)=1\) and
\(\tau(a)^{q-1}=1\) for \(a\ne0\).  In particular
\[
 t_i=\tau\circ\bar\chi_i:G\longrightarrow O^\times
 \tag{JRP.5}
\]
satisfies the original \(r.\mathrm{IsTeich}(t)\), and
\(\bar\chi_i(g)=1\) implies \(t_i(g)=1\).
\begin{proof}
The given embedding makes \(O\) torsion free over
\(\mathbb Z_\ell\); its finite module is therefore free.  Localizing
gives a finite-dimensional domain over \(\mathbb Q_\ell\), hence
the field \(E_0\).  Every element of \(O\) is integral over
\(\mathbb Z_\ell\).  Conversely an element of \(E_0\) integral
over \(\mathbb Z_\ell\) is integral over \(O\), so belongs to the
integrally closed DVR \(O\).  This identifies \(O\) with the
usual integer ring, equivalently the elements of norm at most one.
Finite freeness over the complete ring \(\mathbb Z_\ell\) gives
\(\ell\)-adic completeness.  Since \(\ell\) is a nonzero
nonunit in the DVR, \(\ell O=\mathfrak m_O^e\) for some
\(e>0\); the two adic filtrations are cofinal.  Reduction is a
finite extension of \(\mathbb F_\ell\).

For \(a\in k^\times\), lift it successively to a root of
\(X^{q-1}-1\) modulo \(\mathfrak m_O^N\).  The derivative
\((q-1)X^{q-2}\) is a unit at every such lift.  Given a lift
modulo \(\mathfrak m_O^N\), the equation for a correction in
\(\mathfrak m_O^N/\mathfrak m_O^{N+1}\) is therefore a linear
equation with a unique solution.  Completeness gives a unique root
\(\tau(a)\) reducing to \(a\).  Products of these roots have the
same property for the product of their residues, so uniqueness
proves multiplicativity.  Put \(\tau(0)=0\).  The root reducing
to one is one.  Composing with \(\bar\chi_i\) gives (JRP.5),
with exponent exactly \(\#k-1\), as required in
\(\mathrm{IsTeich}\).  These constructions use the original
coefficient ring throughout.
\end{proof}
% END JRP1

% BEGIN JRP2
\paragraph{JRP2. Algebraic points preserve the prescribed residue map.}
\label{claim:JRP2}
Let \(R\) be an \(O\)-algebra of the original class
\(\mathrm{IsCO}\), and let
\(x:R\to_O\overline{\mathbb Q}_\ell\) be any algebra homomorphism.
Then
\[
 a\in\mathfrak m_R\ \Longrightarrow\ |x(a)|_\ell<1,
 \qquad |x(b)|_\ell\le1\quad(b\in R).
 \tag{JRP.6}
\]
If \(O'\) is the integer ring of a finite coefficient extension
containing \(O\) and \(x(b)\), then
\[
 \overline{x(b)}=j_{k,k'}\bigl(\mathrm{res}_{O,R}(b)\bigr).
 \tag{JRP.7}
\]
Here \(j_{k,k'}:k\hookrightarrow k'\) is the natural residue-field
inclusion; no continuity of \(x\) is assumed.
\begin{proof}
Suppose \(a\in\mathfrak m_R\) and \(|x(a)|_\ell\ge1\).
Then \(\alpha=x(a)\ne0\), and \(\alpha^{-1}\) is integral
over \(O\): it belongs to the integer ring of the finite extension
\(E_0(\alpha)/E_0\), which is integral over \(O\).
Choose a monic equation
\[
 (\alpha^{-1})^d+c_{d-1}(\alpha^{-1})^{d-1}
                    +\cdots+c_0=0,\qquad c_i\in O.
\]
Multiplication by \(\alpha^d\) says that \(x\) kills
\(1+c_{d-1}a+\cdots+c_0a^d\).  This element lies in
\(1+\mathfrak m_R\), so is a unit in the local ring \(R\).
A unital ring map cannot send a unit to zero.  This contradiction
proves the first assertion.

By the \(\mathrm{res\_\allowbreak surj}\) field of \(\mathrm{IsCO}\),
each \(b\in R\) has \(b-b_0\in\mathfrak m_R\) for some
\(b_0\in O\).  Thus \(|x(b)-b_0|_\ell<1\).
JRP1 and the ultrametric inequality give \(|x(b)|_\ell\le1\),
and reduction in any finite coefficient field containing these
elements gives (JRP.7).  The residue of \(b_0\) is exactly
\(\mathrm{res}_{O,R}(b)\).  This argument proves the needed
residue compatibility directly from the original local ring
structure; no analytic boundedness theorem for the entire point is
needed.
\end{proof}
% END JRP2

% BEGIN JRP3
\paragraph{JRP3. The finite prime-to-residue-characteristic inertia calculation.}
\label{claim:JRP3}
Let \(K\) be a finite extension of \(\mathbb Q_p\), with
absolute group \(G_K\) and inertia \(I_K\).  If
\(c:I_K\to\mu_m\subset\overline K^\times\) has open kernel
and \(p\nmid m\), then
\[
 c(g\tau g^{-1})=g(c(\tau))
       \qquad(g\in G_K,\ \tau\in I_K).
 \tag{JRP.8}
\]
In particular a nontrivial character of \(\ell\)-power order,
\(\ell\ne p\), fixed by conjugation by \(g\), forces
\(\epsilon_\ell(g)\equiv1\pmod\ell\).
\begin{proof}
Inertia fixes \(\mu_m\): roots of unity of order prime to \(p\)
are the unique lifts of their residues.  Distinct such roots have
distinct residues, because the derivative of \(X^m-1\) is a unit
at each root.  We first find \(z\ne0\) with
\(\tau(z)=c(\tau)z\).  On the fixed field of \(\ker c\),
the finite group \(I_K/\ker c\) acts faithfully.  Its distinct
field automorphisms are linearly independent.  One proof takes a
relation with the smallest number of nonzero coefficients,
evaluates it on a product, and subtracts a suitable multiple of the
original relation; two distinct automorphisms can be separated by
one field element, contradicting minimality.  Hence the operator
\(\sum_a c(a)^{-1}a\) is nonzero.  Apply it to an element on
which it is nonzero; reindexing the sum proves the desired eigenlaw.

For fixed \(g\) put \(y=g^{-1}(z)/z\).  Every local Galois
automorphism preserves the absolute value, so \(|y|=1\).
As inertia acts trivially on residue fields,
\(|\tau(y)-y|<1\).  The eigenlaw gives
\[
 \tau(y)-y=
 \frac{g^{-1}\bigl(c(g\tau g^{-1})\bigr)-c(\tau)}{c(\tau)}\,y.
\]
The two terms in the numerator are \(m\)-th roots of unity.
Their difference has norm less than one, so their residues, and
hence the roots themselves, agree.  Applying \(g\) proves
(JRP.8).  For the last assertion choose a nontrivial value of
\(c\) and take a power having order \(\ell\).
The cyclotomic definition is
\(g(\zeta)=\zeta^{\epsilon_\ell(g)}\); fixing an element of
order \(\ell\) gives the asserted congruence.  The same proof
applies to any open subgroup of \(G_{\mathbb Q_p}\), using its
finite fixed field.
\end{proof}
% END JRP3

% BEGIN JRP4
\paragraph{JRP4. Scalar residual characteristic polynomials force auxiliary unramifiedness.}
\label{claim:JRP4}
Let \(p\ne\ell\), let \(K/\mathbb Q_p\) be finite, and let
\(O'\) be a complete DVR with finite residue field \(k'\) of
characteristic \(\ell\).  Suppose
\(\rho_0:G_K\to\mathrm{GL}_n(O')\) is continuous modulo every
power of \(\mathfrak m_{O'}\), and
\[
 \det(X-\bar\rho_0(g))=(X-\kappa(g))^n
                  \quad(g\in G_K),
 \tag{JRP.9}
\]
where \(\kappa:G_K\to(k')^\times\) is trivial on inertia.
If some \(g_*\in G_K\) has
\(\bar\epsilon_\ell(g_*)\ne1\), then
\(\rho_0|_{I_K}=1\).  No semisimplicity or bound on \(n\)
is required.
\begin{proof}
Let \(M=\rho_0(g)\), and lift \(\kappa(g)\) to \(b\in O'\).
Cayley--Hamilton applied to (JRP.9) gives
\((\bar M-\bar b I)^n=0\).  Since \(\ell^n\ge n\),
the characteristic-\(\ell\) binomial identity gives
\(M^{\ell^n}\equiv b^{\ell^n}I\pmod{\mathfrak m_{O'}}\).
For \(j\ge1\), if \(M\equiv bI\pmod{\mathfrak m^j}\),
then
\(M^\ell\equiv b^\ell I\pmod{\mathfrak m^{j+1}}\).
Indeed the linear binomial term has a factor \(\ell\in\mathfrak m\),
and every higher term has at least two factors in \(\mathfrak m^j\).
All matrices in this calculation commute with the scalar matrix.
Iteration therefore gives, for each \(N\ge1\),
\[
 \rho_0(g^{\ell^{n+N}})
       \equiv b^{\ell^{n+N}}I\pmod{\mathfrak m_{O'}^N}.
 \tag{JRP.10}
\]

Reduce \(\rho_0\) modulo \(\mathfrak m_{O'}^N\); the target
is finite.  Its inertia image \(H\) has exponent dividing
\(\ell^{n+N}\), by (JRP.10) with \(b=1\).
Cauchy's theorem makes \(H\) an \(\ell\)-group.  If \(H\)
were nontrivial it would have a nontrivial abelian character into
\(\ell\)-power roots of unity.  For completeness, a finite
\(\ell\)-group has nontrivial center by the class equation;
induction on its order gives a nontrivial homomorphism to a cyclic
group of order \(\ell\), unless the group is already abelian,
where the assertion is immediate.  Composing such a character with
the inertia quotient gives a nontrivial
\(c:I_K\to\mu_{\ell^h}\) with open kernel.

The image of \(g_*^{\ell^{n+N}}\) is scalar by (JRP.10), so
it centralizes \(H\) and fixes \(c\).  JRP3 forces
\(\bar\epsilon_\ell(g_*^{\ell^{n+N}})=1\).
But this equals
\(\bar\epsilon_\ell(g_*)^{\ell^{n+N}}\), which cannot be one:
Frobenius is injective on the finite field \(k'\).
Thus \(H=1\).  This holds for every \(N\); separation of the
adic filtration now gives \(\rho_0(\tau)=1\) for every
\(\tau\in I_K\).
\end{proof}
% END JRP4

% BEGIN JRP5
\paragraph{JRP5. The literal auxiliary-place assertion for every original point.}
\label{claim:JRP5}
In the original Hida setting let \(n\ge1\), assume
\(r.\mathrm{V0ok}\), and let
\(x:B\to_O\overline{\mathbb Q}_\ell\) and
\(\rho:F.G_0\to\mathrm{GL}_n(\overline{\mathbb Q}_\ell)\)
satisfy
\[
 \rho\text{ is continuous},\qquad
                 D\otimes_{B,x}\overline{\mathbb Q}_\ell=\det\rho.
 \tag{JRP.11}
\]
Then \(\rho(g)=1\) for every \(g\in F.I_0\).
This is the complete original
\texttt{point\_\allowbreak rep\_\allowbreak unramified\_\allowbreak at\_\allowbreak auxiliary\_\allowbreak place} endpoint.
\begin{proof}
We first explain the coefficient issue.  On any compact local
decomposition group, ADK1 puts the coefficients of the continuous
representation in one finite extension of \(\mathbb Q_\ell\).
Enlarge that field to contain \(E_0\), without changing \(O\)
or the point \(x\).  Compactness bounds all matrix entries.
Consequently
\[
 \mathcal L=\sum_g\rho(g)O'^n
\]
lies between \(O'^n\) and \(\varpi^{-M}O'^n\) for some
\(M\).  It is a finitely generated torsion-free module over the
DVR, hence free of rank \(n\).  A basis conjugates the
representation into \(\mathrm{GL}_n(O')\), continuously modulo
each \(\mathfrak m_{O'}^N\).  Conjugation preserves characteristic
polynomials.  JRP2 applied to every coefficient of
\(D.\chi([g])\), together with \(B.u.\mathrm{red}\), proves
that the residual characteristic polynomial of this lattice is
\(\prod_i(X-\bar\chi_i(g))\), transported through \(k\to k'\).
It is the original residue map, even though this argument temporarily
uses a larger field to choose a stable lattice.

Here is the exact local dictionary at the auxiliary place.  Put
\(F^+=L^+\), let \(v_0\) be the prescribed place, and let
\(p\ne\ell\) be its residue characteristic from the Hida setting.
Choose the specified valuation above \(v_0\).  The decomposition
embedding identifies its local absolute group with an open subgroup
of \(G_{\mathbb Q_p}\), preserving inertia and cyclotomic
characters.  The Hida setting puts this decomposition group in the
image of \(G_L\to G_{L^+}\).  Thus its composite with the
quotient map gives a continuous homomorphism
\(\psi_0:G_{L^+_{v_0}}\to F.G_0\) whose image is \(F.D_0\)
and whose inertia image is \(F.I_0\).  More explicitly,
\(\psi_0(t)=V.\mathrm{pr}(e)\) whenever a decomposition element
\(e\in G_L\) restricts to \(t\); independence follows from
the original quotient, and the formula for its underlying element
is the continuous map \(\mathrm{toGS}(t)\).  This is the source
local bridge, not a newly chosen Galois quotient.

Restrict (JRP.11) along \(\psi_0\).  Since \(n\ge1\), choose
the first residual character.  Formula (JRP.4) makes the local
residual polynomial
\[
                  (X-\kappa(t))^n,
       \qquad\kappa=\bar\chi_0\circ\psi_0,
\]
with \(\kappa\) trivial on inertia.  The last clause of (JRP.4)
supplies a local element with nontrivial residual cyclotomic value;
the residue-field inclusion \(k\to k'\) preserves this
inequality.  JRP4 therefore kills the full local inertia.
Its image is precisely \(F.I_0\), proving the claim.  Neither
the global nor the local representation was assumed semisimple.
\end{proof}
% END JRP5

\subsection{Descent of the full determinant laws}

% BEGIN JRP6
\paragraph{JRP6. Jointly separating points detect polynomial laws and their kernels.}
\label{claim:JRP6}
Let \(A\) be a commutative ring and
\(f_s:A\to K_s\) a family of ring maps such that
\(\bigcap_s\ker f_s=0\).  Equality of determinants of the same
degree over \(A\) can be tested after every \(f_s\).
Moreover, if \(D_A\otimes f_s=\det\rho_s\) and
\(\rho_s(y)=0\) for every \(s\), with coefficients of
\(y\in A[G]\) mapped through \(f_s\), then
\(D_A.\mathrm{Kills}(y)\) in the full sense of (JRP.3).
The corresponding coefficientwise test holds for any natural
polynomial identity, including the polarization identity.
\begin{proof}
Coefficientwise, the maps
\(A[X_1,\ldots,X_m]\to K_s[X_1,\ldots,X_m]\) jointly
separate: every coefficient of a polynomial killed by all of them
is zero.  Given a finite list of group elements \(g_1,\ldots,g_m\),
apply the two determinant laws to the generic element
\(\sum_jX_j[g_j]\).  Naturality identifies coefficient
specialization of the resulting polynomials with the corresponding
laws over \(K_s\).  Their equality at all \(s\) therefore gives
equality over \(A[X_1,\ldots,X_m]\).  Specializing \(X_j\) to
arbitrary elements of an arbitrary \(A\)-algebra proves equality
on every group-algebra element, since each has finite support.

For the kernel assertion, use the polynomial
\[
 D_{A,A[\mathbf X]}
       \left(1+\left(\sum_jX_j[g_j]\right)y\right)-1.
 \tag{JRP.12}
\]
After \(f_s\), the matrix of \(y\) is zero, so this polynomial
is zero.  All its coefficients vanish by separation.  Specialize
the variables to the finitely many coefficients of any
\(h\in C[G]\); naturality proves exactly (JRP.3).
For polarization use instead
\(D((\sum_jX_j[g_j])^*)-D(\sum_jX_j[g_j])\).
The scalar cyclotomic coefficients in \(*\) map compatibly, and
the same argument applies.  No injectivity after tensoring with an
arbitrary algebra, no flatness, and no reducedness assumption is
used.  If the index family is empty, its separation hypothesis
forces \(A=0\), and the same coefficient proof remains valid.
\end{proof}
% END JRP6

% BEGIN JRP7
\paragraph{JRP7. The original pairing gives the entire star identity.}
\label{claim:JRP7}
Suppose \(\rho:G\to\mathrm{GL}_n(K)\) and
\(P\in\mathrm{GL}_n(K)\) satisfy the original
\(F.\mathrm{IsForm}(O,n,\rho,P)\):
\[
 \rho(g)^{\mathsf T}P\rho(g^c)=\epsilon(g)^{1-n}P.
 \tag{JRP.13}
\]
Then, for every commutative \(K\)-algebra \(C\) and every
\(y\in C[G]\),
\[
 \rho_C(y)^{\mathsf T}P_C=P_C\rho_C(y^*),\qquad
       \det\rho_C(y^*)=\det\rho_C(y).
 \tag{JRP.14}
\]
In particular \(\rho(y)=0\) implies \(\rho(y^*)=0\).
\begin{proof}
Multiply (JRP.13) on the right by \(\rho(g^c)^{-1}\).
The result is (JRP.14)'s first equality for \([g]\), with
exactly the exponent \(1-n\) in (JRP.2).  Scalar extension
preserves it; linearity gives it for every finite group-algebra
sum.  Taking determinants and canceling the unit \(\det P_C\)
proves the second identity.  When the left matrix is zero,
invertibility of \(P\) proves the kernel assertion.  No symmetry
or alternating property of \(P\) is required by this argument
or by the original \(\mathrm{RingProps}\).
\end{proof}
% END JRP7

% BEGIN JRP8
\paragraph{JRP8. A flag gives both original ordinary relations, in their order.}
\label{claim:JRP8}
Let \(\rho:H\to\mathrm{GL}_n(K)\) have a full flag with
successive characters \(\xi_0,\ldots,\xi_{n-1}\), in the
literal \(\mathrm{HasFlag}\) convention of upper triangular
matrices.  Then
\[
 \det(X-\rho(h))=\prod_{i=0}^{n-1}(X-\xi_i(h)),
 \qquad
 \prod_{i=0}^{n-1}\bigl(\rho(h_i)-\xi_i(h_i)I\bigr)=0
 \tag{JRP.15}
\]
for every \(h\in H\) and every ordered tuple \((h_i)\).
\begin{proof}
Conjugate once to the upper triangular basis supplied by
\(\mathrm{HasFlag}\).  The first equality is the determinant
of an upper triangular matrix.  Let \(F_j\) be the span of
its first \(j\) basis vectors.  Each
\(\rho(h_i)-\xi_i(h_i)I\) preserves the flag and sends
\(F_{i+1}\) into \(F_i\): the action on
\(F_{i+1}/F_i\) is zero.  The rightmost factor \(i=n-1\)
therefore sends the entire space into \(F_{n-1}\); the next
factor sends this into \(F_{n-2}\), and so on until \(F_0=0\).
This is precisely the ascending written order in (JRP.15).
Conjugating back proves both statements.  For rank zero the
vector space and its identity endomorphism are zero, so the empty
matrix product has the same assertion.  No factors with different
indices were commuted.
\end{proof}
% END JRP8

% BEGIN JRP9
\paragraph{JRP9. The quotient by the specified ideal has separating points.}
\label{claim:JRP9}
Suppose maps \(x_s:B\to_O K\) satisfy
\[
 J\subset\ker x_s\quad(\hbox{all }s),\qquad
       \bigl[x_s(b)=0\ (\hbox{all }s)\bigr]\Longrightarrow b\in J.
 \tag{JRP.16}
\]
They induce jointly separating maps \(f_s:A=B/J\to_O K\).
If \(D\otimes x_s=\det\rho_s\), then
\(D_J\otimes f_s=\det\rho_s\), and the specialized ordinary
characters are \(f_s\psi^J_{v,i}=x_s\psi_{v,i}\).
\begin{proof}
Define \(f_s(\bar b)=x_s(b)\).  The first clause of
(JRP.16) proves independence of the representative.  If every
\(f_s(\bar b)\) is zero, the second clause gives \(b\in J\),
hence \(\bar b=0\).  The quotient and specialization maps
compose to \(x_s\).  Naturality of determinant scalar extension
and the entrywise definition of pushing a unit-valued character
give the final two equalities.  These are equalities on the actual
quotient \(B/J\), rather than on a reduced or characteristic-zero
replacement of it.
\end{proof}
% END JRP9

% BEGIN JRP10
\paragraph{JRP10. Descent of the full polarization field.}
\label{claim:JRP10}
In JRP9, assume each \(\rho_s\) has a matrix satisfying
\(F.\mathrm{IsForm}(O,n,\rho_s,P_s)\).  Then
\(F.\mathrm{IsPolar}(O,D_J)\), with its universal scalar-algebra
quantifier.  Furthermore any \(y\in A[G]\) whose matrix is zero
at all these points satisfies both
\(D_J.\mathrm{Kills}(y)\) and \(D_J.\mathrm{Kills}(y^*)\).
\begin{proof}
JRP7 proves the polarization polynomial identity after each
\(f_s\).  JRP6 descends it coefficientwise to \(A\), then to
all its scalar algebras.  This is exactly \(\mathrm{IsPolar}\).
For the second assertion, JRP7 makes the specialized matrices of
both \(y\) and \(y^*\) zero.  Apply JRP6's kernel test to each.
The proof retains the universal definition of \(\mathrm{Kills}\)
and does not infer it merely from a characteristic polynomial at
\(y\).
\end{proof}
% END JRP10

% BEGIN JRP11
\paragraph{JRP11. Descent of O1 and O2 for the same ordinary characters.}
\label{claim:JRP11}
In JRP9--JRP10, suppose every local restriction
\(\rho_s\circ\phi_v\) has \(\mathrm{HasFlag}\) with the
ordered characters \(x_s\psi_{v,i}\).  Then
\[
 D_J.\chi([\phi_v(g)])=
                  \prod_i\bigl(X-\psi^J_{v,i}(g)\bigr),
 \tag{JRP.17}
\]
and for every original tuple \(g_i\in F.G_v\), both
\(o_v(\psi^J;\mathbf g)\) and
\(o_v(\psi^J;\mathbf g)^*\) are killed by \(D_J\).
\begin{proof}
After each \(f_s\), the characteristic polynomial of the
determinant is the characteristic polynomial of the actual matrix
\(\rho_s(\phi_v(g))\).  JRP8 makes it the indicated product.
Every coefficient of the difference in (JRP.17) is killed by all
\(f_s\), so is zero by JRP9.  This is O1, with the original
pushforward characters.

Specializing (JRP.2) gives the ordered product of matrices in the
second formula of (JRP.15), because coefficient maps preserve each
individual factor and its order.  This product is zero.  JRP10
therefore kills both the original ordinary element and its star.
This is O2 in full, for arbitrary tuples, not only the tuple whose
entries are all equal.
\end{proof}
% END JRP11

% BEGIN JRP12
\paragraph{JRP12. Descent of every prescribed type, O3, and the guarded auxiliary field.}
\label{claim:JRP12}
In JRP9--JRP10, suppose that for every original Teichmuller
system \(t\) satisfying \(r.\mathrm{IsTeich}(t)\), every
\(s,u\), one matrix simultaneously diagonalizes
\(\rho_s|_{I_u}\) as \(\operatorname{diag}(t_i|_{I_u})\).
Then the full \(\mathrm{types}\) and O3 fields of
\(\mathrm{RingProps}\) hold.  If also
\(r.\mathrm{V0ok}\Rightarrow\rho_s|_{I_0}=1\) for every
\(s\), its \(\mathrm{unr0}\) field holds.
\begin{proof}
For any supplied \(t\), conjugation does not change the
determinant polynomial law on \(K[I_u]\).  Thus every
\(f_s\) identifies the restriction of \(D_J\) with the
determinant of the diagonal representation whose entries are the
images of \(t_i\) under \(O\to A\).  JRP6 descends equality
of the entire laws.  This proves \(\mathrm{types}\) for every
\(t\), retaining its quantifier before \(u\).

For O3 use the specific system constructed in JRP1.  If
\(g\in I^p_u\), all \(\bar\chi_i(g)=1\), so all
\(t_i(g)=1\).  The supplied diagonalization then gives
\(\rho_s(g)=1\) for every \(s\).  Apply JRP10 to
\(y=[g]-1\).  Since \(*\) is additive and fixes one,
\(y^*=[g]^*-1\); these are exactly the two O3 kernel
statements.  Under the additional displayed guard, the same
argument for \(g\in I_0\) proves \(\mathrm{unr0}\).
The guard \(\mathrm{V0ok}\) has not been dropped.
\end{proof}
% END JRP12

% BEGIN JRP13
\paragraph{JRP13. Transfer through the original surjection, with the identical kernel.}
\label{claim:JRP13}
Let \(f:B\twoheadrightarrow T\) be the original
\(O\)-algebra map, with \(J=\ker f\), and suppose
\[
                   f(\lambda(a))=\mathrm{scal}_T(a)
                                  \quad(a\in\Lambda).
 \tag{JRP.18}
\]
Assume \(T\) is finite and torsion free over \(\Lambda\),
and a family of good maps \(z:T\to_O K\) jointly separates
\(T\).  If their composite points \(x=z\circ f\) have the
representations, forms, flags, types and guarded inertia properties
of JRP10--JRP12, then \(B.\mathrm{RingProps}(\Lambda,J)\).
\begin{proof}
For this family of composite points, (JRP.16) follows directly:
all points kill \(\ker f\), and vanishing at them forces
\(f(b)=0\) by separation in \(T\).  JRP10--JRP12 therefore
give \(\mathrm{polar}\), O1, O2, O3,
\(\mathrm{types}\) and \(\mathrm{unr0}\) on \(B/J\).

The surjection gives an actual \(O\)-algebra isomorphism
\(e:B/J\simeq T\), \(e(\bar b)=f(b)\).
Equation (JRP.18) makes it \(\Lambda\)-linear for the
original scalar map \(a\mapsto\overline{\lambda(a)}\).
Pulling back a finite generating set for \(T\) proves exactly
the \(\mathrm{fin}\) field.  If \(a\ne0\) and
\(\lambda(a)b\in J\), apply \(f\) to get
\(\mathrm{scal}_T(a)f(b)=0\).  Torsion freeness gives
\(f(b)=0\), hence \(b\in J\), which is precisely
the \(\mathrm{tf}\) field.  No new ideal was defined in this
argument, and no passage to a fraction field was used to establish
these two integral conclusions.
\end{proof}
% END JRP13

% BEGIN JRP14
\paragraph{JRP14. The entire top-ideal branch.}
\label{claim:JRP14}
\leavevmode\par
For the original supplied \(B,\Lambda\), one has
\(B.\mathrm{RingProps}(\Lambda,\top)\).
\begin{proof}
The quotient \(B/\top\) is the zero ring.  Its scalar map from
\(\Lambda\) is finite, generated as a module by the empty set
(or by zero).  The \(\mathrm{tf}\) implication has conclusion
\(b\in\top\), so always holds.  Every commutative algebra over
the zero ring is itself zero, because its unital scalar map sends
\(1=0\) to \(1=0\).  Consequently all values of every
polynomial law agree, every polynomial equality holds, and every
\(\mathrm{Kills}\) assertion is true with its full scalar-algebra
quantifier.  This proves polarization, O1, O2, O3, every type
equality and the guarded \(\mathrm{unr0}\) assertion.  In
particular no nonempty family of arithmetic points is assumed in
this branch.
\end{proof}
% END JRP14

% BEGIN JRP15
\paragraph{JRP15. The full original TJRing endpoint.}
\label{claim:JRP15}
\leavevmode\par
Assume exactly the original \(H'=\mathrm{HypothesesH'}\),
the Hida data \(\ell,L,\Omega,s,n\) with \(n\ge1\),
\(V,Q,c,r,\iota,B,\Lambda\), an original
\(R:\mathrm{UnitaryHeckeData}(Q,\iota,\Lambda)\), and a
level \(K\), with \(V.S'=s.S\).  For every ideal
\(J\subset B\),
\[
             R.\mathrm{IsJ}(B,K,J)
                    \quad\Longrightarrow\quad
             B.\mathrm{RingProps}(\Lambda,J).
 \tag{JRP.19}
\]
\begin{proof}
Expand the original \(\mathrm{IsJ}\) disjunction.  If the
ordinary residual module is zero, its clause gives \(J=\top\),
so JRP14 applies.  Otherwise retain its actual nonzero module,
the given commutativity and \(\mathrm{IsCO}\) structures, and
its surjection
\[
 f:B\twoheadrightarrow
 T=(\mathrm{hset}(Q,R.\tau_0)).TK(R.\theta,R.a,K),
 \qquad J=\ker f,
 \tag{JRP.20}
\]
together with the original \(\mathrm{IsMap}\) witness and
(JRP.18).  No surjection or ideal is reselected.

AFS5 makes this given \(\Lambda\) a domain and supplies
\(\mathrm{DenseWt}(R.\theta,\lambda_0)\) for the zero
dominant weight \(\lambda_0\).  That weight is admitted by
the original \(\mathrm{hset}\).  OHA11 applied at this same
level proves that this same \(T\) is finite over \(\Lambda\)
and that its nonzero scalars act injectively.  The hypotheses
\(\mathrm{IsWt}\), \(\mathrm{IsRes}\) and the actual
ordinary-module axioms are supplied by the fixed \(R\) and
the installation used in HEP12.

Use HEP12's \(\mathrm{Out}\) for the same \(D\otimes f\),
characters \(f\psi_{v,i}\), algebra \(T\), and zero weight.
Its \(\mathrm{sep}\) field, supplied by AFS20, says exactly
that the maps \(z:T\to_O\overline{\mathbb Q}_\ell\) whose
restrictions are good at \(\lambda_0\) jointly separate \(T\).
For every such \(z\), its \(\mathrm{pt}\) field supplies
one actual continuous representation \(\rho_z\), with
\[
 (D\otimes f)\otimes z=\det\rho_z.
\]
Its \(\mathrm{UnitaryPointProperties}\) give the original
form matrix and the simultaneous inertia diagonalizations for
\emph{every} Teichmuller system.  The same point field gives
\(\mathrm{HasFlag}\) with the ordered characters
\(z(f\psi_{v,i})\) under this exact good-point guard.
These are the data used in JRP10--JRP12.  If
\(r.\mathrm{V0ok}\), apply JRP5 to the composite point
\(x=z\circ f\) and to this very \(\rho_z\).  Functoriality
of determinant scalar extension gives (JRP.11), so the missing
auxiliary inertia condition is proved for each selected point.

JRP13 now supplies all eight fields at the original ideal in
(JRP.20).  Their exact dictionary is
\begin{center}
\begin{tabular}{ll}
\texttt{fin}, \texttt{tf} & OHA11, AFS5 and JRP13, over the original \(\Lambda\);\\
\texttt{polar} & the given form and JRP10;\\
\texttt{O1}, \texttt{O2} & the original ordered flags and JRP11;\\
\texttt{O3}, \texttt{types} & every original inertia type and JRP12;\\
\texttt{unr0} & the retained \(\mathrm{V0ok}\) guard, JRP5 and JRP12.
\end{tabular}
\end{center}
The arguments retain \(n\ge1\), every original level and
coefficient object, and the equality \(V.S'=s.S\).
This proves the whole source
\texttt{ring\_\allowbreak props\_\allowbreak of\_\allowbreak is\_\allowbreak j\_\allowbreak unitary\_\allowbreak hecke\_\allowbreak data},
not only its reduction or a property at individual points.
\end{proof}
% END JRP15

\paragraph{Scope.}
JRP15 reconstructs the full \texttt{TJRing} producer with the named
global and ordinary inputs retained.  It proves a component of the
original \texttt{Hida\allowbreak Props}; it does not claim the other arithmetic
purity, classicality, twisting or global counting fields merely from
this ring argument.  The auxiliary assertion JRP5 is independently
available for every original algebraic point satisfying (JRP.11).

\endgroup
\end{document}
